% Consolidation des acquis de la méthodologie de l'exercice sur texte philosophique — Philosophie Tle Tech — Cours signé OnBuch+
% Programme officiel MINESEC. Compiler avec : tectonic N03-consolidation-acquis-methodologie-exercice-texte-philos.tex
\newcommand{\DOCMATIERE}{Philosophie}
\newcommand{\DOCNIVEAU}{Tle Tech}
\newcommand{\DOCMODULE}{Philosophie – classe de Terminale technique}
\newcommand{\DOCLECON}{N03}
\newcommand{\DOCTITRE}{Consolidation des acquis de la méthodologie de l'exercice sur texte philosophique}
% OnBuch+ — préambule « cours signés » ENRICHI (compilé avec tectonic / XeLaTeX)
% Système visuel « Pop » d'OnBuch+ : fond crème (jamais blanc), bordures encre
% épaisses, ombres « dures » décalées, titres Archivo Black en capitales.
% Version MATHÉMATIQUES : graphes (pgfplots/tikz), boîtes pédagogiques
% (définition, propriété, méthode, attention, le savais-tu, exercice résolu,
% exercice, TP, bilan), page de garde et signature OnBuch+ ancrée en bas.
%
% Macros attendues dans main.tex AVANT \input{preamble} :
%   \DOCMATIERE  \DOCNIVEAU  \DOCMODULE  \DOCLECON (numéro)  \DOCTITRE
\documentclass[11pt]{article}

\usepackage{fontspec}
\usepackage[french]{babel}
\usepackage[a4paper,margin=2cm,top=3.4cm,bottom=2.8cm,headheight=1.4cm,footskip=1.3cm]{geometry}
\usepackage{amsmath,amssymb,amsfonts}
\usepackage{mathtools} % \xmapsto, \coloneqq... souvent utilisés par les modèles
\usepackage{cancel} % simplifications barrées (bilans de forces)
% \abs{z} (module, valeur absolue) et \norm{u} (norme), utilisés par les modèles
\providecommand{\abs}{}\let\abs\relax\DeclarePairedDelimiter{\abs}{\lvert}{\rvert}
\providecommand{\norm}{}\let\norm\relax\DeclarePairedDelimiter{\norm}{\lVert}{\rVert}
\usepackage[table]{xcolor} % [table] : fournit aussi \rowcolors (alternance de lignes)
\usepackage{pagecolor}
\usepackage{tikz}
\usetikzlibrary{arrows.meta,positioning,calc,shapes.geometric,shapes.misc,shapes.arrows,decorations.pathmorphing,decorations.pathreplacing,decorations.markings,patterns,shadows,mindmap,trees,fit,backgrounds,matrix,intersections,angles,quotes}
\usepackage{pgfplots}
\usepackage[version=4]{mhchem} % équations nucléaires \ce{^{235}_{92}U}
\usepackage[european,siunitx]{circuitikz} % schémas électriques
\pgfplotsset{compat=1.18}
\usepgfplotslibrary{fillbetween,groupplots} % aires entre courbes (calcul intégral)
\usepackage{enumitem}
\usepackage{fancyhdr}
% [most] charge toutes les bibliothèques tcolorbox (listings, minted, raster,
% documentation...) et gonfle inutilement la mémoire TeX sur un document long
% (~300 boîtes) : on ne charge que ce qui est réellement utilisé.
\usepackage[skins,breakable]{tcolorbox}
\usepackage{graphicx}
\usepackage{colortbl}
\usepackage{booktabs}
\usepackage{tabularx}
\usepackage{multirow}
\usepackage{diagbox} % case d'en-tête diagonale des tableaux à double entrée
% Colonnes extensibles centrée / gauche / droite (C, L, R), souvent utilisées par les modèles
\newcolumntype{C}{>{\centering\arraybackslash}X}
\newcolumntype{L}{>{\raggedright\arraybackslash}X}
\newcolumntype{R}{>{\raggedleft\arraybackslash}X}
\usepackage{array}
\usepackage{multicol}
\usepackage{siunitx}
% parse-numbers=false : accepte aussi \SI{2,5 \times 10^{-3}}{...}
\sisetup{locale=FR,output-decimal-marker={,},group-minimum-digits=4,parse-numbers=false}
\DeclareSIUnit\ohm{\ensuremath{\symup{\Omega}}} % U+2126 absent de la police
\DeclareSIUnit\division{div} % réglages d'oscilloscope (ms/div, V/div)
\DeclareSIUnit\tr{tr} % tours (vitesse de rotation, ex. \SI{1500}{\tr\per\minute})
\DeclareSIUnit\molar{M} % concentration molaire (ex. \SI{3}{\molar}, \SI{1}{\micro\molar})
\DeclareSIUnit\an{an} % année (le modèle confond parfois avec \year, un compteur TeX) — ex. \SI{5}{\kilo\meter\per\an}
\DeclareSIUnit\FCFA{FCFA} % franc CFA (ex. \SI{500}{\FCFA\per\kilo\gram})
\DeclareSIUnit\degreeBrix{\textdegree Brix} % teneur en sucre d'un sirop/jus (réfractométrie)
\DeclareSIUnit\month{mois} % le modèle utilise parfois \month, un compteur TeX, comme unité de durée
\DeclareSIUnit\microliter{\micro\liter} % le modèle écrit parfois \microliter en un seul token
\DeclareSIUnit\million{millions} % dénombrement (ex. \SI{15}{\million\per\mL} spermatozoïdes)
\usepackage{qrcode}
% \degree : commande fréquemment utilisée par les modèles pour un angle ou une
% température en dehors de \SI{}{} (non fournie par les paquets chargés).
\providecommand{\degree}{^{\circ}}
% \qed (fin de démonstration), utilisé par les modèles sans amsthm
\providecommand{\qed}{\hfill$\square$}
% \barre : double barre des valeurs interdites dans un tableau de variations (tkz-tab)
\providecommand{\barre}{\ensuremath{\|}}
% environnement proof (amsthm non chargé) utilisé par les modèles
\ifdefined\proof\else\newenvironment{proof}[1][Démonstration]{\par\noindent\textit{#1.}\ }{\qed\par}\fi
\usepackage{lastpage}
\usepackage{hyperref}
\hypersetup{hidelinks}

% ============================================================
% Palette « Pop » OnBuch+ — mêmes hex que la classe P (plus_ui.dart)
% ============================================================
\definecolor{poporange}{HTML}{F59321}
\definecolor{popdark}{HTML}{E07A0C}
\definecolor{popcream}{HTML}{FFF6E8}
\definecolor{popink}{HTML}{1C1714}
\definecolor{poppurple}{HTML}{7A5AE0}
\definecolor{popgreen}{HTML}{2E9E6A}
\definecolor{poppink}{HTML}{E0457B}
\definecolor{popblue}{HTML}{2F7DE1}
\definecolor{popgold}{HTML}{C98A12}
\definecolor{popmuted}{HTML}{8A7F76}
\definecolor{popline}{HTML}{EADFD2}
\definecolor{popgray}{HTML}{8A7F76}
\colorlet{popgrey}{popgray}
% Alias tolérants (le modèle nomme parfois les couleurs différemment)
\colorlet{popwhite}{white}
\colorlet{popblack}{popink}
\colorlet{popred}{poppink}
\colorlet{popyellow}{popgold}
% Teintes claires (fonds de tableaux/encadrés) — le modèle nomme parfois
% les couleurs avec un suffixe L (ex. popblueL) sans les avoir définies.
\colorlet{poporangeL}{poporange!20}
\colorlet{popdarkL}{popdark!20}
\colorlet{poppurpleL}{poppurple!20}
\colorlet{popgreenL}{popgreen!20}
\colorlet{poppinkL}{poppink!20}
\colorlet{popblueL}{popblue!20}
\colorlet{popgoldL}{popgold!20}
\colorlet{popredL}{popred!20}
\colorlet{popgrayL}{popgray!20}
% Teintes claires pour fonds de tableaux / zones
\colorlet{popblueL}{popblue!10!white}
\colorlet{poporangeL}{poporange!14!white}
\colorlet{popgreenL}{popgreen!12!white}
\colorlet{poppurpleL}{poppurple!12!white}
\colorlet{poppinkL}{poppink!12!white}
\colorlet{popgoldL}{popgold!14!white}

\pagecolor{popcream}

% ============================================================
% Typographie
% ============================================================
\defaultfontfeatures{Ligatures=TeX}
\setmainfont{PlusJakartaSans-Medium.ttf}[Path=../../../fonts/,BoldFont=PlusJakartaSans-Bold.ttf]
% Maths en sans-serif (Fira Math) avec chiffres et lettres droites de Plus
% Jakarta, pour que les formules se fondent dans le texte.
\usepackage[math-style=ISO,bold-style=ISO]{unicode-math}
\setmathfont{FiraMath-Regular.otf}[Path=../../../fonts/]
\setmathfont{PlusJakartaSans-Medium.ttf}[Path=../../../fonts/,range={up/num,up/{Latin,latin},\mathup}]
\usepackage{icomma} % virgule décimale sans espace en mode math
% Caractères mathématiques Unicode absents des polices de texte (Plus Jakarta,
% Archivo Black) : rendus par leur équivalent mathématique, en texte comme en
% formule — sinon ils s'affichent en carré vide (ex. « Arithmétique dans ℤ »).
\usepackage{newunicodechar}
\newunicodechar{ℝ}{\ensuremath{\mathbb{R}}}
\newunicodechar{ℤ}{\ensuremath{\mathbb{Z}}}
\newunicodechar{ℂ}{\ensuremath{\mathbb{C}}}
\newunicodechar{ℕ}{\ensuremath{\mathbb{N}}}
\newunicodechar{ℚ}{\ensuremath{\mathbb{Q}}}
\newunicodechar{𝔻}{\ensuremath{\mathbb{D}}}
\newunicodechar{α}{\ensuremath{\alpha}}
\newunicodechar{β}{\ensuremath{\beta}}
\newunicodechar{γ}{\ensuremath{\gamma}}
\newunicodechar{δ}{\ensuremath{\delta}}
\newunicodechar{ε}{\ensuremath{\varepsilon}}
\newunicodechar{θ}{\ensuremath{\theta}}
\newunicodechar{λ}{\ensuremath{\lambda}}
\newunicodechar{μ}{\ensuremath{\mu}}
\newunicodechar{σ}{\ensuremath{\sigma}}
\newunicodechar{τ}{\ensuremath{\tau}}
\newunicodechar{φ}{\ensuremath{\varphi}}
\newunicodechar{ψ}{\ensuremath{\psi}}
\newunicodechar{ω}{\ensuremath{\omega}}
\newunicodechar{Σ}{\ensuremath{\Sigma}}
% majuscules produites par \MakeUppercase dans les titres (α-aminés -> Α)
\newunicodechar{Α}{\ensuremath{\alpha}}
\newunicodechar{Β}{\ensuremath{\beta}}
\newunicodechar{Γ}{\ensuremath{\Gamma}}
\newunicodechar{Δ}{\ensuremath{\Delta}}
\newunicodechar{Ε}{\ensuremath{\varepsilon}}
\newunicodechar{Θ}{\ensuremath{\Theta}}
\newunicodechar{Λ}{\ensuremath{\Lambda}}
\newunicodechar{Μ}{\ensuremath{\mu}}
\newunicodechar{Τ}{\ensuremath{\tau}}
\newunicodechar{Φ}{\ensuremath{\Phi}}
\newunicodechar{Ψ}{\ensuremath{\Psi}}
\newunicodechar{Ω}{\ensuremath{\Omega}}
\newunicodechar{↦}{\ensuremath{\mapsto}}
\newunicodechar{∈}{\ensuremath{\in}}
\newunicodechar{∉}{\ensuremath{\notin}}
\newunicodechar{∧}{\ensuremath{\wedge}}
\newunicodechar{⇔}{\ensuremath{\Leftrightarrow}}
\newunicodechar{⟺}{\ensuremath{\Longleftrightarrow}}
\newunicodechar{⇒}{\ensuremath{\Rightarrow}}
\newunicodechar{⟹}{\ensuremath{\Longrightarrow}}
\newunicodechar{∘}{\ensuremath{\circ}}
\newunicodechar{∪}{\ensuremath{\cup}}
\newunicodechar{∩}{\ensuremath{\cap}}
\newunicodechar{≡}{\ensuremath{\equiv}}
\newunicodechar{⊂}{\ensuremath{\subset}}
\newunicodechar{⊥}{\ensuremath{\perp}}
\newunicodechar{∃}{\ensuremath{\exists}}
\newunicodechar{∀}{\ensuremath{\forall}}
\newunicodechar{∛}{\ensuremath{\sqrt[3]{\,}}}
\newunicodechar{∜}{\ensuremath{\sqrt[4]{\,}}}
\newunicodechar{⃗}{\ensuremath{{}^{\to}}}
\newunicodechar{ⁿ}{\ifmmode^{n}\else\textsuperscript{\ensuremath{n}}\fi}
\newunicodechar{ˣ}{\ifmmode^{x}\else\textsuperscript{\ensuremath{x}}\fi}
\newunicodechar{ᵇ}{\ifmmode^{b}\else\textsuperscript{\ensuremath{b}}\fi}
\newunicodechar{ʳ}{\ifmmode^{r}\else\textsuperscript{\ensuremath{r}}\fi}
\newunicodechar{ᵘ}{\ifmmode^{u}\else\textsuperscript{\ensuremath{u}}\fi}
\newunicodechar{ʸ}{\ifmmode^{y}\else\textsuperscript{\ensuremath{y}}\fi}
\newunicodechar{ᵃ}{\ifmmode^{a}\else\textsuperscript{\ensuremath{a}}\fi}
\newunicodechar{ᵉ}{\ifmmode^{e}\else\textsuperscript{\ensuremath{e}}\fi}
\newunicodechar{ᵏ}{\ifmmode^{k}\else\textsuperscript{\ensuremath{k}}\fi}
\newunicodechar{ᵖ}{\ifmmode^{p}\else\textsuperscript{\ensuremath{p}}\fi}
\newunicodechar{ᴺ}{\ifmmode^{N}\else\textsuperscript{\ensuremath{N}}\fi}
\newunicodechar{ᶜ}{\ifmmode^{c}\else\textsuperscript{\ensuremath{c}}\fi}
\newunicodechar{ʰ}{\ifmmode^{h}\else\textsuperscript{\ensuremath{h}}\fi}
\newunicodechar{ᵠ}{\ifmmode^{q}\else\textsuperscript{\ensuremath{q}}\fi}
\newunicodechar{ₙ}{\ifmmode_{n}\else\textsubscript{\ensuremath{n}}\fi}
\newunicodechar{ₐ}{\ifmmode_{a}\else\textsubscript{\ensuremath{a}}\fi}
\newunicodechar{ᵢ}{\ifmmode_{i}\else\textsubscript{\ensuremath{i}}\fi}
\newunicodechar{ᵣ}{\ifmmode_{r}\else\textsubscript{\ensuremath{r}}\fi}
\newunicodechar{ₖ}{\ifmmode_{k}\else\textsubscript{\ensuremath{k}}\fi}
\newunicodechar{ᵦ}{\ifmmode_{\beta}\else\textsubscript{\ensuremath{\beta}}\fi}
\newunicodechar{ₓ}{\ifmmode_{x}\else\textsubscript{\ensuremath{x}}\fi}
\newfontfamily\popdisplay{ArchivoBlack-Regular.ttf}[Path=../../../fonts/]
\newfontfamily\poplogo{SpaceGrotesk-Bold.ttf}[Path=../../../fonts/]
\newfontfamily\popsemi{PlusJakartaSans-SemiBold.ttf}[Path=../../../fonts/]
\newfontfamily\popxbold{PlusJakartaSans-ExtraBold.ttf}[Path=../../../fonts/]
\newfontfamily\popxboldls{PlusJakartaSans-ExtraBold.ttf}[Path=../../../fonts/,LetterSpace=3.0]

\setlist[enumerate]{leftmargin=1.7em,itemsep=5pt,topsep=4pt}
\setlist[itemize]{leftmargin=1.7em,itemsep=4pt,topsep=4pt,label={\color{poporange}\textbullet}}
\setlength{\parindent}{0pt}
\setlength{\parskip}{5pt}
\renewcommand{\arraystretch}{1.25}

% pgfplots : style Pop par défaut
\pgfplotsset{
  every axis/.append style={axis line style={popink,line width=1pt},tick style={popink},
    grid style={popline,line width=0.5pt},label style={font=\small},tick label style={font=\footnotesize}},
  popcurve/.style={poporange,line width=1.8pt,no markers,smooth},
}
% Même style disponible dans un \draw TikZ simple (hors pgfplots)
\tikzset{popcurve/.style={poporange,line width=1.8pt}}
\tikzset{popgrid/.style={popline,very thin}}
\colorlet{popcurve}{poporange} % le nom du style est parfois utilisé comme couleur

% ============================================================
% Pill / squiggle
% ============================================================
\newcommand{\poppill}[3]{%
  \tikz[baseline=-0.55ex]\node[draw=popink,line width=1.2pt,rounded corners=2.6pt,%
    fill=#2,inner xsep=6.5pt,inner ysep=3pt]{%
    \popxboldls\fontsize{7}{7}\selectfont\color{#3}\MakeUppercase{#1}};%
}
\newcommand{\popsquiggle}[1]{%
  \tikz[baseline=-0.4ex]{
    \draw[#1,line width=2.6pt,line cap=round]
      plot [smooth] coordinates {(0,0.09) (0.34,0.01) (0.68,0.21) (1.02,0.01) (1.36,0.10)};
  }%
}
% Mot-clé surligné (marqueur orange sous le texte)
\newcommand{\cle}[1]{{\popxbold\color{popdark}#1}}

% ============================================================
% En-tête / pied
% ============================================================
\pagestyle{fancy}
\fancyhf{}
\renewcommand{\headrulewidth}{0pt}
\renewcommand{\footrulewidth}{0pt}
\lhead{\raisebox{-0.15ex}{\poplogo\fontsize{13}{13}\selectfont\color{popink}OnBuch\color{poporange}+}}
\rhead{\poppill{\DOCMATIERE\ \textbullet\ \DOCNIVEAU\ \textbullet\ Leçon \DOCLECON}{white}{popink}}
\lfoot{\popxbold\fontsize{7.6}{7.6}\selectfont\color{popmuted}\MakeUppercase{Cours signé OnBuch+}\ \textbullet\ {\popsemi\DOCTITRE}}
\rfoot{\tikz[baseline=-0.6ex]\node[rounded corners=8.5pt,fill=popink,minimum height=17pt,minimum width=17pt,inner xsep=5pt,inner sep=0]{\popxbold\fontsize{8}{8}\selectfont\color{popcream}\thepage\,/\,\pageref*{LastPage}};}

% ============================================================
% Titres
% ============================================================
\newcommand{\chaptitre}[2]{%
  \begin{tcolorbox}[enhanced,colback=poporange,colframe=popink,boxrule=2.6pt,arc=11pt,
      drop shadow=popink,left=15pt,right=15pt,top=12pt,bottom=11pt,before skip=3pt,after skip=18pt]
    \poppill{#2}{popink}{popcream}\par\vspace{9pt}
    {\raggedright\hyphenpenalty=10000\exhyphenpenalty=10000\popdisplay\fontsize{22}{24}\selectfont\color{popcream}\MakeUppercase{#1}\par}
  \end{tcolorbox}
}
% Section numérotée : pastille orange avec le numéro + titre Archivo
\newcounter{popsec}
\newcommand{\coursec}[1]{%
  \stepcounter{popsec}\setcounter{popsub}{0}%
  \par\vspace{16pt}\noindent
  \begin{minipage}{\linewidth}
  \tikz[baseline=-0.7ex]\node[draw=popink,line width=1.6pt,fill=poporange,rounded corners=4pt,minimum size=22pt,inner sep=0]{\popdisplay\fontsize{12}{12}\selectfont\color{popcream}\thepopsec};%
  \hspace{8pt}{\popdisplay\fontsize{15}{17}\selectfont\color{popink}\MakeUppercase{#1}}\par
  \vspace{4pt}\noindent{\color{poporange}\rule{36pt}{3.6pt}}
  \end{minipage}\par\nopagebreak\vspace{8pt}
}
\newcounter{popsub}
\newcommand{\courssub}[1]{%
  \stepcounter{popsub}%
  \par\vspace{9pt}\noindent{\popxbold\fontsize{11.5}{13}\selectfont\color{popink}{\color{poporange}\thepopsec.\thepopsub}\hspace{6pt}#1}\par\nopagebreak\vspace{4pt}
}

% ============================================================
% Boîtes pédagogiques — même recette : fond blanc, bordure épaisse colorée,
% ombre dure encre, badge plein. Titre optionnel [#1].
% ============================================================
\tcbset{popbase/.style={breakable,enhanced,colback=white,boxrule=2.3pt,arc=9pt,
  shadow={3pt}{-3pt}{0pt}{popink},left=14pt,right=14pt,top=11pt,bottom=11pt,before skip=11pt,after skip=15pt,
  fontupper=\popsemi\fontsize{10.6}{14.5}\selectfont\color{popink},
  fontlower=\popsemi\fontsize{10.6}{14.5}\selectfont\color{popink}}}
% Coche dessinée (le glyphe ✓ n'existe pas dans les polices du document)
\newcommand{\popcheck}[1]{\tikz[baseline=-0.5ex]\draw[#1,line width=1.6pt,line cap=round,line join=round](-0.13,0.0)--(-0.04,-0.09)--(0.13,0.1);}
\providecommand{\coche}{\popcheck{popgreen}}
\providecommand{\resultat}[1]{\par\smallskip\noindent\fcolorbox{popgreen}{popgreenL}{\parbox{\dimexpr\linewidth-2\fboxsep-2\fboxrule\relax}{\textbf{Résultat.} #1}}\par\smallskip}
\newcommand{\popboxhead}[3]{\poppill{#1}{#2}{white}\ifx\relax#3\relax\else\hspace{6pt}{\popxbold\fontsize{10.6}{12}\selectfont\color{popink}#3}\fi\par\vspace{7pt}}

\newtcolorbox{aretenir}[1][]{popbase,colframe=popgreen,before upper={\popboxhead{À retenir}{popgreen}{#1}}}
\newtcolorbox{exemplebox}[1][]{popbase,colframe=poporange,before upper={\popboxhead{Exemple}{poporange}{#1}}}
\newtcolorbox{definition}[1][]{popbase,colframe=popblue,before upper={\popboxhead{Définition}{popblue}{#1}}}
\newtcolorbox{propriete}[1][]{popbase,colframe=poppurple,before upper={\popboxhead{Propriété}{poppurple}{#1}}}
\newtcolorbox{methode}[1][]{popbase,colframe=popgold,before upper={\popboxhead{Méthode}{popgold}{#1}}}
\newtcolorbox{attention}[1][]{popbase,colframe=poppink,before upper={\popboxhead{Attention}{poppink}{#1}}}
\newtcolorbox{experience}[1][]{popbase,colframe=popblue,colback=popblueL,before upper={\popboxhead{Expérience}{popblue}{#1}}}
% « Le savais-tu ? » : carte violette pleine (contexte, culture, Cameroun)
\newtcolorbox{savaistu}[1][]{popbase,colback=poppurple,colframe=popink,
  fontupper=\popsemi\fontsize{10.4}{14.2}\selectfont\color{white},
  before upper={\poppill{Le savais-tu\,?}{white}{poppurple}\ifx\relax#1\relax\else\hspace{6pt}{\popxbold\color{white}#1}\fi\par\vspace{7pt}}}
% Exercice résolu : énoncé en haut, \tcblower puis la solution sur fond crème
\newcounter{popexo}\newcounter{popexr}
\newtcolorbox{exoresolu}[1][]{popbase,colframe=poporange,colbacklower=popcream,
  segmentation style={popink,line width=1.2pt,dashed},
  before upper={\stepcounter{popexr}\popboxhead{Exercice résolu \thepopexr}{poporange}{#1}},
  before lower={\poppill{Solution}{popink}{popcream}\par\vspace{6pt}}}
% Exercice (non résolu en place) — la correction est dans la partie Corrigés
\newtcolorbox{exercice}[1][]{popbase,colframe=popink,
  before upper={\stepcounter{popexo}\popboxhead{Exercice \thepopexo}{popink}{#1}}}
% Corrigé
\newtcolorbox{corrige}[1][]{popbase,colframe=popgreen,colback=popgreenL,
  before upper={\popboxhead{Corrigé}{popgreen}{#1}}}
% Objectifs / prérequis / bilan
\newtcolorbox{objectifs}{popbase,colframe=popink,colback=poporangeL,
  before upper={\popboxhead{Objectifs de la leçon}{popink}{}}}
\newtcolorbox{prerequis}{popbase,colframe=popink,colback=popblueL,
  before upper={\popboxhead{Prérequis}{popblue}{}}}
\newtcolorbox{bilan}{popbase,colframe=popink,colback=white,boxrule=2.8pt,
  before upper={\popboxhead{Fiche bilan}{poporange}{L'essentiel à connaître pour l'examen}}}
% Figure encadrée avec légende
\newtcolorbox{popfigure}[1][]{enhanced,breakable=false,colback=white,colframe=popline,boxrule=1.4pt,arc=8pt,
  left=10pt,right=10pt,top=10pt,bottom=8pt,before skip=10pt,after skip=12pt,halign=center,
  lower separated=false,halign lower=center,
  fontlower=\popsemi\fontsize{9}{11.5}\selectfont\color{popmuted},
  #1}
\newcommand{\legende}[1]{\par\vspace{6pt}{\popsemi\fontsize{9}{11.5}\selectfont\color{popmuted}#1}}

% ============================================================
% Signature OnBuch+
% ============================================================
\newcommand{\poplogoinline}[1]{{\poplogo\fontsize{#1}{#1}\selectfont\color{popink}OnBuch\color{poporange}+}}
\newcommand{\signatureonbuch}{%
  \begin{tcolorbox}[enhanced,colback=popink,colframe=popink,boxrule=2.2pt,arc=10pt,
      drop shadow=poporange,left=14pt,right=14pt,top=10pt,bottom=10pt,before skip=0pt,after skip=0pt]
    \tikz[baseline=-0.7ex]\node[circle,fill=popgreen,draw=popcream,line width=1.3pt,minimum size=22pt,inner sep=0]{\popcheck{white}};%
    \hspace{9pt}\begin{minipage}[c]{0.62\linewidth}
      {\popxbold\fontsize{12}{13}\selectfont\color{popcream}Signé\ \textbullet\ {\poplogo OnBuch\color{poporange}+}}\par\vspace{2pt}
      {\raggedright\popsemi\fontsize{8}{10.5}\selectfont\color{popline}\MakeUppercase{Équipe pédagogique OnBuch+}\\\MakeUppercase{Conforme au programme officiel MINESEC}\par}
    \end{minipage}\hfill
    {\poplogo\fontsize{22}{22}\selectfont\color{popcream}OnBuch\color{poporange}+}
  \end{tcolorbox}%
}

% ============================================================
% Page de garde
% #1 titre · #2 matière · #3 niveau · #4 module · #5 n° leçon · #6 durée ·
% #7 description / objectifs (texte ou itemize)
% ============================================================
\newcommand{\pagegarde}[7]{%
  \thispagestyle{empty}
  \newgeometry{a4paper,margin=1.7cm,top=1.6cm,bottom=1.4cm}%
  \begin{tikzpicture}[remember picture,overlay]
    \fill[poporange!18] ([xshift=-3.2cm,yshift=-3.6cm]current page.north east) circle (4.6cm);
    \fill[poppurple!14] ([xshift=2.4cm,yshift=7.6cm]current page.south west) circle (3.8cm);
  \end{tikzpicture}%
  {\poplogo\fontsize{30}{30}\selectfont\color{popink}OnBuch\color{poporange}+}\hfill
  \raisebox{0.6ex}{\poppill{Cours signé \textbullet\ Premium}{popink}{popcream}}\par
  \vspace{1pt}\popsquiggle{poppurple}\par
  \vspace{8pt}
  \begin{tcolorbox}[enhanced,colback=poporange,colframe=popink,boxrule=2.8pt,arc=13pt,
      drop shadow=popink,left=18pt,right=18pt,top=12pt,bottom=13pt,after skip=10pt]
    \poppill{#2 \textbullet\ #3}{popink}{popcream}\hspace{5pt}\poppill{#4}{white}{popink}\par\vspace{8pt}
    {\popdisplay\fontsize{14}{14}\selectfont\color{popink}LEÇON #5}\par\vspace{4pt}
    {\raggedright\hyphenpenalty=10000\exhyphenpenalty=10000\popdisplay\fontsize{25}{27}\selectfont\color{popcream}\MakeUppercase{#1}\par}
    \vspace{8pt}
    {\popxbold\fontsize{9.5}{9.5}\selectfont\color{popink}\MakeUppercase{Durée conseillée : #6}}
  \end{tcolorbox}
  \begin{tcolorbox}[enhanced,colback=white,colframe=popink,boxrule=2.2pt,arc=10pt,
      drop shadow=popink,left=16pt,right=16pt,top=10pt,bottom=10pt,after skip=8pt]
    \poppill{Dans cette leçon}{poppurple}{white}\par\vspace{5pt}
    \popsemi\fontsize{9.3}{12.6}\selectfont\color{popink}#7
  \end{tcolorbox}
  \noindent\begin{minipage}[t]{0.487\linewidth}
    \begin{tcolorbox}[enhanced,colback=popgreen,colframe=popink,boxrule=2.2pt,arc=10pt,
        drop shadow=popink,left=11pt,right=11pt,top=10pt,bottom=10pt,height=3.1cm,valign=center]
      \tcbox[colback=white,colframe=popink,boxrule=1.3pt,arc=5pt,boxsep=0pt,nobeforeafter,
        left=3pt,right=3pt,top=3pt,bottom=3pt,valign=center]{\qrcode[height=1.9cm]{https://play.google.com/store/apps/details?id=com.luvvix.onbuch&hl=fr}}%
      \hspace{8pt}\begin{minipage}[c]{0.5\linewidth}
        {\popdisplay\fontsize{11}{12.5}\selectfont\color{white}\MakeUppercase{Télécharge OnBuch}}\par\vspace{4pt}
        {\raggedright\popsemi\fontsize{8.4}{11}\selectfont\color{white}Gratuit \textbullet\ Google Play\\Tuteur IA, annales, concours, résultats\par}
      \end{minipage}
    \end{tcolorbox}
  \end{minipage}\hfill
  \begin{minipage}[t]{0.487\linewidth}
    \begin{tcolorbox}[enhanced,colback=poppurple,colframe=popink,boxrule=2.2pt,arc=10pt,
        drop shadow=popink,left=11pt,right=11pt,top=10pt,bottom=10pt,height=3.1cm,valign=center]
      \tikz[baseline=-0.7ex]\node[circle,fill=white,draw=popink,line width=1.3pt,minimum size=1.6cm,inner sep=0]{\popdisplay\fontsize{24}{24}\selectfont\color{poppurple}+};%
      \hspace{8pt}\begin{minipage}[c]{0.6\linewidth}
        {\popdisplay\fontsize{11}{12.5}\selectfont\color{white}\MakeUppercase{Passe à OnBuch+}}\par\vspace{4pt}
        {\raggedright\popsemi\fontsize{8.4}{11}\selectfont\color{white}Tous les cours signés, Tuteur IA illimité, fiches PDF — onglet OnBuch+ de l'app\par}
      \end{minipage}
    \end{tcolorbox}
  \end{minipage}
  \vspace{8pt}
  \popcontenu
  \vfill
  \signatureonbuch
  \restoregeometry
  \newpage
}

% Rangée « Ce cours contient » (page de garde)
\newcommand{\poptile}[3]{% #1 couleur #2 gros texte #3 légende
  \begin{tcolorbox}[enhanced,colback=#1,colframe=popink,boxrule=2pt,arc=9pt,shadow={2.5pt}{-2.5pt}{0pt}{popink},
      left=6pt,right=6pt,top=9pt,bottom=9pt,halign=center,valign=center,height=1.75cm,width=\linewidth,nobeforeafter]
    {\popdisplay\fontsize{12.5}{13}\selectfont\color{white}\MakeUppercase{#2}}\par\vspace{3pt}
    {\centering\popsemi\fontsize{7.8}{9.5}\selectfont\color{white}#3\par}
  \end{tcolorbox}}
\newcommand{\popcontenu}{%
  \par\noindent{\popxboldls\fontsize{7.5}{7.5}\selectfont\color{popmuted}CE COURS SIGNÉ CONTIENT}\par\vspace{7pt}
  \noindent\begin{minipage}[t]{0.235\linewidth}\poptile{popblue}{Cours}{complet, illustré, pas à pas}\end{minipage}\hfill
  \begin{minipage}[t]{0.235\linewidth}\poptile{poporange}{Méthodes}{et exercices résolus}\end{minipage}\hfill
  \begin{minipage}[t]{0.235\linewidth}\poptile{poppink}{Exercices}{type examen + corrigés}\end{minipage}\hfill
  \begin{minipage}[t]{0.235\linewidth}\poptile{popgreen}{Fiche bilan}{l'essentiel à retenir}\end{minipage}\par}

% Sommaire
\newtcolorbox{sommaire}{enhanced,colback=white,colframe=popink,boxrule=2.2pt,arc=10pt,
  drop shadow=popink,left=18pt,right=18pt,top=13pt,bottom=13pt,before skip=6pt,after skip=20pt,
  before upper={\poppill{Sommaire}{poppurple}{white}\par\vspace{9pt}\popsemi\fontsize{10.6}{14.5}\selectfont\color{popink}}}

% Signature de fin de cours — poussée tout en bas de la dernière page
\newcommand{\findecours}{%
  \par\vfill\nopagebreak
  \begin{center}\popsquiggle{poporange}\end{center}\vspace{4pt}
  \signatureonbuch
}
\AtBeginDocument{\def\llbracket{\mathopen{[\mkern-3.5mu[}}\def\rrbracket{\mathclose{]\mkern-3.5mu]}}} % intervalles d'entiers (glyphe ⟦ absent de la police)
% Glyphes absents de Fira Math : redéfinis à partir de glyphes disponibles
\AtBeginDocument{%
  \def\setminus{\mathbin{\backslash}}\let\smallsetminus\setminus
  \def\vdots{\mathord{\vbox{\baselineskip3.2pt\lineskiplimit0pt\kern4pt\hbox{.}\hbox{.}\hbox{.}}}}%
  \def\ddots{\mathinner{\mkern1mu\raise6.5pt\hbox{.}\mkern2mu\raise3.6pt\hbox{.}\mkern2mu\raise0.7pt\hbox{.}\mkern1mu}}%
  \def\triangle{\mathord{\scalebox{0.9}{$\symup{\Delta}$}}}%
  \def\nabla{\mathord{\raisebox{\depth}{\scalebox{1}[-1]{$\symup{\Delta}$}}}}%
  \def\checkmark{\popcheck{popdark}}%
  \def\star{\mathbin{\ast}}\def\bigstar{\mathord{\ast}}%
  \def\diamond{\mathbin{\scalebox{0.6}{$\Diamond$}}}%
  \def\bigcup{\mathop{\mathchoice{\raisebox{-0.3em}{\scalebox{1.7}{$\cup$}}}{\scalebox{1.25}{$\cup$}}{\cup}{\cup}}\displaylimits}%
  \def\bigcap{\mathop{\mathchoice{\raisebox{-0.3em}{\scalebox{1.7}{$\cap$}}}{\scalebox{1.25}{$\cap$}}{\cap}{\cap}}\displaylimits}%
  \def\lesseqqgtr{\mathrel{\lessgtr}}\def\bigoplus{\mathop{\oplus}\displaylimits}%
}
\providecommand{\ssi}{si et seulement si} % abréviation fréquente
\AtBeginDocument{\providecommand{\wideparen}{\overparen}} % arc de cercle

\begin{document}

\pagegarde{Consolidation des acquis de la méthodologie de l'exercice sur texte philosophique}{Philosophie}{Tle Tech}{Philosophie – classe de Terminale technique}{N03}{3 séances (fiche de progression harmonisée)}{Cette leçon de consolidation vise à ancrer durablement la méthodologie de l'exercice sur texte philosophique (épreuve obligatoire du Baccalauréat technique) à travers des mises en situation progressives, l'analyse de copies types et des entraînements ciblés sur les trois questions types (explication, analyse, discussion).
\vspace{4pt}
\begin{multicols}{2}\small
\begin{itemize}
\item À la fin de cette leçon, je sais identifier avec certitude les attendus précis de chacune des trois questions de l'épreuve de texte (Q1 : explicitation, Q2 : analyse/argumentation, Q3 : discussion/critique).
\item Je sais structurer une réponse de Q1 en distinguant le sens littéral, les concepts clés et l'articulation logique de l'extrait sans faire de hors-sujet ni de paraphrase.
\item Je sais construire une réponse de Q2 en dégageant la thèse, les arguments de l'auteur et leur enchaînement logique, en utilisant le vocabulaire technique approprié (thèse, argument, présupposé, concept).
\item Je sais rédiger une Q3 (discussion) en adoptant une posture critique structurée (concession, limite, contre-argumentation, ouverture) et en mobilisant des références philosophiques et des exemples ancrés dans le réel camerounais.
\item Je sais gérer mon temps d'épreuve (3h) en allouant durées proportionnées à chaque question (environ 45min, 45min, 90min) et en réservant 15min à la relecture.
\item Je sais éviter les pièges classiques : confusion explication/commentaire, absence de thèse en Q2, discussion hors sujet ou sans nuance en Q3, style oral ou approximatif.
\end{itemize}\end{multicols}}

\begin{sommaire}
\begin{multicols}{2}
\begin{enumerate}
\item 1. Architecture de l'épreuve et gestion stratégique du temps (3h)
\item 2. Maîtrise de la Question 1 : L'explicitation rigoureuse du sens (Savoir)
\item 3. Maîtrise de la Question 2 : L'analyse de l'argumentation (Savoir-faire)
\item 4. Maîtrise de la Question 3 : La discussion critique structurée (Compétence supérieure)
\item 5. Atelier de consolidation : Simulation Bac \& Auto-évaluation croisée
\item 6. Transfert \& Citoyenneté : L'esprit critique hors les murs
\item Activité d'intégration
\item Méthodes et exercices résolus
\item Exercices
\item Corrigés des exercices
\item Fiche bilan
\end{enumerate}
\end{multicols}
\end{sommaire}

\begin{objectifs}
\begin{itemize}
\item À la fin de cette leçon, je sais identifier avec certitude les attendus précis de chacune des trois questions de l'épreuve de texte (Q1 : explicitation, Q2 : analyse/argumentation, Q3 : discussion/critique).
\item Je sais structurer une réponse de Q1 en distinguant le sens littéral, les concepts clés et l'articulation logique de l'extrait sans faire de hors-sujet ni de paraphrase.
\item Je sais construire une réponse de Q2 en dégageant la thèse, les arguments de l'auteur et leur enchaînement logique, en utilisant le vocabulaire technique approprié (thèse, argument, présupposé, concept).
\item Je sais rédiger une Q3 (discussion) en adoptant une posture critique structurée (concession, limite, contre-argumentation, ouverture) et en mobilisant des références philosophiques et des exemples ancrés dans le réel camerounais.
\item Je sais gérer mon temps d'épreuve (3h) en allouant durées proportionnées à chaque question (environ 45min, 45min, 90min) et en réservant 15min à la relecture.
\item Je sais éviter les pièges classiques : confusion explication/commentaire, absence de thèse en Q2, discussion hors sujet ou sans nuance en Q3, style oral ou approximatif.
\item Je sais auto-évaluer ma copie à l'aide d'une grille de correction type jury (cohérence, pertinence, rigueur conceptuelle, qualité de la langue).
\item Je sais transférer ces compétences méthodologiques à l'analyse de tout document philosophique (article de presse, extrait de discours, proverbe) dans une démarche d'esprit critique citoyen.
\end{itemize}
\end{objectifs}

\begin{prerequis}
\begin{itemize}
\item Connaissance des grands concepts du programme de Terminale (Sujet, Liberté, Société, Histoire, Science, Technique, Langage).
\item Maîtrise de la structure standard d'une dissertation philosophique (introduction, développement en parties équilibrées, conclusion) acquise en Première.
\item Notions de base de logique formelle : distinction concept/proposition/argument, types de raisonnement (déduction, induction, analogie), principales fautes logiques (pétition de principe, homme de paille).
\item Habitude de la lecture active : surlignage, annotation, repérage des connecteurs logiques (donc, cependant, car, or...).
\item Culture générale minimale sur le contexte camerounais (institutions, enjeux sociaux, diversité culturelle, actualité récente).
\end{itemize}
\end{prerequis}

\newpage

\coursec{Architecture de l'épreuve et gestion stratégique du temps (3h)}

Au lycée technique de Bafoussam, un élève brillant en série F4 rate l'épreuve de texte philosophique au Bac parce qu'il a confondu « expliquer le sens » et « commenter la thèse », perdant ainsi 6 points sur 20. À Douala, une candidate en STT réussit brillamment l'explication de texte mais échoue à la question de discussion pour avoir manqué d'exemples concrets tirés du vécu camerounais. Ces deux situations révèlent que la maîtrise du cours ne suffit pas : c'est la rigueur méthodologique qui fait la différence entre l'admis et le recalé.

\emph{Problématique : Comment transformer les 180 minutes de l'épreuve en un levier de performance maximale en maîtrisant l'architecture de l'examen, la grille d'évaluation et une gestion stratégique du temps ?}

\courssub{Présentation officielle de l'épreuve : coefficient, durée, nature du support}

L'épreuve de philosophie au Probatoire et au Baccalauréat technique (toutes séries confondues : F, STT, G, etc.) est une \cle{épreuve écrite} d'une durée de \textbf{3 heures (180 minutes)}. Elle porte un \textbf{coefficient 2} (souvent coefficient 3 selon les séries et les sessions, vérifiez votre convocation), ce qui en fait une épreuve déterminante pour l'obtention du diplôme.

Le support est un \cle{texte philosophique} unique, d'une longueur comprise entre \textbf{400 et 600 mots}. L'auteur peut être un \textbf{classique} (Platon, Aristote, Descartes, Kant, Marx, Bergson, Sartre, etc.) ou un \textbf{contemporain} (auteurs africains comme Fabien Eboussi Boulaga, Paulin Hountondji, ou occidentaux récents). Le texte n'est jamais un résumé de cours : c'est un extrait d'œuvre où l'auteur \emph{argumente} pour défendre une \cle{thèse}.

\begin{definition}[Les quatre piliers de l'analyse textuelle]
\textbf{Thèse} : La position précise que l'auteur défend en réponse à un problème philosophique. Ce n'est pas un avis, c'est une réponse argumentée. \par
\textbf{Argument} : Une raison logique ou une preuve invoquée pour soutenir la thèse. Un texte en contient plusieurs, articulés entre eux. \par
\textbf{Concept} : Une notion opératoire (ex. : liberté, conscience, justice, technique, État) que l'auteur utilise avec un sens technique précis, souvent différent du sens commun. \par
\textbf{Présupposé} : Une idée non dite, non démontrée dans le texte, mais nécessaire pour que l'argumentation tienne (ex. : « L'homme est bon par nature » chez Rousseau).
\end{definition}

\courssub{Découpage temporel type : La stratégie des 180 minutes}

Ne subissez pas le temps, \textbf{gelez-le}. Voici le découpage optimal validé par les correcteurs expérimentés. Respectez-le à la minute près lors de vos entraînements.

\begin{popfigure}
\centering
\begin{tikzpicture}[
    scale=0.85,
    task/.style={draw=popdark, thick, rectangle, rounded corners=3pt, minimum height=1.2cm, font=\small\bfseries, text width=3.2cm, align=center},
    buffer/.style={draw=poporange, thick, dashed, rectangle, rounded corners=3pt, minimum height=1.2cm, font=\small\itshape, text width=1.8cm, align=center, fill=poporange!10},
    axis/.style={->, >=Stealth, thick, popdark},
    time/.style={font=\small\bfseries, popblue},
    label/.style={font=\tiny, popmuted, anchor=north},
]
% Axes temps
\draw[axis] (0,0) -- (19,0) node[right, label] {Temps (minutes)};
\foreach \x/\t in {0/0, 15/15, 30/30, 45/45, 60/60, 75/75, 90/90, 105/105, 120/120, 135/135, 150/150, 165/165, 180/180}
    \draw (\x, -0.1) -- (\x, 0.1) node[below, label] {\t};

% Blocs principaux
\node[task, fill=popblue!20, align=center] (L1) at (1.5, 2){Lecture 1\\Repérage global\\(10 min)};
\node[task, fill=popblue!20, align=center] (L2) at (1.5, 0.5){Lecture 2\\Analyse fine\\(15 min)};
\node[task, fill=popblue!20, align=center] (L3) at (1.5, -1){Lecture 3\\Vérification\\(5 min)};
\node[buffer, align=center] (B1) at (1.5, -2.5){Tampon\\Lecture};

\node[task, fill=popgreen!20, align=center] (Q1) at (6.5, 2){Brouillon Q1\\Plan détaillé\\(10 min)};
\node[task, fill=popgreen!20, align=center] (Q1w) at (6.5, 0.5){Rédaction Q1\\Explicitation\\(35 min)};
\node[buffer] (B2) at (6.5, -1) {Tampon Q1};

\node[task, fill=poporange!20, align=center] (Q2b) at (12, 2){Brouillon Q2\\Structure args\\(10 min)};
\node[task, fill=poporange!20, align=center] (Q2w) at (12, 0.5){Rédaction Q2\\Analyse args\\(35 min)};
\node[buffer] (B3) at (12, -1) {Tampon Q2};

\node[task, fill=poppurple!20, align=center] (Q3b) at (16.5, 2.5){Brouillon Q3\\Problématique,\\Plan dialectique,\\Exemples CM\\(20 min)};
\node[task, fill=poppurple!20, align=center] (Q3w) at (16.5, 0.5){Rédaction Q3\\Discussion\\(65 min)};
\node[buffer] (B4) at (16.5, -1.5) {Tampon Q3};

\node[task, fill=poppink!20, align=center] (REL) at (18.5, -2.5){Relecture\\Globale\\(10 min)};

% Flèches de liaison
\draw[axis, popmuted] (L1.east) -- ++(1,0) |- (Q1.west);
\draw[axis, popmuted] (Q1w.east) -- ++(1,0) |- (Q2b.west);
\draw[axis, popmuted] (Q2w.east) -- ++(1,0) |- (Q3b.west);
\draw[axis, popmuted] (Q3w.east) -- ++(1,0) |- (REL.west);

% Légendes phases
\node[font=\bfseries\large, popblue, anchor=west] at (0, 3.5) {PHASE 1 : APPROPRIATION (30 min)};
\node[font=\bfseries\large, popgreen, anchor=west] at (0, 1.5) {PHASE 2 : Q1 - EXPLICITATION (45 min)};
\node[font=\bfseries\large, poporange, anchor=west] at (0, -0.5) {PHASE 3 : Q2 - ARGUMENTATION (45 min)};
\node[font=\bfseries\large, poppurple, anchor=west] at (0, -2.5) {PHASE 4 : Q3 - DISCUSSION (90 min)};
\node[font=\bfseries\large, poppink, anchor=west] at (0, -4) {PHASE 5 : RELECTURE (10 min)};
\end{tikzpicture}
\legende{Diagramme de Gantt : Répartition optimale des 180 minutes. Les zones « Tampon » (pointillés orange) absorbent les dépassements imprévus. Ne les sacrifiez jamais.}
\end{popfigure}

\begin{methode}[La « Lecture active en 3 temps »]
\textbf{Temps 1 : Repérage global (10 min).} Lisez d'une traite, crayon en main. Objectif : identifier l'\cle{auteur} (si connu), le \cle{thème} général (la justice, la technique, la liberté...), la \cle{thèse probable} (pour ou contre ?), le \cle{ton} (dogmatique, ironique, dialogué, phénoménologique) et la \cle{structure paragraphique} (combien de paragraphes ? transition ?). Soulignez les \textbf{mots-clés} récurrents.
\textbf{Temps 2 : Analyse fine (15 min).} Relisez paragraphe par paragraphe. Numérotez les \cle{arguments} (Arg 1, Arg 2...). Entourez les \cle{connecteurs logiques} (donc, or, mais, cependant, en effet). Repérez les \cle{concepts} techniques (définis ou utilisés implicitement). Identifiez les \cle{présupposés}. Sur le brouillon, schématisez : \emph{Thèse $\leftarrow$ Arg 1 (Concept A) + Arg 2 (Concept B) + ...}.
\textbf{Temps 3 : Vérification croisée (5 min).} Relisez une dernière fois \emph{uniquement} pour vérifier : « Ai-je bien compris le sens de \emph{ce} mot dans \emph{ce} contexte ? », « La thèse que j'ai notée correspond-elle vraiment à la conclusion du texte ? », « Y a-t-il un contre-argument intégré dans le texte (concession) ? ».
\end{methode}

\begin{attention}[Erreurs fatales de la première heure]
\begin{itemize}
    \item \textbf{Lire une seule fois et écrire.} C'est l'erreur n°1. Le texte résiste à la lecture unique. La Q1 exige une fidélité que seule la 3e lecture garantit.
    \item \textbf{Confondre Thème et Thèse.} Le thème est large (« La technique »), la thèse est précise (« La technique aliène l'homme en le transformant en rouage » — Ellul). Noter « Thème : La technique » en Q1 vaut 0 point sur la justesse de la thèse.
    \item \textbf{Prendre des notes sur la copie d'examen.} Tout se passe sur le \cle{brouillon} (fourni ou papier personnel). La copie finale doit être propre, structurée, sans ratures.
\end{itemize}
\end{attention}

\courssub{Analyse de la grille d'évaluation type jury MINESEC}

Comprendre comment le correcteur note change votre façon d'écrire. La grille officielle (non publiée mais stable) pondère ainsi sur 20 points :

\begin{center}
\begin{tabularx}{\linewidth}{|l|X|X|c|}\hline
\rowcolor{popblueL}
\textbf{Question} & \textbf{Compétence visée} & \textbf{Critères majeurs (Points forts / Faibles)} & \textbf{Pondération approx.} \\ \hline
\textbf{Q1} & \cle{Savoir} (Explicitation) & \textbf{Justesse} de la thèse, \textbf{précision} des concepts, \textbf{fidélité} au texte (pas d'apport extérieur), \textbf{clarté} de l'exposé. & \textbf{4 à 5 pts} \\ \hline
\textbf{Q2} & \cle{Savoir-faire} (Analyse) & \textbf{Exhaustivité} des arguments, \textbf{rigueur} de l'enchaînement logique, \textbf{identification} des concepts/présupposés, \textbf{qualité} de la reformulation. & \textbf{6 à 7 pts} \\ \hline
\textbf{Q3} & \cle{Compétence supérieure} (Discussion) & \textbf{Problématisation} (dépasser le texte), \textbf{structure dialectique} (Thèse/Antithèse/Synthèse), \textbf{pertinence des exemples} (dont camerounais), \textbf{philosophicité} du raisonnement, \textbf{langue}. & \textbf{8 à 9 pts} \\ \hline
\end{tabularx}
\end{center}

\begin{exemplebox}[Statistiques Bac Session 2023 — Séries Techniques (Données académiques anonymisées)]
\begin{center}
\begin{tabular}{|c|c|c|c|}\hline
\rowcolor{popblueL}
\textbf{Question} & \textbf{Moyenne nationale} & \textbf{Écart-type} & \textbf{Taux de réussite ($>$ 50\%)} \\ \hline
Q1 (Explicitation) & 3,2 / 5 & 1,1 & 68\% \\ \hline
Q2 (Argumentation) & 4,1 / 7 & 1,4 & 55\% \\ \hline
Q3 (Discussion) & 5,5 / 8 & 1,8 & 42\% \\ \hline
\textbf{Total} & \textbf{12,8 / 20} & -- & -- \\ \hline
\end{tabular}
\end{center}
\emph{Analyse :} L'écart se creuse dramatiquement sur la Q3. Les candidats maîtrisent l'explicitation (cours appris), mais échouent à \emph{produire} une pensée critique structurée et illustrée. C'est là que se joue l'admission.
\end{exemplebox}

\courssub{Différence fondamentale avec la dissertation : Le « contrat de fidélité »}

\begin{center}
\begin{tabularx}{\linewidth}{|l|X|X|}\hline
\rowcolor{popblueL}
\textbf{Critère} & \textbf{Dissertation} & \textbf{Explication de texte} \\ \hline
\textbf{Sujet} & Question ou citation courte (libre) & Texte long imposé (contraint) \\ \hline
\textbf{Guide} & Plan libre (dialectique, analytique...) & 3 Questions \textbf{imposées} et \textbf{ordonnées} \\ \hline
\textbf{Q1 / Q2} & N'existent pas & \textbf{Obligation de fidélité stricte} : \emph{« Que dit le texte ? »}, \emph{« Comment le dit-il ? »}. Hors-texte = Hors-sujet. \\ \hline
\textbf{Q3} & Corps de la dissertation & \textbf{Discussion} : \emph{« Qu'en pensez-vous ? »}. Liberté \textbf{après} fidélité. On quitte le texte pour le confronter. \\ \hline
\textbf{Exemples} & Illustres la pensée du candidat & \textbf{Q1/Q2 : Interdits} (sauf définition de concept). \textbf{Q3 : Obligatoires} (propres + camerounais). \\ \hline
\end{tabularx}
\end{center}

\cle{Règle d'or} : En Q1 et Q2, vous êtes \textbf{l'interprète fidèle} de l'auteur. En Q3, vous devenez \textbf{le philosophe critique} qui dialogue avec lui.

\courssub{Utilisation efficace du brouillon : L'architecture invisible de la réussite}

Le brouillon n'est pas une « copie sale ». C'est le \textbf{chantier de l'ingénieur}. Un brouillon réussi garantit une copie finale fluide. Voici la structure type à reproduire mécaniquement :

\begin{popfigure}
\centering
\begin{tikzpicture}[
    mindmap,
    concept color=popblue,
    level 1 concept/.append style={level distance=4cm, sibling angle=90, font=\small\bfseries},
    level 2 concept/.append style={level distance=3cm, sibling angle=45, font=\tiny},
    level 3 concept/.append style={level distance=2.5cm, font=\tiny},
    every node/.style={scale=0.85},
    text=popdark,
    edge from parent/.style={draw=popline, thick}
]
\node[concept, fill=popblue!30] {BROUILLON MAÎTRE}
    child[concept color=popgreen] {node[concept] {LECTURES (30')}
        child {node {Auteur / Date / Contexte}}
        child {node {Thème / Thèse (1 phrase)}}
        child {node {Structure : §1, §2, §3...}}
        child {node {Concepts clés (liste)}}
        child {node {Présupposés repérés}}}
    child[concept color=poporange] {node[concept] {Q1 : EXPLICITATION (45')}
        child {node {Intro : Accroche, Auteur, Thèse, Plan}}
        child {node {§1 : Thèse détaillée}}
        child {node {§2 : Concept A expliqué}}
        child {node {§3 : Concept B expliqué}}
        child {node {Concl : Portée du sens}}}
    child[concept color=poppurple] {node[concept] {Q2 : ARGUMENTATION (45')}
        child {node {Intro : Thèse + Annonce args}}
        child {node {Arg 1 : Énoncé + Analyse + Concept}}
        child {node {Arg 2 : Énoncé + Analyse + Concept}}
        child {node {Arg 3 : Énoncé + Analyse + Concept}}
        child {node {Articulation logique (Enchaînement)}}}
    child[concept color=poppink] {node[concept] {Q3 : DISCUSSION (90')}
        child {node {Problématique (dépassement)}}
        child {node {Plan Dialectique : Thèse / Antithèse / Synthèse}}
        child {node {Ex 1 : Philo classique}}
        child {node {Ex 2 : Vécu Camerounais (OBLIGATOIRE)}}
        child {node {Ex 3 : Actu / Tech / Science}}
        child {node {Intro / Concl rédigées}}};
\end{tikzpicture}
\legende{Arbre de décision : Structure type du brouillon. Chaque branche correspond à une phase temporelle. La branche Q3 est la plus fournie : c'est là que se gagnent les points décisifs.}
\end{popfigure}

\begin{exoresolu}[Simulation de brouillon : Extrait fictif type Bac]
\textbf{Extrait (simulé) :} « La technique n'est pas un simple outil neutre. Elle structure notre rapport au monde. En automatisant la production, elle libère l'homme de la peine, mais l'enferme dans la consommation. Le progrès technique devient ainsi une promesse de bonheur qui se retourne en aliénation. » (Inspiré de H. Marcuse / J. Ellul).

\textbf{Brouillon Q1 (10 min) :}
\begin{itemize}
    \item \textbf{Thèse} : La technique n'est pas neutre ; elle aliène par la consommation ce qu'elle libère par l'automatisation.
    \item \textbf{Concepts} : Neutralité (refusée), Automatisation, Aliénation, Progrès, Bonheur (promesse trompeuse).
    \item \textbf{Plan Q1} : 1. Refus de la neutralité (structure monde). 2. Double mouvement : libération peine / enfermement conso. 3. Renversement du progrès en aliénation.
\end{itemize}

\textbf{Brouillon Q2 (10 min) :}
\begin{itemize}
    \item \textbf{Arg 1} : « Structure notre rapport au monde » $\rightarrow$ Argument ontologique : la technique médiatise l'être-au-monde (Concept : Médiation technique).
    \item \textbf{Arg 2} : « Automatisant... libère de la peine » $\rightarrow$ Argument positif (concession) : gain de temps/énergie.
    \item \textbf{Arg 3} : « Mais l'enferme dans la consommation » $\rightarrow$ Argument critique principal : création de besoins artificiels (Concept : Société de consommation / Fétichisme).
    \item \textbf{Arg 4} : « Promesse... se retourne » $\rightarrow$ Argument dialectique : dialectique du progrès (maîtrise nature $\rightarrow$ asservissement homme).
    \item \textbf{Présupposé} : L'homme est un être de besoins ; le bonheur vrai n'est pas la satisfaction illimitée des besoins.
\end{itemize}

\textbf{Brouillon Q3 (20 min) :}
\begin{itemize}
    \item \textbf{Problématique} : La technique est-elle nécessairement aliénante ou peut-elle être réappropriée comme instrument d'émancipation ?
    \item \textbf{Thèse (Texte)} : Aliénation structurelle (Marcuse/Ellul). Ex : Ouvrier chaîne montage / Smartphone addiction (Douala/Yaoundé).
    \item \textbf{Antithèse} : Technique libératrice / neutre. Ex : Machinisme agricole (Nord-Cameroun : gain temps $\rightarrow$ éducation enfants) ; Médecine technique (vaccins, IRM) ; Logiciels libres / Open Source (réappropriation).
    \item \textbf{Synthèse} : Ni neutralité absolue, ni déterminisme fatal. La technique est un \emph{pharmakon} (remède/poison). L'enjeu : \textbf{gouvernance technique} (démocratie technique, éthique de l'IA, choix énergétiques au Cameroun).
    \item \textbf{Exemples Camerounais sélectionnés} : 1. Poussée de l'orpaillage artisanal vs industriel (East Region) : technique destructrice. 2. Mobile Money (Orange Money/MTN MoMo) : technique d'inclusion financière (émancipation). 3. Barrage de Nachtigal/Lom Pangar : technique énergie vs déplacement populations.
\end{itemize}
\tcblower
\textbf{Solution détaillée / Analyse du correcteur :}
Ce brouillon montre la \textbf{hiérarchisation} : Q1/Q2 restent collés au texte (citations implicites). Q3 \textbf{s'envole} avec une problématique qui déborde le texte (« réappropriation », « émancipation »), un plan dialectique clair, et surtout \textbf{trois exemples camerounais distincts} (économie informelle, fintech, grand ouvrage) qui ancrent la philosophie dans le réel. C'est ce « ancrage local » que le jury MINESEC valorise de plus en plus.
\end{exoresolu}

\courssub{Stratégie de la première lecture : Le « Scanner cognitif »}

Dès la distribution du sujet, activez votre \cle{scanner cognitif} en 3 passes rapides (intégrées dans les 10 min de Lecture 1) :

\begin{enumerate}
    \item \textbf{Identification de l'Auteur/Époque} : Nom connu ? $\rightarrow$ Activez le « fichier auteur » (thèses majeures, concepts fétiches, vocabulaire). Inconnu ? $\rightarrow$ Indices internes : vocabulaire (phénoménologique ? marxiste ? analytique ?), style (rigoureux ? littéraire ?), références citées. Cela oriente l'interprétation des concepts.
    \item \textbf{Repérage de la Thèse (Thèse = Conclusion)} : Cherchez la phrase qui résout la tension. Souvent en début (thèse d'emblée) ou fin (thèse aboutie). Mots-clés : \emph{« Il faut conclure que... », « En définitive... », « C'est pourquoi... », « Ainsi... »}.
    \item \textbf{Cartographie Paragraphique} : Numérotez les paragraphes. Pour chacun : \emph{Une fonction} (Thèse, Argument 1, Objection, Réponse, Exemple, Conclusion). Tracez des flèches entre paragraphes : $\rightarrow$ (enchaînement), $\dashrightarrow$ (opposition/concession), $\rightsquigarrow$ (approfondissement).
\end{enumerate}

\begin{savaistu}[Le secret du correcteur : L'introduction de la Q3]
Le correcteur lit souvent \textbf{en premier l'introduction de la Question 3}. Pourquoi ? C'est là que se révèle la \emph{maturité philosophique} : la capacité à \textbf{problématiser} (transformer le thème du texte en question philosophique ouverte) et à \textbf{annoncer un plan dialectique maîtrisé}. Une intro Q3 brillante (problématique claire, plan net, enjeu annoncé) crée un \emph{a priori positif} qui rejaillit sur la notation de Q1 et Q2 (bienveillance sur la fidélité). Une intro Q3 faible (« Je vais discuter le texte... ») signale un candidat en difficulté, incitant à la sévérité sur les approximations de Q1/Q2. \textbf{Soignez votre intro Q3 comme votre visage au jury.}
\end{savaistu}

\begin{exercice}[Entraînement chronométré « Bac Blanc »]
Prenez un sujet zéro ou un sujet de session antérieure (disponibles sur OnBuch+ ou site MINESEC).
\begin{enumerate}
    \item Mettez un minuteur sur \textbf{180 minutes exactes}.
    \item Appliquez \textbf{strictement} le diagramme de Gantt (Section 1.2).
    \item Interdiction d'écrire sur la copie finale avant la minute 30 (fin Phase 1).
    \item À la fin, notez votre ressenti : « Ai-je manqué de temps où ? » « Mon brouillon était-il assez structuré ? »
    \item Faites corriger par un pair ou un prof avec la grille de la Section 1.3.
\end{enumerate}
\emph{Objectif :} Automatiser la gestion du temps pour que le jour J, le cerveau soit 100\% disponible pour la \emph{pensée}, pas pour la \emph{montre}.
\end{exercice}

\coursec{Maîtrise de la Question 1 : L'explicitation rigoureuse du sens (Savoir)}

Dans la section précédente, nous avons appris à \emph{architecturer} notre épreuve : lire le texte deux fois, répartir les trois heures, réserver environ 45 minutes à la première question. Mais savoir gérer son temps ne suffit pas : encore faut-il savoir quoi écrire pendant ces 45 minutes. C'est l'objet de cette section. La Question 1 est la \cle{porte d'entrée} de l'exercice sur texte : elle vérifie un \textbf{savoir}, c'est-à-dire votre capacité à comprendre rigoureusement ce que l'auteur dit et veut dire. Une Q1 ratée condamne presque toujours la copie entière, car les questions 2 et 3 s'appuient sur elle. Inversement, une Q1 réussie installe la confiance du correcteur et prépare le terrain de toute la suite.

\courssub{Comprendre la consigne : ce que l'examinateur attend vraiment}

\textbf{Les formulations types de la Question 1}\par

Au Probatoire et au Baccalauréat technique, la Question 1 se présente sous des formes voisines :

\begin{itemize}
\item \og Expliquez le sens du texte. \fg
\item \og Dégagez la thèse de l'auteur et les concepts principaux. \fg
\item \og Quel est le problème posé par le texte et quelle réponse l'auteur y apporte-t-il ? \fg
\end{itemize}

Ces trois formulations demandent la même compétence : \textbf{expliciter}, c'est-à-dire rendre explicite ce qui est plié (\emph{ex-plicare} : dé-plier) dans le texte. L'examinateur ne vous demande ni de répéter le texte, ni de le juger, ni de parler du thème en général. Il vous demande de montrer que vous avez \emph{compris} : la \cle{thèse} (ce que l'auteur soutient), les \cle{concepts} (les outils de pensée qu'il utilise) et l'\cle{enchaînement logique} (comment les idées s'articulent).

\begin{attention}[La consigne cache un contrat]
Derrière \og expliquez le sens du texte \fg, le correcteur lit : \og prouvez-moi que vous savez ce qu'est un concept philosophique, que vous savez définir les termes de l'auteur, et que vous voyez la logique interne du passage \fg. Toute réponse qui ne remplit pas ce contrat, même bien écrite, est hors sujet.
\end{attention}

\textbf{Sens littéral et sens philosophique : la distinction cruciale}\par

Voici l'intuition de départ. Quand un ami vous dit \og il pleut des cordes \fg, vous comprenez deux choses : le sens littéral (des gouttes d'eau tombent) et le sens réel (il pleut très fort, peut-être faut-il remettre la sortie au champ). En philosophie, cette distinction devient décisive.

\begin{definition}[Sens littéral et sens philosophique]
Le \cle{sens littéral} est ce que le texte dit explicitement, phrase par phrase : les mots, les exemples, les affirmations de surface. Le \cle{sens philosophique} est ce que le texte \emph{veut dire} : la thèse défendue, le problème traité, les implications de pensée (ce qui suit nécessairement des affirmations de l'auteur). Expliciter, c'est passer du premier au second sans jamais trahir le premier.
\end{definition}

\begin{exemplebox}[Arendt et la technique]
Supposons qu'un texte d'Hannah Arendt affirme : \og l'homme moderne ne fabrique plus des outils pour le monde, il fabrique un monde pour ses outils \fg. Le sens littéral : une inversion entre outil et monde. Le sens philosophique : la technique moderne a cessé d'être un simple moyen au service de fins humaines ; elle impose désormais ses propres fins et reconfigure l'existence humaine (ce qu'on appellera la \cle{raison instrumentale} devenue autonome). C'est ce second niveau que la Q1 doit atteindre.
\end{exemplebox}

\courssub{Le travail sur les concepts : le cœur du savoir}

\textbf{Qu'est-ce qu'un concept philosophique ?}\par

\begin{definition}[Concept philosophique]
Un \cle{concept philosophique} est un outil de pensée opératoire : une idée générale, construite et définie rigoureusement, qui permet de penser un problème. Il ne se confond pas avec le mot du dictionnaire courant. Ainsi \og aliénation \fg ne signifie pas seulement \og être dépossédé \fg, mais désigne le processus par lequel le sujet devient étranger à lui-même, à son travail ou à sa communauté ; \og technicité \fg désigne le caractère propre de ce qui relève de la technique ; \og communauté \fg désigne un mode d'être-ensemble fondé sur l'appartenance et non sur le simple intérêt.
\end{definition}

L'erreur de l'élève pressé est de définir les mots du texte avec le sens du marché ou de la rue. Or le correcteur attend la définition \emph{philosophique}. Si le texte de Fabien Eboussi Boulaga parle de \og l'existence \fg, il ne s'agit pas du fait banal d'être vivant, mais de la manière dont un sujet se pose, assume sa condition et s'interroge sur son être dans une situation historique donnée — ici, la situation camerounaise et africaine postcoloniale.

\textbf{Articuler les concepts entre eux}\par

Définir les concepts ne suffit pas : il faut montrer leurs \textbf{relations}. Trois types de relations reviennent constamment :

\begin{itemize}
\item \textbf{Hiérarchie} : un concept est subordonné à un autre (la technique comme \emph{moyen} subordonné à une \emph{fin}).
\item \textbf{Opposition} : deux concepts s'excluent ou se tensionnent (raison instrumentale / raison communicative chez Habermas ; être / paraître).
\item \textbf{Complémentarité} : deux concepts se définissent l'un par l'autre (liberté et responsabilité ; individu et communauté).
\end{itemize}

\begin{propriete}[Principe de fidélité textuelle]
Toute affirmation produite en Q1 doit être étayée par une référence précise au texte : un mot, une tournure, une ligne citée brièvement entre guillemets, puis \emph{expliquée}. On n'affirme jamais \og l'auteur pense que... \fg sans montrer \emph{où} et \emph{comment} le texte le dit. Ce principe est la règle d'or de l'exercice : il distingue l'explication de texte de l'improvisation.
\end{propriete}

\courssub{La méthode de préparation sur brouillon}

Avant de rédiger, il faut extraire la matière du texte. Voici la technique professionnelle, à réaliser au brouillon en cinq minutes maximum.

\begin{methode}[La technique \og Mots-clés / Concepts / Liens logiques \fg]
\begin{enumerate}
\item \textbf{Repérer les mots-clés} : soulignez dans le texte les termes porteurs de sens (noms abstraits, verbes forts, répétitions). Ne retenez que 4 à 6 mots.
\item \textbf{Traduire en concepts} : dans une deuxième colonne, transformez chaque mot-clé en concept défini philosophiquement en une phrase.
\item \textbf{Identifier les liens logiques} : dans une troisième colonne, notez les connecteurs et articulations du texte (\og car \fg, \og mais \fg, \og donc \fg, \og si... alors \fg) et ce qu'ils établissent (cause, opposition, conséquence).
\item \textbf{Dégager la thèse} : formulez en une seule phrase ce que l'auteur soutient, en combinant les concepts et les liens.
\item \textbf{Découper en parties} : le tableau révèle naturellement 2 ou 3 blocs de sens ; ce seront vos parties.
\end{enumerate}
\end{methode}

\begin{exemplebox}[Application de la grille à un extrait]
Pour un extrait où Eboussi Boulaga écrit que l'Africain a longtemps reçu son image de l'extérieur et doit désormais se penser lui-même, la grille donne :

\begin{center}
\begin{tabularx}{\linewidth}{|l|X|X|}
\hline
\rowcolor{popblueL} \textbf{Mots-clés} & \textbf{Concepts définis} & \textbf{Liens logiques} \\
\hline
\og image reçue \fg & \cle{Aliénation} : être défini par le regard d'autrui & opposition (\og longtemps... désormais \fg) \\
\hline
\og se penser soi-même \fg & \cle{Subjectivité} : capacité du sujet à se définir & conséquence (\og doit \fg : exigence) \\
\hline
\og existence \fg & \cle{Existence} : manière assumée d'être-au-monde & complémentarité (existence et subjectivité) \\
\hline
\end{tabularx}
\end{center}

La thèse s'en déduit : l'existence authentique exige de rompre avec l'aliénation du regard extérieur pour accéder à une subjectivité assumée.
\end{exemplebox}

\courssub{La structure type d'une réponse réussie}

\textbf{Les trois moments de la réponse}\par

Une Q1 n'est pas un bloc informe : c'est une mini-dissertation \emph{interne au texte}, en trois moments.

\textbf{1. L'introduction (3 à 5 lignes).} Elle comprend : une \textbf{accroche contextuelle} (situer brièvement l'auteur et le problème : \og Dans cet extrait, Hannah Arendt s'interroge sur le rapport de l'homme moderne à la technique \fg) ; la \textbf{citation de la thèse} (formulée clairement : \og l'auteur soutient que... \fg) ; l'\textbf{annonce du plan} (\og nous verrons d'abord que..., ensuite que..., enfin que... \fg).

\textbf{2. Le développement (2 ou 3 parties).} Chaque partie suit la structure interne du texte : une partie = un mouvement de pensée de l'auteur. Chaque partie définit un concept, cite brièvement le texte, explique la citation, et montre le lien avec la partie suivante.

\textbf{3. La conclusion (2 à 3 lignes).} Elle synthétise le sens global (\og ainsi, pour l'auteur, ... \fg) et prépare la transition vers la Q2 (\og reste à examiner par quels arguments cette thèse est établie \fg).

\begin{popfigure}
\begin{center}
\begin{tikzpicture}[node distance=0.55cm,
box/.style={draw=popblue, fill=popblueL, rounded corners, text width=9.5cm, align=center, font=\small, thick},
part/.style={draw=poporange, fill=poporangeL, rounded corners, text width=9.5cm, align=center, font=\small, thick},
art/.style={draw=popgreen, fill=popgreenL, rounded corners, text width=9.5cm, align=center, font=\small, thick},
fleche/.style={-{Stealth[length=3mm]}, thick, popdark}]
\node[box] (intro) {\textbf{Introduction} : accroche contextuelle + thèse citée + annonce du plan};
\node[part, below=of intro] (pA) {\textbf{Partie A} : Concept 1 défini + citation expliquée};
\node[part, below=of pA] (pB) {\textbf{Partie B} : Concept 2 défini + citation expliquée};
\node[art, below=of pB] (artic) {\textbf{Articulation} : hiérarchie, opposition ou complémentarité des concepts};
\node[box, below=of artic] (concl) {\textbf{Conclusion} : synthèse du sens global + transition vers la Q2};
\draw[fleche] (intro) -- (pA);
\draw[fleche] (pA) -- (pB);
\draw[fleche] (pB) -- (artic);
\draw[fleche] (artic) -- (concl);
\end{tikzpicture}
\end{center}
\legende{Schéma heuristique de la structure d'une Q1 réussie : chaque étage repose sur le précédent, et l'articulation des concepts fait la charnière entre le développement et la conclusion.}
\end{popfigure}

\courssub{Les interdits formels : ce qui fait perdre des points}

La Q1 est un exercice codifié. Quatre fautes sont formellement interdites :

\begin{itemize}
\item \textbf{La paraphrase} : réécrire le texte avec d'autres mots, phrase après phrase. C'est la faute la plus répandue : elle prouve qu'on a lu, pas qu'on a compris.
\item \textbf{Le commentaire critique} : donner son avis (\og je pense que l'auteur a raison... \fg). Le jugement personnel est réservé à la Q3.
\item \textbf{La dissertation hors texte} : développer tout ce qu'on sait sur le thème (la technique, la liberté...) sans se référer à \emph{ce} passage précis.
\item \textbf{La citation abusive} : recopier de longs extraits sans les expliquer. Une citation non commentée vaut zéro.
\end{itemize}

\begin{attention}[L'erreur la plus sanctionnée]
Faire une dissertation sur le \emph{thème} du texte au lieu d'expliquer \emph{ce texte-là}. Si le texte porte sur la technique chez Arendt et que vous récitez votre leçon générale sur la technique, le correcteur applique le hors-sujet partiel : la note tombe sous 2/5, quelle que soit la qualité de la langue. Règle de survie : chaque paragraphe doit contenir au moins une référence explicite au texte.
\end{attention}

\begin{exemplebox}[Statistique de correction]
Sur 20 copies types d'une classe de Terminale technique, 12 font de la paraphrase en Q1 : elles tournent autour de 1,5/5. Seules les 8 copies qui définissent les concepts, citent brièvement le texte et montrent l'enchaînement logique dépassent 3,5/5. L'écart ne tient pas à l'intelligence, mais à la méthode.
\end{exemplebox}

Le tableau suivant fixe clairement la frontière entre paraphrase et explicitation.

\begin{popfigure}
\begin{center}
\begin{tikzpicture}
\node[draw=popdark, thick, rounded corners, inner sep=6pt, align=center]{
\begin{tabular}{|p{3.6cm}|p{3.6cm}|p{4.6cm}|}
\hline
\textbf{Extrait original} & \textbf{Paraphrase interdite} & \textbf{Explicitation attendue} \\
\hline
\og L'homme moderne fabrique un monde pour ses outils \fg & L'homme d'aujourd'hui construit un environnement adapté à ses instruments. & Inversion de la hiérarchie moyen/fin : la \cle{technicité} n'est plus subordonnée à des fins humaines, c'est elle qui dicte ses fins. Concept : \cle{raison instrumentale} autonomisée. \\
\hline
\og L'Africain doit se penser lui-même \fg & Le Camerounais doit réfléchir par lui-même à sa propre personne. & Revendication de la \cle{subjectivité} : rompre avec l'\cle{aliénation} du regard colonial pour fonder une existence assumée. \\
\hline
\end{tabular}};
\end{tikzpicture}
\end{center}
\legende{Tableau comparatif : la paraphrase déplace les mots sans déplier la pensée ; l'explicitation nomme le concept et dégage l'implication philosophique.}
\end{popfigure}

\courssub{Entraînement guidé sur un extrait du programme}

\begin{exoresolu}[Explication d'un extrait sur la technique]
\textbf{Énoncé.} Texte : \og L'outil servait la main de l'homme ; la machine moderne commande le geste de l'ouvrier. Ce qui était moyen est devenu maître : la technique ne répond plus seulement à nos besoins, elle en crée de nouveaux que nous croyons choisir librement. \fg Expliquez le sens de ce texte.
\tcblower
\textbf{Solution détaillée.}

\emph{Introduction.} Ce texte s'inscrit dans la réflexion philosophique contemporaine sur la technique, illustrée notamment par Arendt. L'auteur y soutient que la technique moderne a renversé son statut : de moyen subordonné à l'homme, elle est devenue puissance qui le détermine. Nous verrons d'abord le constat du renversement, puis l'analyse de l'illusion de liberté qu'il engendre.

\emph{Partie A.} Le texte oppose deux régimes techniques : l'\og outil \fg, qui \og servait la main \fg, et la \og machine \fg, qui \og commande le geste \fg. Le concept central est celui de \cle{raison instrumentale} : la technique comme pur calcul de moyens. Tant que l'outil sert, la fin reste humaine ; quand la machine commande, la hiérarchie moyen/fin s'inverse — c'est le sens de la formule \og ce qui était moyen est devenu maître \fg.

\emph{Partie B.} Le texte enfonce l'analyse : la technique \og crée de nouveaux besoins que nous croyons choisir librement \fg. Apparaît ici le concept d'\cle{aliénation} : le sujet devient étranger à ses propres désirs, qui lui sont dictés de l'extérieur tout en lui laissant le sentiment d'être libre. L'opposition entre \og croyons \fg et la réalité du conditionnement marque l'illusion au cœur de la consommation moderne.

\emph{Conclusion.} Le sens global du texte est donc une critique de l'autonomisation de la technique : elle ne sert plus l'homme mais le configure, jusqu'à fabriquer ses désirs. Reste à examiner, dans la Q2, par quels arguments et quels procédés l'auteur établit cette thèse.
\end{exoresolu}

\begin{exercice}[À vous de jouer]
Texte : \og Une communauté n'est pas une somme d'individus qui se tolèrent ; elle est ce par quoi chacun découvre qu'il existe. \fg Rédigez la Q1 complète (introduction, deux parties, conclusion) en appliquant la grille \og Mots-clés / Concepts / Liens logiques \fg au brouillon avant de rédiger.
\end{exercice}

\begin{corrige}[Exercice]
\emph{Grille attendue au brouillon} : mots-clés \og somme d'individus \fg / \og tolèrent \fg / \og découvre qu'il existe \fg ; concepts : \cle{individu} (sujet isolé), \cle{communauté} (être-ensemble constitutif), \cle{existence} (être reconnu et assumé) ; liens : opposition (\og n'est pas... elle est \fg), causalité (\og ce par quoi \fg).

\emph{Introduction} : le texte traite du rapport individu/communauté ; thèse : la communauté n'est pas un agrégat extérieur aux individus mais la condition de leur existence même ; annonce du plan.

\emph{Partie A} : réfutation de la conception sommative — la tolérance mutuelle ne fait qu'juxtaposer des individus fermés ; concept d'individualisme critiqué.

\emph{Partie B} : thèse positive — la communauté est constitutive : c'est dans le lien aux autres que le sujet accède à sa propre existence (complémentarité individu/communauté).

\emph{Conclusion} : synthèse (l'être humain est fondamentalement relationnel) + transition vers l'analyse argumentative de la Q2.
\end{corrige}

\begin{savaistu}[Gagner du temps en Q1, c'en gagner en Q2]
Une bonne Q1 contient \emph{implicitement} le plan de la Q2 : en dégageant la thèse et les concepts, vous avez déjà repéré les arguments que la Q2 vous demandera d'analyser. Les correcteurs le constatent : les copies fortes en Q1 le sont presque toujours en Q2. Investir sérieusement 45 minutes en Q1, c'est sécuriser la moitié de l'épreuve.
\end{savaistu}

\begin{aretenir}[L'essentiel de la section]
\begin{itemize}
\item La Q1 vérifie un \textbf{savoir} : comprendre la thèse, définir les concepts, montrer l'enchaînement logique.
\item Toujours distinguer \cle{sens littéral} (ce qui est dit) et \cle{sens philosophique} (ce qui est voulu et impliqué).
\item Structure obligatoire : introduction (accroche + thèse + plan), 2 ou 3 parties suivant le texte, conclusion avec transition vers la Q2.
\item Principe de fidélité textuelle : aucune affirmation sans référence précise au texte.
\item Interdits absolus : paraphrase, avis personnel, dissertation hors texte, citation non expliquée.
\end{itemize}
\end{aretenir}

Ayant ainsi sécurisé la compréhension du texte, nous pouvons passer à l'étape suivante : analyser \emph{comment} l'auteur construit sa démonstration, c'est-à-dire la Question 2, consacrée à l'argumentation.

\coursec{Maîtrise de la Question 2 : L'analyse de l'argumentation (Savoir-faire)}

Après avoir maîtrisé l'explicitation du sens (Question 1), où il s'agissait de \emph{restituer} la pensée de l'auteur en la reformulant avec ses propres mots, nous changeons radicalement de registre. La Question 2 ne vous demande plus \og{}que dit le texte ?\fg{}, mais \og{}comment le texte le dit-il ?\fg{} et \og{}pourquoi cela tient-il debout ?\fg{}. Vous passez du rôle d'\emph{exégète} (celui qui interprète) à celui d'\emph{analyste logique} (celui qui dissèque la machinerie interne du raisonnement). C'est l'épreuve reine du \emph{savoir-faire} philosophique : montrer que vous savez reconnaître une structure argumentative, en évaluer la solidité et en rendre compte avec un vocabulaire technique précis.

\courssub{La consigne type et l'enjeu cognitif : passer du « quoi » au « comment »}

Les formulations officielles sont stéréotypées, il faut les reconnaître instantanément pour ne pas perdre de temps :
\begin{itemize}
    \item \og{}Montrez comment l'auteur argue / justifie sa thèse.\fg{}
    \item \og{}Analysez le raisonnement de l'auteur.\fg{}
    \item \og{}De quels arguments l'auteur s'appuie-t-il pour établir sa thèse ?\fg{}
\end{itemize}

\begin{attention}[Piège mortel : La rechute en Q1]
Ne \textbf{jamais} refaire de l'explication de texte (Q1) dans la Q2. Si vous écrivez \og{}L'auteur explique que la technique n'est pas neutre...\fg{}, vous êtes en Q1. En Q2, vous devez écrire : \og{}L'auteur \textbf{argumente} en montrant que la technique porte en elle une finalité propre (argument ontologique), \textbf{s'appuie} sur l'exemple du barrage hydroélectrique (argument factuel/illustratif), \textbf{anticipe} l'objection de l'instrumentalisme pour \textbf{y répondre} par la distinction moyen/fin (dialectique interne).\fg{} La Q2 est une \textbf{analyse de la forme logique}, pas une paraphrase du fond.
\end{attention}

\courssub{Identification de la structure argumentative : l'architecture du texte}

Tout texte philosophique de bac suit une architecture quasi standardisée. Votre première tâche sur brouillon est de dresser la \emph{carte d'identité} du raisonnement.

\begin{definition}[Vocabulaire technique obligatoire de la Q2]
Pour obtenir les points de \og{}vocabulaire technique\fg{} (1 pt sur 7), vous devez utiliser naturellement ces opérateurs logiques :
\begin{itemize}
    \item \textbf{Thèse centrale} : \cle{L'auteur soutient que...}, \cle{Il défend l'idée que...}, \cle{Sa thèse principale est que...}
    \item \textbf{Argumentation} : \cle{Il argumente en montrant que...}, \cle{Il établit que... en invoquant...}, \cle{Son raisonnement repose sur l'idée que...}
    \item \textbf{Illustration/Exemple} : \cle{Il s'appuie sur l'exemple de...}, \cle{Il illustre son propos par le cas de...}, \cle{Il mobilise la référence à...}
    \item \textbf{Dialectique interne (Objection/Réponse)} : \cle{Il anticipe l'objection selon laquelle... pour y répondre que...}, \cle{Il concède que... mais rétorque que...}, \cle{Il écarte l'idée que... au motif que...}
    \item \textbf{Articulation/Enchaînement} : \cle{Dans un premier temps... dans un second temps...}, \cle{Par ailleurs...}, \cle{En outre...}, \cle{Corrélativement...}, \cle{En définitive...}
\end{itemize}
\end{definition}

Voici la représentation graphique type de cette architecture. Mémorisez cette forme arborescente : elle doit devenir votre \og{}squelette\fg{} mental pour toute Q2.

\begin{popfigure}
\begin{tikzpicture}[
    grow'=down,
    level distance=2.2cm,
    sibling distance=3.5cm,
    edge from parent/.style={very thick, draw=popblue, -{Stealth[length=3mm]}},
    every node/.style={align=center, font=\small, inner sep=4pt},
    thesis/.style={rectangle, rounded corners, fill=popblue, text=white, font=\bfseries\small, minimum width=6cm},
    arg/.style={rectangle, rounded corners, fill=popblueL, draw=popblue, minimum width=5cm},
    ex/.style={ellipse, fill=popgreenL, draw=popgreen, minimum width=4cm},
    obj/.style={diamond, fill=poporangeL, draw=poporange, aspect=2, minimum width=4cm},
    rep/.style={rectangle, rounded corners, fill=poppinkL, draw=poppink, minimum width=4.5cm}
]
\node[thesis, align=center] (thesis){THÈSE CENTRALE\\ \textit{(Ce que l'auteur veut prouver)}}
    child { node[arg, align=center] (arg1){ARGUMENT 1\\ \textit{(Logique / Ontologique)}}
        child { node[ex] (ex1) {EXEMPLE / ILLUSTRATION 1} }
        child { node[obj] (obj1) {OBJECTION ANTICIPÉE 1} }
        child { node[rep] (rep1) {RÉPONSE / RÉFUTATION} }
    }
    child { node[arg, align=center] (arg2){ARGUMENT 2\\ \textit{(Factuel / Historique)}}
        child { node[ex] (ex2) {EXEMPLE / ILLUSTRATION 2} }
    }
    child { node[arg, align=center] (arg3){ARGUMENT 3\\ \textit{(Pragmatique / Éthique)}}
        child { node[ex] (ex3) {EXEMPLE / ILLUSTRATION 3} }
    };
% Annotations
\node[left=0.5cm of thesis, text width=3cm, align=right, font=\itshape\small, text=popdark] {SOMMET\\DU TEXTE};
\node[left=0.5cm of arg1, text width=3cm, align=right, font=\itshape\small, text=popdark] {PILIERS\\MAJEURS};
\node[left=0.5cm of ex1, text width=3cm, align=right, font=\itshape\small, text=popdark] {ANCRAGE\\CONCRET};
\node[left=0.5cm of obj1, text width=3cm, align=right, font=\itshape\small, text=popdark] {DIALECTIQUE\\INTERNE};
\end{tikzpicture}
\legende{Arbre argumentatif type : la thèse domine, les arguments se ramifient, les exemples ancrent, la dialectique interne solidifie.}
\end{popfigure}

\courssub{La distinction cruciale : Argument vs Explication}

C'est la compétence discriminante au Bac. Confondre les deux vaut 0 point sur l'argument concerné.

\begin{definition}[Typologie des arguments (à connaître par cœur)]
\begin{itemize}
    \item \textbf{Argument logique (déductif)} : Raisonnement nécessaire. Si prémisses vraies $\rightarrow$ conclusion vraie. Ex: \og{}Tout homme est mortel, Socrate est un homme, donc Socrate est mortel.\fg{} (Syllogisme).
    \item \textbf{Argument logique (inductif)} : Généralisation à partir de cas particuliers. \og{}Tous les cygnes observés sont blancs $\rightarrow$ Tous les cygnes sont blancs.\fg{} (Probable, non certain).
    \item \textbf{Argument factuel (ou empirique)} : S'appuie sur un fait observable, une donnée scientifique, un événement historique. \og{}Le réchauffement climatique est réel car la température moyenne a augmenté de 1,1°C depuis 1850.\fg{}
    \item \textbf{Argument d'autorité} : Renvoie à un expert, un texte fondateur, une tradition. \og{}Comme l'écrit Kant dans \emph{Fondements de la métaphysique des mœurs}...\fg{} \textbf{Attention} : En philosophie, l'autorité ne \emph{prouve} jamais, elle \emph{appuie}. La force vient de la pertinence de la citation, pas du nom.
    \item \textbf{Argument analogique} : Transfère une propriété d'un domaine connu vers un domaine inconnu par similarité de structure. \og{}L'État est comme un navire, le capitaine doit savoir gouverner (Platon).\fg{} Force : heuristique. Faiblesse : la dissimilarité cachée.
    \item \textbf{Argument pragmatique (ou consequentialiste)} : Justifie une thèse par ses effets utiles ou néfastes. \og{}La liberté d'expression doit être protégée car elle permet le progrès social.\fg{}
\end{itemize}
\end{definition}

\begin{propriete}[Règle d'or : Argument vs Explication]
\begin{itemize}
    \item \textbf{L'Argument} : Vise à \textbf{prouver} la thèse. Il répond à \og{}Pourquoi cette thèse est-elle vraie ?\fg{}. Il a une force probante. \emph{Marqueur} : \og{}Donc\fg{}, \og{}Car\fg{}, \og{}Puisque\fg{}, \og{}C'est pourquoi\fg{}.
    \item \textbf{L'Explication} : Vise à \textbf{éclairer} un concept, à en déployer le sens. Elle répond à \og{}Que veut dire ce concept ?\fg{}. Elle n'a \textbf{aucune force probante} pour la thèse. \emph{Marqueur} : \og{}C'est-à-dire\fg{}, \og{}Autrement dit\fg{}, \og{}En d'autres termes\fg{}, \og{}Il faut entendre par là\fg{}.
\end{itemize}
\textbf{Test pratique} : Si vous supprimez le passage, la thèse tombe-t-elle ? $\rightarrow$ Argument. La thèse tient-elle encore mais le concept reste obscur ? $\rightarrow$ Explication.
\end{propriete}

\begin{exemplebox}[Discrimination sur un extrait fictif]
\textit{« La technique n'est pas un simple outil neutre. \textbf{Car} elle impose sa propre logique d'efficacité à la société (Argument logique/ontologique). \textbf{Par exemple}, l'algorithme de notation sociale en Chine structure les comportements citoyens (Argument factuel/illustratif). \textbf{Autrement dit}, la technique devient un environnement total (Explication du concept d'environnement technique). \textbf{On objectera} que l'homme reste maître de ses fins (Objection). \textbf{Mais} cette maîtrise est illusoire car les moyens déterminent les fins (Réponse/Argument). »}

\begin{center}
\begin{tabularx}{\linewidth}{|l|X|X|}\hline
\textbf{Segment} & \textbf{Type} & \textbf{Justification} \\ \hline
\og{}Car elle impose sa logique...\fg{} & \textbf{Argument (Ontologique)} & Prouve la non-neutralité par une nécessité structurelle. \\ \hline
\og{}Par exemple, l'algorithme...\fg{} & \textbf{Argument (Factuel/Illustratif)} & Ancre la thèse dans le réel observable. \\ \hline
\og{}Autrement dit, la technique devient...\fg{} & \textbf{Explication} & Déploie le sens de \og{}environnement total\fg{}, ne prouve rien. \\ \hline
\og{}On objectera que...\fg{} / \og{}Mais...\fg{} & \textbf{Dialectique (Objection/Réponse)} & Renforce la thèse par réfutation adverse. \\ \hline
\end{tabularx}
\end{center}
\end{exemplebox}

\courssub{Traitement des concepts opératoires : l'auteur comme stratège conceptuel}

En Q2, vous devez montrer comment l'auteur \emph{manipule} ses concepts pour les faire fonctionner comme des pièces de son raisonnement. Un concept n'est pas une étiquette fixe ; l'auteur le \textbf{redéfinit}, le \textbf{restreint}, l'\textbf{étend} ou le \textbf{couple} à un autre.

\begin{exemplebox}[Redéfinition stratégique : \og{}Technique\fg{} chez Heidegger vs Ellul]
\begin{itemize}
    \item \textbf{Heidegger (\emph{La question de la technique})} : Redéfinit la technique non pas comme \og{}moyen\fg{} ni \og{}activité humaine\fg{}, mais comme \cle{mode de révélation de la vérité} (\emph{Gestell} / \og{}Engeancement\fg{}). \emph{Enjeu Q2} : L'argument ontologique repose sur ce glissement : si technique = révélation, alors elle échappe à la maîtrise humaine.
    \item \textbf{Ellul (\emph{La technique ou l'enjeu du siècle})} : Définit la technique comme \cle{l'ensemble des méthodes rationnellement arrivées à un maximum d'efficacité dans tout domaine d'activité humaine} (\emph{La Technique} au singulier, système autonome). \emph{Enjeu Q2} : L'argument systémique (autonomie, auto-augmentation) ne tient que si on accepte cette définition totalisante.
\end{itemize}
\textbf{Consigne Q2} : Écrivez \og{}L'auteur \textbf{redéfinit} le concept de technique non plus comme... mais comme..., ce qui lui permet d'argumenter que...\fg{}
\end{exemplebox}

\courssub{Méthodologie sur brouillon : le Tableau d'analyse argumentative}

Ne commencez \textbf{jamais} la rédaction propre sans ce tableau. Il garantit l'exhaustivité sélective (voir Propriété ci-dessous) et la structure de votre développement.

\begin{methode}[Le Tableau d'analyse argumentative (sur brouillon, 10 min max)]
\begin{enumerate}
    \item \textbf{Colonne 1 : N°} (Ordre d'apparition dans le texte).
    \item \textbf{Colonne 2 : Argument (Citation courte + Reformulation)}. Notez les mots-clés du texte (\og{}car\fg{}, \og{}donc\fg{}, \og{}or\fg{}).
    \item \textbf{Colonne 3 : Type} (Logique déductif / Inductif / Factuel / Analogique / Autorité / Pragmatique / Dialectique).
    \item \textbf{Colonne 4 : Force intrinsèque} (Solide / Probable / Faible / Fallacieux). \emph{Évaluez la logique interne, pas votre avis personnel}.
    \item \textbf{Colonne 5 : Limite interne / Objection de l'auteur} (L'auteur la signale-t-il ? Y répond-il ? La laisse-t-il pendante ?).
\end{enumerate}
\textbf{Règle de rédaction} : Un paragraphe développé = Une ligne du tableau (un argument majeur). Les exemples s'intègrent \emph{dans} le paragraphe de l'argument qu'ils soutiennent.
\end{methode}

\begin{propriete}[Règle de l'exhaustivité sélective]
\begin{itemize}
    \item \textbf{Exhaustivité sur le majeur} : Vous \textbf{devez} traiter tous les arguments \emph{structurants} (ceux qui portent la thèse). En rater un = perte de points sèche (souvent 1 pt par argument manquant).
    \item \textbf{Sélectivité sur le mineur} : Ignorez les \emph{connecteurs} simples, les \emph{transitions}, les \emph{explications pures} (sauf si elles redéfinissent un concept clé pour l'argument), les \emph{détails anecdotiques}.
    \item \textbf{Signalement des failles assumées} : Si l'auteur dit \og{}Certes, on pourrait objecter X, mais je ne traiterai pas cela ici\fg{}, vous \textbf{devez} le signaler : \og{}L'auteur écarte délibérément l'objection X, ce qui constitue une limite assumée de sa démonstration.\fg{}
\end{itemize}
\end{propriete}

\courssub{Exercice pratique : Reconstitution du raisonnement sur Kant}

Appliquons la méthode sur l'incipit célèbre de \emph{Qu'est-ce que les Lumières ?} (1784). C'est un classique du Bac technique.

\begin{exoresolu}[Analyse argumentative : Kant, \emph{Qu'est-ce que les Lumières ?} (Extrait)]
\textbf{Énoncé :} \textit{« Les Lumières, c'est la sortie de l'homme hors de sa minorité dont il est lui-même responsable. La minorité, c'est l'incapacité de se servir de son entendement sans la direction d'un autre. [...] Ayez le courage de vous servir de votre propre entendement ! Voilà la devise des Lumières. »}

\textbf{Étape 1 : Identification de la Thèse}
\cle{Thèse} : Les Lumières sont un processus d'émancipation intellectuelle par l'autonomie de la raison (\og{}Sortie de la minorité\fg{}).

\textbf{Étape 2 : Remplissage du Tableau (Brouillon)}

\begin{center}
\begin{tabularx}{\linewidth}{|c|X|c|c|X|}\hline
\textbf{N°} & \textbf{Argument (Citation / Reformulation)} & \textbf{Type} & \textbf{Force} & \textbf{Limite / Dialectique} \\ \hline
1 & \og{}La minorité [...] dont il est lui-même responsable\fg{} (Définition + Imputation responsabilité). & \textbf{Logique (Définitionnal)} / \textbf{Éthique} & \textbf{Forte} : Lie condition (minorité) et cause (responsabilité). Fonde l'impératif moral. & Aucune objection anticipée ici. Pose le postulat anthropologique. \\ \hline
2 & \og{}Incapacité de se servir de son entendement sans la direction d'un autre\fg{} (Analyse de la cause : paresse/lâcheté). & \textbf{Psychologique / Factuel (Observation anthropologique)} & \textbf{Moyenne} : Généralisation sur \og{}l'homme\fg{}. Nécessite l'exemple des \og{}tuteurs\fg{} (v. suite texte). & Kant anticipe : \og{}Il est si commode d'être mineur...\fg{} (Argument pragmatique : confort). \\ \hline
3 & \og{}Ayez le courage de vous servir de votre propre entendement\fg{} (Impératif / Devise). & \textbf{Pragmatique / Normatif (Appel à l'action)} & \textbf{Forte (rhétoriquement)} : Passe du descriptif au prescriptif. \og{}Sapere aude\fg{}. & Exige une condition politique (liberté publique) traitée dans la suite. \\ \hline
\end{tabularx}
\end{center}

\textbf{Étape 3 : Rédaction type Q2 (Niveau Bac - 7/7)}

\og{}Dans cet extrait, \textbf{Kant soutient que les Lumières consistent en une sortie de la minorité}, entendue comme l'incapacité à user de sa raison propre, minorité dont l'homme est \textbf{responsable} par paresse et lâcheté.

\textbf{Il argumente d'abord par une analyse définitionnelle et éthique} : il \textbf{définit} la minorité non comme un déficit intellectuel naturel, mais comme un \textbf{défaut moral} (\og{}dont il est lui-même responsable\fg{}). Ce glissement conceptuel (minorité = faute) \textbf{fonde la possibilité même de la sortie} : ce qui est choisi peut être défait.

\textbf{Dans un second temps, il mobilise un argument psychologique et factuel} : il \textbf{identifie} les causes subjectives de cette minorité, la \og{}paresse\fg{} et la \og{}lâcheté\fg{}, qui expliquent pourquoi les hommes acceptent volontiers la \og{}direction d'un autre\fg{} (tuteurs, pasteurs, médecins). \textbf{Il anticipe l'objection du confort} (\og{}Il est si commode d'être mineur\fg{}) pour mieux la retourner : ce confort est une aliénation consentie.

\textbf{Enfin, il conclut par un argument pragmatique normatif}, lançant l'impératif \og{}Sapere aude\fg{} (\og{}Aie le courage de savoir\fg{}), qui transforme l'analyse descriptive en \textbf{devise programmatique}. \textbf{La structure du raisonnement est donc déductive} : de la définition de l'état (minorité) à la cause (responsabilité), pour aboutir à l'exigence normative (courage/autonomie).\fg{}

\tcblower
\textbf{Corrigé commenté (Grille 7 pts) :}
\begin{itemize}
    \item \textbf{2 pts Thèse} : Claire, précise, cite le texte (\og{}sortie de la minorité\fg{}, \og{}responsable\fg{}).
    \item \textbf{3 pts Arguments} : 3 arguments majeurs traités (Définitionnel/Éthique ; Psychologique/Factuel ; Pragmatique/Normatif) avec citations intégrées et types nommés.
    \item \textbf{1 pt Vocabulaire} : \og{}argumente\fg{}, \textbf{définit}, \textbf{identifie}, \textbf{anticipe l'objection}, \textbf{structure déductive}.
    \item \textbf{1 pt Structure/Liens} : \og{}Dans un second temps\fg{}, \textbf{Enfin}, \og{}pour aboutir à\fg{}. Enchaînement logique visible.
\end{itemize}
\end{exoresolu}

\courssub{Processus d'évaluation de la force argumentative}

Au-delà de l'identification, le candidat d'excellence évalue \emph{interno} la solidité des arguments. C'est ce qui fait la différence entre 5/7 et 7/7.

\begin{popfigure}
\begin{tikzpicture}[
    node distance=1.8cm and 2.5cm,
    box/.style={rectangle, rounded corners, draw=popblue, thick, fill=popblueL, minimum height=1.2cm, text width=3.2cm, align=center, font=\small},
    decision/.style={diamond, draw=poporange, thick, fill=poporangeL, aspect=2, minimum width=3.5cm, text width=3cm, align=center, font=\small},
    eval/.style={rectangle, rounded corners, draw=popgreen, thick, fill=popgreenL, minimum height=1.2cm, text width=3.5cm, align=center, font=\small}
]
\node[box, align=center] (rep){\textbf{1. REPÉRER}\\\og{}Car\fg{}, \og{}Donc\fg{}, \og{}Or\fg{}, \og{}Par exemple\fg{}, \og{}Certes... Mais\fg{}};
\node[decision, right=3cm of rep, align=center] (type){\textbf{2. IDENTIFIER LE TYPE}\\\textit{Logique (Déd/Ind)}\\ \textit{Factuel}\\ \textit{Analogique}\\ \textit{Autorité}\\ \textit{Pragmatique}\\ \textit{Dialectique}};
\node[eval, right=3cm of type, align=center] (eval){\textbf{3. ÉVALUER LA FORCE}\\\textit{Validité logique ?}\\ \textit{Vérité prémisses ?}\\ \textit{Pertinence exemple ?}\\ \textit{Réfutation objection ?}\\ \textit{Redéfinition concept ?}};
\node[box, below=1.5cm of type, xshift=-3cm, align=center] (redac){\textbf{4. RÉDIGER EN Q2}\\\textit{Vocabulaire technique}\\ \textit{Structure par arguments}\\ \textit{Articulation fluide}\\ \textit{Exhaustivité sélective}};
\draw[very thick, -{Stealth[length=3mm]}, poppurple] (eval) -- (redac);
\draw[very thick, -{Stealth[length=3mm]}, poppurple] (type) |- (redac);
\end{tikzpicture}
\legende{Diagramme de flux : le processus mental de l'analyste en Q2. La flèche violette finalise la rédaction.}
\end{popfigure}

\begin{aretenir}[Barème indicatif Q2 (sur 7 points) — Ciblez le 7/7]
\begin{itemize}
    \item \textbf{2 points} : Thèse clairement identifiée, formulée avec les mots du texte + reformulation précise.
    \item \textbf{3 points} : Analyse des arguments majeurs (généralement 3). \textbf{1 pt par argument} : Citation + Type nommé + Rôle dans la démonstration + Articulation avec le suivant.
    \item \textbf{1 point} : Vocabulaire technique maîtrisé (\og{}argumente\fg{}, \og{}s'appuie\fg{}, \og{}anticipe\fg{}, \og{}redéfinit\fg{}, \og{}articule\fg{}).
    \item \textbf{1 point} : Structure globale cohérente (Intro thèse $\rightarrow$ Développement arguments $\rightarrow$ Synthèse structurelle), connecteurs logiques fluides, \textbf{zéro explication de texte (Q1)}.
\end{itemize}
\end{aretenir}

\begin{savaistu}[Le secret des 7/7 : Les connecteurs du texte = Votre plan]
Les meilleurs candidats ne \emph{inventent} pas un plan artificiel (I, II, III). Ils \textbf{suivent la trace des connecteurs logiques du texte lui-même}.
\begin{itemize}
    \item Repérez \og{}Premièrement... Deuxièmement...\fg{} $\rightarrow$ Votre plan : \og{}Dans un premier temps... Dans un second temps...\fg{}
    \item Repérez \og{}Certes... Cependant...\fg{} $\rightarrow$ Votre paragraphe : \og{}L'auteur concède... mais il rétorque...\fg{}
    \item Repérez \og{}C'est pourquoi...\fg{} $\rightarrow$ Votre transition : \og{}Cette argumentation conduit l'auteur à conclure que...\fg{}
\end{itemize}
\textbf{Astuce} : Surlignez \textbf{tous} les connecteurs logiques à la première lecture. Ils sont le \emph{squelette} ; votre rédaction n'est que la \emph{chair} mise dessus.
\end{savaistu}

\begin{attention}[Erreurs fréquentes sanctionnées (Check-list de relecture)]
\begin{enumerate}
    \item \textbf{Confusion Q1/Q2} : Phrases commençant par \og{}L'auteur explique que...\fg{} au lieu de \og{}L'auteur argue que...\fg{}.
    \item \textbf{Arguments fantômes} : Inventer des arguments qui ne sont pas dans le texte (ex: \og{}Il cite Platon\fg{} alors que non).
    \item \textbf{Liste à puces} : \og{}Argument 1 : ... Argument 2 : ...\fg{} sans phrases de liaison, sans articulation. \textbf{Zéro point structure}.
    \item \textbf{Jugement personnel} : \og{}Cet argument est faible car je ne suis pas d'accord.\fg{} $\rightarrow$ Hors sujet. On évalue la \textbf{cohérence interne}, pas la vérité absolue.
    \item \textbf{Oubli de la dialectique} : Ne pas traiter l'objection/réponse si elle est explicite dans le texte = perte de 1 pt argument.
    \item \textbf{Concepts non analysés} : Traiter \og{}liberté\fg{} comme un mot transparent sans voir que l'auteur en donne une définition spécifique (ex: liberté = autonomie vs liberté = arbitraire).
\end{enumerate}
\end{attention}

\courssub{Entraînement dirigé : Application sur un extrait court}

Pour ancrer la méthode, traitez cet extrait de \textbf{Rousseau, \emph{Discours sur l'origine de l'inégalité} (Partie II)} en 20 min chrono (Brouillon + Rédaction).

\begin{exercice}[Entraînement Q2 — Rousseau]
\textit{« Le premier qui, ayant enclos un terrain, s'avisa de dire \og{}Ceci est à moi\fg{} et trouva des gens assez simples pour le croire, fut le vrai fondateur de la société civile. Combien de crimes, de guerres, de meurtres, combien de misères et d'horreurs n'eût point épargnés au genre humain celui qui, arrachant les pieux ou comblant le fossé, eut crié à ses semblables : \og{}Gardez-vous d'écouter cet imposteur ; vous êtes perdus si vous oubliez que les fruits sont à tous et que la terre n'est à personne !\fg{} »}

\textbf{Consigne} : \og{}Montrez comment Rousseau argue pour dénoncer l'origine de la propriété privée et de la société civile.\fg{}

\textbf{Indications pour le brouillon} :
\begin{enumerate}
    \item Thèse ? (Négation de la légitimité naturelle).
    \item Argument 1 : Type ? (Factuel/Historique fictionnel \og{}Le premier qui...\fg{}). Force ?
    \item Argument 2 : Type ? (Argument par l'absurde / Contrefactuel \og{}Combien de crimes...\fg{}). Force ?
    \item Argument 3 : Type ? (Appel à la nature / Droit naturel \og{}les fruits sont à tous\fg{}).
    \item Dialectique ? (L'imposteur / les \og{}gens assez simples\fg{}).
    \item Vocabulaire technique à placer : \og{}fiction fondatrice\fg{}, \og{}argument contrefactuel\fg{}, \og{}dénonce\fg{}, \og{}légitime\fg{}.
\end{enumerate}
\end{exercice}

\begin{corrige}[Exercice n°3 — Rousseau (Proposition de rédaction 7/7)]
\og{}Dans cet extrait, \textbf{Rousseau soutient que la propriété privée et la société civile ne sont pas des institutions naturelles ou légitimes, mais le fruit d'une usurpation historique fondée sur la ruse et acceptée par naïveté}.

\textbf{Il argumente d'abord par une fiction historique fondatrice} : il \textbf{met en scène} \og{}le premier qui, ayant enclos un terrain...\fg{}. Cet \textbf{argument factuel contrefactuel} (reconstitution hypothétique de l'origine) vise à \textbf{dénaturaliser} la propriété : elle a un début, donc elle n'est pas éternelle. La force réside dans la simplicité du geste (\og{}enclos\fg{}, \og{}dit\fg{}, \og{}croit\fg{}) qui dévoile la violence symbolique originelle.

\textbf{Dans un second temps, il déploie un argument par l'absurde (ou contrefactuel) de portée éthique et politique} : il \textbf{évalue les conséquences} de cet acte initial par une énumération gradatoire (\og{}crimes, guerres, meurtres, misères, horreurs\fg{}). \textbf{Il anticipe l'objection de l'ordre social} (la propriété comme fondement de la paix) pour la retourner : c'est la propriété qui \textbf{engendre} le désordre. La force est rhétorique et morale : le bilan catastrophique discrédite l'origine.

\textbf{Enfin, il articule un argument de droit naturel} par la voix du \og{}cri\fg{} du sage : \og{}Les fruits sont à tous et la terre n'est à personne\fg{}. Il \textbf{redéfinit} le rapport à la terre : non pas \og{}appropriation\fg{} (exclusion), mais \og{}communalité\fg{} (usage partagé). \textbf{La structure globale est une généalogie critique} : elle déconstruit la légitimité présente en exposant l'illégitimité de l'origine (ruse + naïveté = imposture).\fg{}

\tcblower
\textbf{Points clés du corrigé} :
\begin{itemize}
    \item Thèse : \og{}Usurpation / Illégitimité\fg{} (2 pts).
    \item Arg 1 : Fiction historique / Dénaturalisation (1 pt).
    \item Arg 2 : Contrefactuel / Bilan catastrophique / Anticipation objection (1 pt).
    \item Arg 3 : Droit naturel / Redéfinition communalité (1 pt).
    \item Vocabulaire : \og{}met en scène\fg{}, \textbf{dénaturalise}, \textbf{évalue}, \textbf{anticipe}, \textbf{redéfinit}, \textbf{généalogie critique} (1 pt).
    \item Structure : \og{}D'abord... Second temps... Enfin... Structure globale\fg{} (1 pt).
\end{itemize}
\end{corrige}

\coursec{Maîtrise de la Question 3 : La discussion critique structurée (Compétence supérieure)}

Après avoir maîtrisé l'explicitation du sens (Q1) et l'analyse de l'argumentation (Q2), vous abordez l'épreuve reine : la discussion critique. Là où la Q1 demande « \emph{Qu'a dit l'auteur ?} » et la Q2 « \emph{Comment l'a-t-il démontré ?} », la Q3 vous somme de prendre la parole : \emph{« Et toi, qu'en penses-tu, arguments à l'appui ? »}. C'est le passage du rôle d'exégète à celui de philosophe. Au Probatoire comme au Baccalauréat technique, cette question vaut souvent 8 points : elle sanctionne votre capacité à penser par vous-même en vous appuyant sur la culture philosophique acquise.

\courssub{La consigne type et l'exigence dialectique}

Les libellés varient — \emph{« Discutez la thèse de l'auteur »}, \emph{« Cette thèse est-elle recevable ? »}, \emph{« Qu'en pensez-vous ? »} —, mais l'attente est unique : une \cle{discussion dialectique} en trois temps, non une simple liste d'opinions. Le plan « Pour / Contre » binaire est un héritage du collège ; en Terminale technique, on exige une \cle{problématisation} qui fait avancer la pensée.

\begin{definition}[Discussion critique philosophique]
Opération rationnelle consistant à soumettre une thèse à l'épreuve de la raison en reconnaissant sa part de vérité (concession), en en dévoilant les limites ou contradictions (critique), pour ouvrir une perspective plus féconde (dépassement/synthèse). Elle ne nie pas la thèse, elle la \cle{complexifie}.
\end{definition}

\begin{propriete}[Structure impérative en 3 temps dialectiques]
Toute copie de Q3 recevable au Bac technique doit articuler visiblement :
\begin{enumerate}
    \item \textbf{Concession (Thèse) :} Reconnaissance argumentée de la pertinence, de la portée ou de la nécessité de la thèse dans son contexte. On montre \emph{pourquoi} l'auteur a raison sur ce point.
    \item \textbf{Critique / Limites (Antithèse) :} Objections fondées : présupposés non examinés, contre-exemples réels, distinctions conceptuelles ignorées, apories logiques, portée limitée. Chaque objection = un argument + une illustration.
    \item \textbf{Dépassement / Synthèse :} Nouvelle distinction, condition de validité, articulation des deux pôles, ouverture sur un problème connexe. On ne « tranche » pas arbitrairement, on \emph{hiérarchise}.
\end{enumerate}
\end{propriete}

\begin{popfigure}
\begin{tikzpicture}[
    node distance=1.8cm and 2.5cm,
    box/.style={draw=popdark, thick, rounded corners, fill=popcream, text width=3.2cm, align=center, font=\small},
    arrow/.style={-Stealth, thick, popblue},
    labelarrow/.style={midway, fill=white, inner sep=2pt, font=\tiny, text=popdark}
]
\node[box, align=center] (thesis){\textbf{THÈSE DE L'AUTEUR}\\(Q1/Q2 : Comprise \& Analysée)};
\node[box, right=of thesis, align=center] (concession){\textbf{1. CONCESSION}\\\emph{« Il a raison car... »}\\\textit{Pertinence, portée, nécessité}};
\node[box, right=of concession, align=center] (critique){\textbf{2. CRITIQUE / LIMITES}\\\emph{« Mais il omet que... »}\\\textit{Présupposés, contre-exemples,\\distinctions, apories}};
\node[box, right=of critique, align=center] (depassement){\textbf{3. DÉPASSEMENT / SYNTHÈSE}\\\emph{« Dès lors, il faut distinguer... »}\\\textit{Nouvelle perspective,\\condition de validité, ouverture}};

\draw[arrow] (thesis) -- node[labelarrow] {Appui Q1/Q2} (concession);
\draw[arrow] (concession) -- node[labelarrow] {Reconnaissance} (critique);
\draw[arrow] (critique) -- node[labelarrow] {Dépassement} (depassement);

% Flèche retour courbe pour montrer la dialectique
\draw[arrow, poporange, dashed, bend left=40] (depassement) to node[labelarrow, above, poporange] {Nouvelle thèse ?} (thesis);
\end{tikzpicture}
\legende{Schéma dialectique de la Q3 : le mouvement de la pensée ne s'arrête pas à la critique, il produit une synthèse qui peut devenir nouvelle thèse.}
\end{popfigure}

\courssub{Le principe de charité : la condition éthique de la critique}

On ne critique pas une caricature. La rigueur philosophique impose de présenter la thèse adverse sous sa forme la plus forte, la plus défendable, avant de l'attaquer. C'est le \cle{principe de charité} : vous avez fait ce travail en Q1 et Q2 ; la Q3 en est le fruit.

\begin{propriete}[Principe de charité (exigible en Q3)]
On ne critique une thèse qu'après l'avoir comprise et présentée sous sa forme la plus forte (acquise en Q1/Q2). Critiquer un homme de paille (\emph{straw man}) — une version affaiblie ou déformée — vaut 0 point à la partie critique.
\end{propriete}

\begin{attention}[Piège : L'ad hominem et la critique de style]
\begin{itemize}
    \item \textbf{Ad hominem :} « Kant était célibataire, donc sa morale est abstraite. » $\rightarrow$ \textbf{Hors sujet.}
    \item \textbf{Style :} « L'auteur écrit difficilement. » $\rightarrow$ \textbf{Irrelevant.}
    \item \textbf{Oubli de la concession :} Attaquer d'emblée. $\rightarrow$ \textbf{Démarche dogmatique, note plafonnée.}
\end{itemize}
La concession prouve à l'examinateur que vous avez \emph{entendu} l'auteur avant de lui répondre.
\end{attention}

\courssub{La « Boîte à outils » critique : votre capital intellectuel}

Une discussion de qualité ne s'improvise pas. Elle puise dans une \cle{boîte à outils} que vous construisez tout au long de l'année. Au programme technique, trois registres sont indispensables et se complètent.

\begin{popfigure}
\begin{tikzpicture}[
    mindmap,
    concept color=popblue,
    level 1 concept/.append style={level distance=4.5cm, sibling angle=90, font=\small\bfseries},
    level 2 concept/.append style={level distance=3cm, sibling angle=45, font=\scriptsize},
    level 3 concept/.append style={level distance=2.5cm, font=\tiny},
    every node/.style={scale=0.85}
]
\node[concept, align=center]{Ma Boîte\\ à Outils Q3}
    child[concept color=popgreen] {node[concept, align=center]{Distinctions\\Conceptuelles}
        child {node[concept] {Légalité / Légitimité}}
        child {node[concept] {Technique / Technologie}}
        child {node[concept] {Individu / Personne}}
        child {node[concept] {Pouvoir / Autorité}}
        child {node[concept] {Égalité / Équité}}
        child {node[concept] {Liberté / Licence}}
    }
    child[concept color=poporange] {node[concept, align=center]{Auteurs\\Référence}
        child {node[concept, align=center]{Programme :\\Platon, Aristote, Descartes,\\Kant, Marx, Bergson,\\Sartre, Ricoeur, Arendt}}
        child {node[concept, align=center]{Hors programme\\pertinents :\\Rawls, Habermas,\\Foucault, Simondon,\\Mbembe, Wiredu}}
    }
    child[concept color=poppink] {node[concept, align=center]{Exemples\\Cameroun (Réalité vécue)}
        child {node[concept, align=center]{Coutumes : Sawa, Bamiléké,\\Beti, Peul (héritage, chefferie)}}
        child {node[concept, align=center]{Droit : OHADA, Code civil,\\Constitution 1996/2008}}
        child {node[concept, align=center]{Numérique : Mobile Money,\\e-gouv, fracture numérique}}
        child {node[concept, align=center]{Environnement : Biodiversité,\\forêt du Bassin du Congo}}
        child {node[concept, align=center]{Vivre-ensemble : bilinguisme,\\décentralisation, ONG}}
    }
    child[concept color=poppurple] {node[concept, align=center]{Logique\\Formelle}
        child {node[concept, align=center]{Syllogisme,\\modus ponens/tollens}}
        child {node[concept, align=center]{Dilemme, réduction\\à l'absurde}}
        child {node[concept, align=center]{Définition par genre\\et différence spécifique}}
    };
\end{tikzpicture}
\legende{Carte mentale : « Ma boîte à outils Q3 ». Les distinctions conceptuelles sont le moteur ; les auteurs donnent l'autorité ; les exemples camerounais ancrent le réel ; la logique assure la rigueur.}
\end{popfigure}

\begin{definition}[Distinction conceptuelle (Opération philosophique majeure)]
Séparer rigoureusement deux notions que le sens commun confond, pour faire apparaître une différence de nature, de portée ou de valeur. C'est le principal levier de la critique (ex: distinguer \emph{Légalité} — conformité à la loi écrite — et \emph{Légitimité} — acceptation morale/fondée par le peuple — pour discuter une thèse sur l'obéissance à l'État).
\end{definition}

\begin{savaistu}[L'exemple camerounais, levier de distinction]
Un exemple camerounais \emph{bien analysé} vaut mieux que trois citations d'auteurs mal comprises.
\begin{itemize}
    \item \textbf{Propriété (Locke) :} Confronter la théorie du travail comme fondement de la propriété à la gestion des terres ancestrales (droit coutumier : la terre appartient aux ancêtres/vivants/morts, inaliénable) vs Code foncier 1974/2001 (immatriculation, titre foncier). \emph{Distinction : Propriété privée libérale vs Propriété communautaire/fonction sociale.}
    \item \textbf{Technique (Heidegger/Simondon) :} Opposer l'importation de machines « clé en main » (technique comme asservissement) à l'innovation locale (ex: \emph{presse à huile artisanale améliorée}, \emph{four à brique écologique} au Nord) — \emph{Technologie vs Savoir-faire endogène}.
    \item \textbf{Liberté vs Coutumes :} Discuter la liberté individuelle (Sartre) face aux rites d'initiation (Ngondo Sawa, Tso'o Bamiléké) ou au mariage forcé. \emph{Distinction : Liberté négative (absence de contrainte) vs Liberté positive (autonomie dans l'appartenance).}
\end{itemize}
\end{savaistu}

\courssub{Méthode : La « Fiche de préparation Q3 » sur brouillon (15 min)}

Ne commencez jamais la rédaction finale sans ce brouillon structuré. C'est la garantie de la cohérence dialectique.

\begin{methode}[La Fiche de préparation Q3 sur brouillon]
\emph{Matériel :} Une feuille A4 pliée en 3 colonnes verticales. \emph{Durée :} 15 minutes max.
\begin{enumerate}
    \item \textbf{Colonne 1 : « Points forts de la thèse » (Pour la Concession).}
    Relevez 2 à 3 arguments centraux de l'auteur (issus de votre Q2). Pour chacun, notez : \textit{« Pourquoi est-il pertinent ? Dans quel contexte ? Quelle expérience le confirme ? »} $\rightarrow$ Cela rédigera votre Partie 1.
    \item \textbf{Colonne 2 : « Objections solides » (Pour la Critique).}
    Pour chaque point fort, cherchez la faille :
    \begin{itemize}
        \item \textit{Distinction conceptuelle :} L'auteur confond-il deux notions ? (ex: Pouvoir/Autorité).
        \item \textit{Présupposé :} Quel postulat non prouvé ? (ex: « L'homme est naturellement bon/mauvais »).
        \item \textit{Contre-exemple réel (Cameroun) :} Une situation concrète qui contredit la thèse.
        \item \textit{Auteur opposé :} Qui a dit le contraire avec force ? (ex: Hobbes vs Rousseau ; Marx vs Smith).
    \end{itemize}
    Sélectionnez les 2 ou 3 objections les plus \emph{philosophiques} (pas des anecdotes). $\rightarrow$ Cela rédigera votre Partie 2.
    \item \textbf{Colonne 3 : « Pistes de dépassement » (Pour la Synthèse).}
    Comment articuler les deux ? 
    \begin{itemize}
        \item \textit{Condition de validité :} « La thèse tient SI... » (ex: La démocratie délibérative fonctionne SI les citoyens sont éduqués/informés).
        \item \textit{Nouvelle distinction :} « Il ne faut pas opposer X et Y, mais distinguer X1 et X2... »
        \item \textit{Ouverture :} « Cela pose le problème de... » (ex: la justice transitionnelle après la critique du droit positif).
    \end{itemize}
    $\rightarrow$ Cela rédigera votre Partie 3 et votre Conclusion.
\end{enumerate}
\end{methode}

\courssub{Soin de la forme : l'architecture visible de la copie}

L'examinateur lit 100 copies en 3h. La vôtre doit être \emph{lisible au premier coup d'œil}.

\begin{aretenir}[Architecture type d'une copie Q3 (8 pts)]
\begin{itemize}
    \item \textbf{Introduction (1 pt) :} \textbf{Amorce} (actualité, citation, fait social camerounais) $\rightarrow$ \textbf{Rappel précis de la thèse} (reprise Q1) $\rightarrow$ \textbf{Problème} (Pourquoi cette thèse est-elle discutable malgré sa force ?) $\rightarrow$ \textbf{Plan annoncé} (ex: « Nous accorderons d'abord... puis nous montrerons les limites... enfin nous proposerons... »).
    \item \textbf{Développement (5 pts) :} 3 paragraphes distincts (saut de ligne + alinéa). \textbf{Transitions fortes} : « Cependant », « Toutefois », « C'est ici que la distinction X/Y devient cruciale », « Au terme de cette critique, il apparaît que... ».
    \item \textbf{Conclusion (1 pt) :} \textbf{Bilan synthétique} (pas de simple liste) $\rightarrow$ \textbf{Réponse nuancée au problème} $\rightarrow$ \textbf{Ouverture} (nouvelle question, enjeu citoyen, perspective technique/professionnelle).
    \item \textbf{Langue (1 pt) :} Zéro faute d'orthographe grave, vocabulaire précis (\emph{distinction, présupposé, aporie, réquisit, hétéronomie/autonomie}), phrases complètes.
\end{itemize}
\end{aretenir}

\begin{attention}[Pièges rédactionnels fréquents]
\begin{itemize}
    \item \textbf{Style familier :} « C'est vrai que... », « Moi je pense que... » $\rightarrow$ Remplacez par : « Il convient de reconnaître que... », « L'analyse montre que... ».
    \item \textbf{Dogmatisme :} « C'est faux. » $\rightarrow$ « Cette thèse atteint ses limites lorsque... ».
    \item \textbf{Accumulation d'exemples non analysés :} « Au Cameroun, il y a le Ngondo, le Ngondo, le Ngondo... » $\rightarrow$ Un seul exemple, \emph{décortiqué} : description $\rightarrow$ analyse conceptuelle $\rightarrow$ portée philosophique.
    \item \textbf{Hors-sujet (Thème vs Thèse) :} Discuter « La liberté » en général au lieu de discuter « La thèse de Sartre sur la liberté radicale dans \emph{L'Être et le Néant} ».
\end{itemize}
\end{attention}

\courssub{Exercice résolu : Simulation Q3 (8 pts)}

\textbf{Sujet :} \textit{Extrait de John Locke, \emph{Second Traité du gouvernement civil} (Chap. V) : « Le travail [...] met une distinction entre [les choses] et le commun ; car ce qui est au-delà de ce que l'usage de la vie demande, est une propriété que l'homme a acquise par son travail. »}
\textit{Consigne :} \textbf{Discutez la thèse de l'auteur sur l'origine de la propriété.}

\begin{exoresolu}[Q3 : Discussion de la thèse lockéenne sur la propriété]
\textbf{Énoncé :} Discutez la thèse selon laquelle le travail est l'origine et la justification de la propriété privée.

\textbf{Solution détaillée :}

\tcblower

\textbf{Introduction}
[Amorce] Au Cameroun, le conflit foncier entre promoteurs immobiliers titulaires de titres fonciers (loi n° 74-1 du 6 juillet 1974) et communautés villageoises gérant la terre selon le droit coutumier (terre des ancêtres, inaliénable) illustre violemment la question de l'origine de la propriété. [Rappel thèse] John Locke, dans le \emph{Second Traité}, fonde la propriété privée sur le travail : en mélangeant son travail aux choses communes, l'homme s'approprie légitimement une portion de la nature. [Problème] Cette thèse, fondatrice du libéralisme, est-elle universellement recevable ou masque-t-elle des présupposés anthropologiques et historiques contestables ? [Plan] Nous accorderons d'abord à Locke la pertinence de lier propriété et responsabilité productive (I), avant de montrer les limites anthropologiques et historiques de ce « mythe originaire » (II), pour enfin distinguer propriété lucrative et propriété d'usage, condition d'une justice foncière (III).

\textbf{Développement}

\textbf{I. Concession : La force heuristique du critère travail.}
Locke a raison de lier propriété et \emph{labor}. Premièrement, le travail transforme la nature en ressource utile : défricher, cultiver, bâtir crée de la valeur là où il n'y avait que potentialité. Cela fonde une \emph{responsabilité} : le propriétaire est celui qui « améliore » le bien (théorie de l'amélioration). Deuxièmement, ce critère offre une limite naturelle : la « clause de non-gaspillage » (n'approprions que ce qu'on peut utiliser avant qu'il ne pourrisse) et la clause « assez et aussi bon » pour les autres. Dans une économie de subsistance pré-monétaire, cette thèse légitime l'effort paysan et protège contre l'accaparement oisif. Au Cameroun, le paysan qui met en valeur une jachère par son travail (cacao, café) acquiert une légitimité morale reconnue par la coutume comme par le droit moderne : le travail \emph{fait} la propriété.

\textbf{II. Critique : Les limites du fondement par le travail.}
Cependant, trois objections majeures fissurent cette thèse.
\textit{1. L'invisibilisation du travail collectif et féminin (Distinction conceptuelle).} Locke pose un \emph{individu} abstrait, propriétaire de sa personne et de sa force de travail. Or, la mise en valeur de la terre est historiquement collective (lignage, groupe d'âge) et genrée (femmes cultivateuses sans droit foncier coutumier). La distinction \emph{Individu / Personne} (Mauss, Ricoeur) et \emph{Travail productif / Travail de care} (Fraser) montre que le « mélange du travail » efface les contributeurs invisibles. Au Cameroun, la femme qui cultive la terre de son mari n'en acquiert pas la propriété coutumière.
\textit{2. Le leurre de l'origine : l'accumulation primitive (Marx).} La thèse suppose un état de nature où la terre est « commune » et libre. L'histoire (Amérique, Australie, Cameroun colonial) montre que la propriété privée naît souvent de la \emph{spoliation} (conquête, décret colonial, accaparement) légitimée \emph{a posteriori} par le mythe du travail. Marx parle d'« accumulation primitive » : le travail ne crée pas le titre, il le justifie rétroactivement.
\textit{3. L'aporie de l'argent et de l'accumulation illimitée.} Locke admet que l'invention de la monnaie (or, argent) permet d'accumuler sans gaspillage (l'argent ne pourrit pas). La « clause assez et aussi bon » devient caduque. La thèse légitime ainsi l'inégalité radicale : le travail de l'ouvrier enrichit le propriétaire du capital. La distinction \emph{Travail / Capital} (Marx, Piketty) révèle que la propriété issue du travail propre se mue en propriété sur le travail d'autrui.

\textbf{III. Dépassement : Vers une distinction opératoire — Propriété d'usage vs Propriété lucrative.}
Le dépassement ne consiste pas à nier la propriété, mais à en \emph{typer} les régimes. Il faut distinguer (distinction conceptuelle clé) :
\begin{itemize}
    \item La \textbf{propriété d'usage / de subsistance} (terre cultivée, logement, outils de travail) : elle relève du droit naturel lockeen, protégée par le travail effectif et la destination universelle des biens (Doctrine sociale de l'Église, Art. 17 DUDH). Au Cameroun, le certificat d'occupation foncière ou le titre foncier d'habitation relèvent de ce régime.
    \item La \textbf{propriété lucrative / spéculative} (foncier nu en attente de plus-value, parts sociales, brevets bloquants) : elle ne relève pas du « mélange du travail » mais de la rente ou du pouvoir financier. Elle appelle une régulation forte (taxe foncière progressive, droit de préemption urbain, plafonnement des surfaces).
\end{itemize}
La thèse de Locke reste valide pour le premier régime ; elle devient idéologique si elle sert à sacraliser le second.

\textbf{Conclusion}
[Bilan] La thèse de Locke révèle une vérité anthropologique : l'appropriation par le travail fonde l'autonomie et la responsabilité. Mais elle occulte la dimension sociale, historique et genrée de la propriété. [Réponse] Elle est recevable comme fondement de la \emph{propriété d'usage}, inopérante comme justification de la \emph{propriété capitalistique}. [Ouverture] Cette distinction impose au juriste technique (notaire, géomètre, urbaniste) de ne pas confondre \emph{titrer} (acte juridique) et \emph{légitimer} (justice foncière). Le défi camerounais actuel (réforme foncière, décentralisation) est d'inventer des formes de propriété hybrides (communautaires, coopératives) qui honorent le travail sans figer l'exclusion.
\end{exoresolu}

\begin{exemplebox}[Barème indicatif Q3 (sur 8 pts) — Auto-évaluation]
\begin{center}
\begin{tabularx}{\linewidth}{|l|X|c|}\hline
\textbf{Critère} & \textbf{Indicateurs de réussite} & \textbf{Pts} \\ \hline
Introduction / Plan & Amorce contextualisée (Cameroun si possible), thèse reprise, problème explicite, plan dialectique annoncé. & 1 \\ \hline
Concession (Partie 1) & 1 à 2 arguments forts de la thèse reconnus, expliqués, illustrés (ex: clause non-gaspillage). Ton charitable. & 2 \\ \hline
Critique (Partie 2) & 2 à 3 objections \textbf{structurées} (Distinction conceptuelle + Auteur + Exemple CM). Pas de « pour/contre » plat. & 3 \\ \hline
Dépassement (Partie 3) & Nouvelle distinction opératoire, condition de validité, ou articulation synthétique. Ouverture pertinente. & 1 \\ \hline
Conclusion / Langue & Bilan nuancé, réponse au problème, ouverture. Français correct, vocabulaire précis, 0 faute grave. & 1 \\ \hline
\end{tabularx}
\end{center}
\legende{Grille d'auto-évaluation à garder en tête pendant l'épreuve. Visez les 3 pts de la critique : c'est là que se joue la note.}
\end{exemplebox}

\courssub{Transfert immédiat : Entraînement guidé}

Appliquez la méthode sur ce sujet type Bac Technique :

\begin{exercice}[Entraînement Q3 — 30 min chrono]
\textbf{Sujet :} \textit{« La technique n'est pas un outil neutre ; elle porte en elle une rationalité spécifique qui structure nos vies. » (D'après Heidegger, \emph{Essai sur la technique}).}
\textbf{Consigne :} Discutez cette thèse.

\textbf{Consignes de travail :}
\begin{enumerate}
    \item Remplissez votre \textbf{Fiche de préparation Q3} (3 colonnes) en 10 min.
    \item Rédigez l'\textbf{Introduction complète} (amorce CM, thèse, problème, plan).
    \item Rédigez la \textbf{Transition I $\rightarrow$ II} et la \textbf{Transition II $\rightarrow$ III}.
    \item Écrivez la \textbf{Conclusion avec ouverture}.
\end{enumerate}
\textit{Indices boîte à outils :} Distinction \emph{Technique / Technologie} (Ellul/Simondon) ; \emph{Moyen / Fin} ; Exemple CM : Mobile Money (Orange Money/MTN MoMo) comme infrastructure de pouvoir vs usage citoyen ; Auteurs : Heidegger, Ellul, Simondon, Feenberg, \emph{législation OHADA et nationale sur les transactions électroniques}.
\end{exercice}

\begin{corrige}[Exercice Q3 — Éléments attendus]
\textbf{Fiche brouillon (extrait) :}
\begin{itemize}
    \item \textbf{Col 1 (Concession) :} Technique = moyen efficace (moyen/fin) ; libère de la peine (ex: pompe à eau vs puisage) ; rationalité calculante nécessaire pour ingénieur/technicien.
    \item \textbf{Col 2 (Critique) :} 1. \emph{Distinction Technique/Technologie (Ellul)} : La technique devient \emph{système} (technicien) autonomisant l'efficacité comme valeur suprême (ex: algorithmes de notation sociale, IA prédictive). 2. \emph{Non-neutralité axiologique (Heidegger)} : « L'essence de la technique n'est rien de technique » (Gestell) — elle nous \emph{en-cadre} (ex: Mobile Money formate l'épargne, la dette, l'identité numérique). 3. \emph{Contre-exemple CM :} Savoir-faire local (four à brique, pharmacopée) = technique \emph{non} asservissante, intégrée au cosmos social.
    \item \textbf{Col 3 (Dépassement) :} Distinction \emph{Technique asservissante / Technique libératrice} (Feenberg) ou \emph{Instrumentalisation / Appropriation}. Condition : \emph{Gouvernance technique} (participation usagers, éthique by design, droit à la déconnexion, logiciels libres). Ouverture : Responsabilité de l'ingénieur technicien (déontologie).
\end{itemize}
\textbf{Introduction type :} « Lorsque le paysan du Nord-Cameroun utilise son téléphone pour vérifier le cours du coton via Mobile Money avant de vendre, il croit maîtriser un outil. Pourtant, l'algorithme qui fixe les frais de transaction, le taux de change appliqué, l'historique de crédit qui conditionnera son prochain prêt, tout cela le précède et le structure. Heidegger affirme que la technique n'est pas neutre... »
\textbf{Transition I$\rightarrow$II :} « Certes, la technique comme moyen libère l'homme de la nécessité naturelle. \textbf{Toutefois}, cette efficacité même devient une fin en soi lorsqu'elle s'autonomise en \emph{système technicien} (Ellul). »
\textbf{Transition II$\rightarrow$III :} « \textbf{Au terme de cette critique}, il ne s'agit pas de refuser la technique (impossible et indésirable), mais de distinguer les régimes de technicité... »
\textbf{Conclusion :} « La thèse de Heidegger est recevable comme critique de l'aveuglement instrumental. Elle nous oblige, futurs techniciens supérieurs, à concevoir des techniques \emph{appropriables}, soumises au débat démocratique, non subies. La neutralité n'est pas dans l'objet, elle est dans la gouvernance. »
\end{corrige}

\begin{aretenir}[Check-list finale Q3 — À relire la veille de l'épreuve]
\begin{itemize}
    \item [ ] Ai-je fait ma \textbf{Fiche 3 colonnes} sur brouillon ?
    \item [ ] Mon \textbf{introduction} a-t-elle une amorce \textbf{camerounaise} ou d'actualité technique ?
    \item [ ] Ma \textbf{concession} cite-t-elle explicitement des arguments de la Q2 ?
    \item [ ] Mes \textbf{objections} utilisent-elles des \textbf{distinctions conceptuelles} (pas juste des avis) ?
    \item [ ] Ai-je au moins \textbf{un exemple camerounais analysé} (pas cité) ?
    \item [ ] Mon \textbf{dépassement} propose-t-il une \textbf{nouvelle distinction} ou une \textbf{condition de validité} ?
    \item [ ] Mes \textbf{transitions} font-elles avancer la dialectique (mots-charnières) ?
    \item [ ] Ma \textbf{conclusion ouvre-t-elle} (pas de simple résumé) ?
    \item [ ] Relu : \textbf{Zéro faute grave}, vocabulaire précis (\emph{présupposé, aporie, réquisit, hétéronomie}).
\end{itemize}
\end{aretenir}

\coursec{Atelier de consolidation : Simulation Bac \& Auto-évaluation croisée}

Dans la section précédente, nous avons appris à construire une discussion critique structurée (Q3) : thèse de l'auteur, argumentation, objection fondée, réponse nuancée. Mais savoir faire chaque question séparément ne suffit pas : au Baccalauréat technique, tout se joue dans \emph{une seule copie}, en \emph{un temps limité}, sous \emph{pression}. Cette section est donc un atelier grandeur nature : nous allons simuler l'épreuve, corriger comme le jury, échanger nos copies, et fabriquer l'outil de révision ultime. Objectif : transformer le savoir méthodologique en \cle{réflexe d'examen}.

\courssub{Mise en situation d'examen blanc : le sujet inédit}

\begin{definition}[La simulation Bac]
Une \cle{simulation Bac} (ou examen blanc) est une épreuve réalisée dans les conditions réelles : sujet inédit, temps limité chronométré, copie individuelle, sans document, sans téléphone. Elle ne vise pas seulement à « vérifier les connaissances », mais à entraîner la \emph{gestion du temps}, la \emph{gestion du stress} et l'\emph{enchaînement des trois questions}.
\end{definition}

\textbf{Protocole de la séance.} Deux formats sont possibles selon l'emploi du temps :
\begin{itemize}
\item \textbf{Format court (1 h 30)} : traitement des questions 1 et 2 uniquement (explicitation du sens + analyse de l'argumentation). C'est le minimum vital.
\item \textbf{Format complet (3 h)} : les trois questions, dans les conditions exactes du Baccalauréat technique. C'est l'idéal en fin de séquence.
\end{itemize}

\begin{exemplebox}[Sujet inédit type Bac Tech — extrait de Hans Jonas, \emph{Le Principe responsabilité}]
«~La nature humaine, modifiée par la technique, est devenue un objet de la responsabilité.\dots Le type nouveau d'action humaine exige une éthique nouvelle, avec des règles nouvelles, une éthique de la responsabilité à long terme, coextensive à l'étendue de notre pouvoir.\dots Agis de façon que les effets de ton action soient compatibles avec la permanence d'une vie authentiquement humaine sur terre.~»
\par\medskip
\textbf{Question 1} : Dégagez l'idée principale du texte et expliquez les notions de \emph{responsabilité} et de \emph{pouvoir}. \textbf{Question 2} : Analysez l'argumentation de l'auteur (thèse, arguments, enjeu). \textbf{Question 3} : La technique moderne oblige-t-elle à inventer une morale nouvelle ?
\end{exemplebox}

\begin{attention}[Erreur fréquente : sacrifier la Q3]
Beaucoup de candidats, mal organisés, rédigent longuement Q1 et Q2, puis n'ont plus que dix minutes pour Q3 — question pourtant la plus valorisée, car elle mesure la compétence supérieure de discussion personnelle. \textbf{Solution de secours} : si le temps presse, écrivez \emph{d'abord au brouillon} l'introduction et le plan détaillé de la Q3 (thèse de l'auteur, deux arguments, une objection, une réponse nuancée, conclusion). Même si la rédaction finale est courte, le correcteur verra une pensée structurée. Une Q3 brève mais architecturée vaut mieux qu'une Q3 absente ou brouillonne.
\end{attention}

\courssub{Correction guidée collective : penser comme le jury}

Après la simulation, le professeur projette \textbf{trois copies anonymisées} : une copie faible, une copie moyenne, une copie excellente. La classe applique alors la grille officielle d'évaluation. L'exercice est capital : \emph{celui qui sait corriger sait rédiger}. En identifiant ce qui distingue une copie excellente d'une copie moyenne, l'élève intériorise les \cle{critères discriminants}.

\begin{definition}[La copie « type jury »]
Une copie \cle{type jury} est une copie qui facilite le travail du correcteur et met en valeur la pensée du candidat. Elle est : \textbf{lisible} (écriture soignée, sans ratures excessives), \textbf{structurée} (paragraphes séparés par des sauts de ligne, alinéas visibles), avec les \textbf{concepts en gras ou soulignés}, les \textbf{citations intégrées dans la phrase} («~Jonas affirme que "la nature humaine est devenue un objet de la responsabilité"~», et non une citation jetée isolément), et \textbf{sans «~je pense que~»} en Q1 et Q2, où l'on doit restituer la pensée de l'auteur, pas la sienne. Le «~je~» n'apparaît qu'en Q3, dans la discussion.
\end{definition}

\begin{exemplebox}[Critères discriminants observés sur les trois copies]
\begin{itemize}
\item \textbf{Copie faible} : paraphrase du texte (l'élève répète les phrases de Jonas avec d'autres mots), aucune notion définie, Q3 absente. Le correcteur ne trouve aucun savoir vérifiable.
\item \textbf{Copie moyenne} : idée principale dégagée, notions définies mais sans lien avec le texte, argumentation repérée mais non analysée, Q3 réduite à une opinion personnelle non argumentée.
\item \textbf{Copie excellente} : idée principale formulée en une phrase précise, notions définies \emph{puis réinvesties} dans l'explication du texte, thèse et arguments explicités avec connecteurs logiques («~d'abord\dots ensuite\dots c'est pourquoi~\dots~»), Q3 avec objection réelle (peut-on moraliser la technique sans freiner le progrès ?) et réponse nuancée.
\end{itemize}
\end{exemplebox}

\begin{popfigure}
\begin{center}
\renewcommand{\arraystretch}{1.3}
\begin{tabularx}{\linewidth}{|l|X|X|X|X|}
\hline
\rowcolor{popblueL} \textbf{Critère} & \textbf{Non acquis} & \textbf{En cours} & \textbf{Acquis} & \textbf{Expert} \\
\hline
\textbf{Fidélité au texte} & Paraphrase ou hors-sujet & Idée principale approximative & Idée principale exacte et citée & Sens restitué avec finesse, citations intégrées \\
\hline
\textbf{Concepts} & Absents ou faux & Cités sans définition & Définis correctement & Définis, distingués, réinvestis dans l'analyse \\
\hline
\textbf{Argumentation (Q2)} & Non repérée & Thèse trouvée, arguments confus & Thèse + arguments + enjeu & Structure logique reconstituée avec connecteurs \\
\hline
\textbf{Critique (Q3)} & Absente & Opinion sans argument & Objection + réponse & Discussion nuancée, position personnelle justifiée \\
\hline
\textbf{Langue} & Nombreuses fautes, illisible & Fautes gênantes & Correct et clair & Style précis, vocabulaire philosophique maîtrisé \\
\hline
\textbf{Temps} & Q3 non traitée & Fin bâclée & Les 3 questions traitées & Répartition équilibrée et relecture faite \\
\hline
\end{tabularx}
\end{center}
\legende{Grille d'auto-évaluation croisée. À utiliser en binôme : chaque élève cote la copie de son camarade ligne par ligne, puis justifie oralement chaque niveau attribué.}
\end{popfigure}

\courssub{Travail en binômes : l'évaluation croisée au code couleur}

\begin{methode}[L'annotation croisée en trois couleurs]
\begin{enumerate}
\item \textbf{Échange} : chaque binôme échange ses copies de simulation. On corrige la copie de l'autre, jamais la sienne (le regard extérieur voit ce que l'auteur ne voit plus).
\item \textbf{Annotation au code couleur} : \textcolor{popgreen}{\textbf{vert}} pour ce qui est réussi (définition exacte, argument bien repéré, belle formulation) ; \textcolor{poporange}{\textbf{orange}} pour ce qui est imprécis (notion citée sans définition, citation non intégrée, connecteur absent) ; \textcolor{popink}{\textbf{rouge}} pour l'erreur ou le hors-sujet (contresens sur l'auteur, «~je pense que~» en Q1, paraphrase).
\item \textbf{Retour argumenté} : chaque annotation rouge ou orange doit être accompagnée d'une justification écrite en marge («~ici tu répètes le texte au lieu de l'expliquer~», «~cette notion n'est pas définie~»). Une annotation sans justification n'aide pas.
\item \textbf{Restitution orale} : le correcteur explique à l'auteur ses trois points forts et ses trois axes de progrès, en deux minutes. L'auteur écoute sans se justifier, puis pose ses questions.
\end{enumerate}
\end{methode}

\begin{attention}[Piège de l'évaluation croisée]
Deux écueils : la \emph{complaisance} (tout mettre en vert pour ne pas blesser son camarade) et la \emph{sévérité gratuite} (tout mettre en rouge sans explication). Règle d'or : annoter comme on voudrait être annoté — avec exigence sur le fond et bienveillance sur la forme. Le but n'est pas de noter l'autre, mais de le faire progresser.
\end{attention}

\courssub{Rédaction collaborative : la Q3 « modèle » au tableau}

L'atelier se poursuit par une construction collective. Le professeur interroge la classe : quelles sont les meilleures idées produites dans vos Q3 ? Les réponses sont notées au tableau, triées, hiérarchisées, puis mises en forme.

\begin{exoresolu}[Construction collective d'une Q3 modèle sur le texte de Jonas]
Énoncé : «~La technique moderne oblige-t-elle à inventer une morale nouvelle ?~» Construisez une discussion structurée à partir des apports de la classe.
\tcblower
\textbf{Solution détaillée (rédaction finale soignée).}
\par\medskip
\emph{Introduction.} Jonas soutient que la puissance inédite de la technique moderne rend notre responsabilité «~coextensive à l'étendue de notre pouvoir~». Faut-il en conclure que les morales traditionnelles sont dépassées et qu'il faut en inventer une entièrement nouvelle ?
\par
\emph{Thèse de l'auteur et arguments.} Pour Jonas, la morale ancienne supposait des actions aux effets proches et réversibles : elle ne visait que le voisin immédiat. Or la technique moderne — industrie, nucléaire, biotechnologies — produit des effets planétaires et irréversibles, qui engagent les générations futures. Dès lors, seule une éthique nouvelle, fondée sur la responsabilité envers l'humanité à venir, peut guider l'action : «~Agis de façon que les effets de ton action soient compatibles avec la permanence d'une vie authentiquement humaine sur terre.~»
\par
\emph{Objection.} On peut objecter que les grandes morales traditionnelles contiennent déjà l'exigence de ne pas nuire autrui : le commandement «~tu ne tueras pas~» ou l'impératif kantien d'universalité suffiraient, étendus aux générations futures. Inventer une morale nouvelle serait inutile ; il suffirait d'élargir l'ancienne.
\par
\emph{Réponse nuancée.} Mais l'objection sous-estime la nouveauté de la situation : les morales anciennes ne prévoyaient pas un pouvoir capable de détruire les conditions mêmes de la vie humaine. La responsabilité jonasienne n'abolit pas les anciens devoirs, elle les \emph{complète} par une dimension inédite : le devoir envers des êtres qui n'existent pas encore et ne peuvent rien réclamer.
\par
\emph{Conclusion.} La technique moderne n'oblige donc pas à jeter la morale ancienne, mais à la prolonger par un principe nouveau : la responsabilité envers l'avenir de l'humanité et de la nature. Au Cameroun, la question de l'exploitation des forêts ou des ressources minières illustre concrètement cette exigence : nos choix techniques d'aujourd'hui engagent la vie des générations de demain.
\end{exoresolu}

\courssub{La relecture « à froid » et la Fiche Mémo Méthodo}

\begin{methode}[La relecture « à froid » en 3 passes]
Réservez impérativement les dix dernières minutes de l'épreuve à la relecture, en trois passes distinctes — on ne peut pas tout vérifier en même temps :
\begin{enumerate}
\item \textbf{Passe 1 — Sens global et respect des consignes} : ai-je répondu aux trois questions ? L'idée principale est-elle formulée ? Y a-t-il du hors-sujet ?
\item \textbf{Passe 2 — Vocabulaire philosophique et connecteurs} : les notions sont-elles définies ? Les connecteurs logiques («~car~», «~en effet~», «~cependant~», «~par conséquent~») marquent-ils le raisonnement ? Les citations sont-elles intégrées ?
\item \textbf{Passe 3 — Orthographe, grammaire, ponctuation} : accords (sujet-verbe, nom-adjectif), majuscules aux concepts personnifiés ou aux noms propres (\emph{l'État}, \emph{la Nature}, Kant, Jonas), ponctuation des citations.
\end{enumerate}
\end{methode}

\begin{savaistu}[Le capital sympathie du correcteur]
Le correcteur évalue aussi la \emph{qualité de l'expression}, qui pèse souvent 2 à 3 points dans la note globale. Une copie propre, sans ratures, bien paragraphe, avec une écriture lisible, commence avec un «~capital sympathie~» : à contenu égal, elle est mieux notée qu'une copie sale et désorganisée. Le correcteur lit des centaines de copies en quelques jours : faciliter sa lecture, c'est déjà plaider pour soi.
\end{savaistu}

\begin{exemplebox}[Un résultat chiffré qui doit vous convaincre]
Le suivi des classes de Terminale technique ayant pratiqué des examens blancs montre un \textbf{gain moyen de 2 à 3 points} entre la première simulation et l'épreuve réelle du Bac, pour les élèves ayant effectué \textbf{trois simulations chronométrées corrigées en classe}. Le gain ne vient pas de connaissances nouvelles, mais de la maîtrise du temps, de la structure et de la relecture. La méthode, c'est des points.
\end{exemplebox}

\begin{popfigure}
\begin{center}
\begin{tikzpicture}[
  every node/.style={font=\small},
  racine/.style={rectangle, rounded corners, draw=popdark, fill=popblueL, thick, align=center, font=\small\bfseries, minimum width=4.2cm},
  branche/.style={rectangle, rounded corners, draw=popblue, fill=popcream, align=center, font=\footnotesize, text width=4.6cm},
  lien/.style={thick, popmuted}
]
\node[racine, align=center] (r) at (0,0){FICHE MÉMO MÉTHODO BAC\\(format A5 recto-verso)};
\node[branche, align=center] (q1) at (-5.2,2.6){\textbf{Q1 — Le sens}\\Idée principale en 1 phrase\\Définir chaque notion\\Expliquer = relier notions + texte\\Jamais de «~je~»};
\node[branche, align=center] (q2) at (5.2,2.6){\textbf{Q2 — L'argumentation}\\Thèse / arguments / enjeu\\Connecteurs logiques\\Citations intégrées dans la phrase\\Rester fidèle à l'auteur};
\node[branche, align=center] (q3) at (-5.2,-2.6){\textbf{Q3 — La discussion}\\Thèse de l'auteur $\rightarrow$ objection\\$\rightarrow$ réponse nuancée\\Position personnelle justifiée\\Conclusion ouverte};
\node[branche, align=center] (org) at (5.2,-2.6){\textbf{Organisation}\\Lecture + plan : 30 min\\Q1 : 45 min \quad Q2 : 45 min\\Q3 : 45 min \quad Relecture : 15 min\\Plan de Q3 au brouillon d'abord !};
\draw[lien] (r) -- (q1);
\draw[lien] (r) -- (q2);
\draw[lien] (r) -- (q3);
\draw[lien] (r) -- (org);
\end{tikzpicture}
\end{center}
\legende{Modèle de «~Fiche Mémo Méthodo Bac~» : une carte mentale compacte sur carte A5 recto-verso. Recto : les trois questions et leurs exigences. Verso : la répartition du temps, les connecteurs utiles et les pièges à éviter. À réviser la veille de l'examen — et rien d'autre.}
\end{popfigure}

\courssub{Gestion du stress et stratégies de secours}

Le jour du Bac, l'imprévu est la règle, pas l'exception. Voici les quatre situations critiques et leurs parades :

\begin{itemize}
\item \textbf{Texte très court.} Ne paniquez pas : un texte court est dense. Chaque phrase compte. Analysez les termes un à un, dégagez les présupposés (ce que l'auteur admet sans le dire), exploitez la moindre opposition de mots. La brièveté du texte autorise une explication plus fine.
\item \textbf{Texte très long.} Ne lisez pas tout au même niveau : repérez d'abord la première phrase (souvent la thèse) et la dernière (souvent la conclusion), puis les articulations («~mais~», «~donc~», «~cependant~»). Sélectionnez trois ou quatre passages clés et construisez l'explication autour d'eux.
\item \textbf{Auteur inconnu.} C'est sans importance : l'exercice porte sur le \emph{texte}, pas sur l'auteur. La méthode est identique — idée principale, notions, arguments. Ne perdez pas une minute à chercher qui est l'auteur ; le nom en bas de page suffit.
\item \textbf{Thème non révisé.} Revenez aux notions fondamentales transversales : la liberté, le devoir, la vérité, la technique, l'État. Presque tout texte philosophique mobilise l'une d'elles. Et rappelez-vous : la méthode (définir, analyser, discuter) vaut pour \emph{tous} les sujets.
\end{itemize}

\begin{attention}[Le stress n'est pas l'ennemi]
Un stress modéré améliore la concentration. L'ennemi, c'est la panique. Technique de secours : si vous bloquez, posez le stylo, respirez profondément trois fois, relisez la première phrase du texte et demandez-vous simplement : «~Que veut dire cette phrase ?~» Écrivez la réponse au brouillon. Le mouvement de la pensée est relancé.
\end{attention}

\begin{propriete}[Focus Bac — Sujets probables 2026-27]
D'après les orientations de la fiche de progression harmonisée MINESEC, quatre thématiques sont à réviser en priorité pour l'exercice sur texte :
\begin{itemize}
\item \textbf{Technique et Environnement} : Hans Jonas (\emph{Le Principe responsabilité}) — la responsabilité envers les générations futures ;
\item \textbf{Intelligence artificielle et Sujet} : Hans Jonas, Günther Anders — l'homme peut-il rester maître de ses créations techniques ?
\item \textbf{Démocratie et réseaux sociaux} : Jürgen Habermas (l'espace public), Byung-Chul Han (la critique de la transparence numérique) ;
\item \textbf{Identité et mondialisation} : Kwame Anthony Appiah, Achille Mbembe — l'identité entre enracinement et ouverture au monde.
\end{itemize}
Pour chaque thème, préparez : deux définitions de notions, un argument d'auteur, une objection classique.
\end{propriete}

\begin{aretenir}[L'essentiel de l'atelier]
\begin{itemize}
\item La simulation chronométrée est le seul entraînement qui prépare vraiment à l'épreuve : trois simulations corrigées = 2 à 3 points gagnés en moyenne.
\item Corriger la copie d'un camarade avec la grille officielle et le code couleur (vert / orange / rouge) apprend à rédiger mieux qu'un cours magistral.
\item Une copie «~type jury~» est lisible, structurée, avec concepts soulignés, citations intégrées, et sans «~je~» en Q1-Q2.
\item La relecture à froid se fait en trois passes : sens, vocabulaire, orthographe.
\item La Fiche Mémo A5 recto-verso est le seul document à réviser la veille du Bac.
\item Face à l'imprévu (texte court, long, auteur inconnu, thème non révisé), la méthode reste identique : c'est elle votre véritable assurance.
\end{itemize}
\end{aretenir}

\begin{exercice}[Pour s'entraîner à la maison]
Choisissez un texte philosophique d'une page (par exemple un extrait d'Achille Mbembe sur l'identité). Traitez les trois questions en 2 h 30 chronométrées, puis appliquez à votre propre copie la grille d'auto-évaluation de la figure 1 et la relecture en trois passes. Notez-vous honnêtement sur 20 et identifiez vos deux axes de progrès prioritaires avant la prochaine simulation.
\end{exercice}

\coursec{Transfert \& Citoyenneté : L'esprit critique hors les murs}

Après l'intensité de la simulation du Probatoire (Section 5), où vous avez mesuré votre maîtrise technique de l'exercice sur texte, il est temps de faire le pas de côté décisif : comprendre que la méthode philosophique n'est pas un « outil d'examen » que l'on range au placard après l'épreuve, mais une \textbf{disposition permanente de l'esprit}. Les trois gestes que vous avez répétés — \emph{expliciter}, \emph{analyser}, \emph{discuter} — sont exactement ceux que tout citoyen lucide et tout technicien responsable doit mobiliser face aux flux d'informations, aux choix éthiques et aux dilemmes professionnels. Cette section ancre vos acquis dans le réel camerounais d'aujourd'hui.

\courssub{De l'épreuve au réel : Analyser un article d'opinion camerounais}

La méthode « Texte $\to$ Explication $\to$ Analyse $\to$ Discussion » s'applique \emph{mutatis mutandis} à tout discours argumentatif : éditorial de \emph{Le Jour} ou \emph{Cameroon Tribune}, tribune sur un blog influent, post viral sur les réseaux sociaux. L'enjeu : ne pas subir l'opinion, mais la \emph{décortiquer}.

\begin{methode}[Analyse d'article de presse en 4 questions]
Appliquez cette grille à tout article d'opinion (éducation, corruption, jeunesse, diaspora, environnement) :
\begin{enumerate}
    \item \textbf{Quelle est la thèse implicite ?} Repérez la conclusion que l'auteur veut vous faire admettre sans toujours la formuler explicitement. Cherchez les marqueurs : « Il est urgent de… », « Force est de constater que… », « La solution est… ».
    \item \textbf{Quels arguments ?} Classez-les : \emph{factuels} (chiffres, rapports, observations vérifiables), \emph{émotionnels} (appel à la peur, à la fierté, à l'indignation), \emph{d'autorité} (citation d'un expert, d'un chef traditionnel, d'un texte sacré), \emph{logiques} (raisonnement cause-effet, analogie, \emph{reductio ad absurdum}).
    \item \textbf{Quels présupposés idéologiques ?} Quelles valeurs non dites structurent le raisonnement ? (Ex. : « Le développement passe nécessairement par l'industrialisation lourde » présuppose un productivisme ; « La jeunesse est l'avenir » peut masquer une injonction à la patience sans partage du pouvoir).
    \item \textbf{Ma discussion argumentée.} En mobilisant vos connaissances (philosophie, droit, économie, sciences), produisez : un \emph{accord partiel} (argument solide), une \emph{objection} (contre-exemple, faille logique, angle mort), une \emph{ouverture} (autre perspective, conséquence à long terme).
\end{enumerate}
\end{methode}

\begin{exemplebox}[Application : Éditorial fictif « L'urgence du tout-numérique à l'école »]
\textbf{Texte :} « Nos écoles manquent de tableaux, de craie, de manuels. La solution ? Équiper chaque élève d'une tablette. Le Rwanda l'a fait, pourquoi pas nous ? Le numérique est l'avenir, refusons le retard ! »
\begin{enumerate}
    \item \textbf{Thèse implicite :} L'équipement numérique \emph{suffit} à résoudre la crise éducative ; le modèle rwandais est transposable tel quel.
    \item \textbf{Arguments :} Factuel (manque de manuels) ; d'autorité/analogie (Rwanda) ; émotionnel (peur du retard, fierté nationale).
    \item \textbf{Présupposés :} L'infrastructure (électricité, internet, maintenance) existe ou suivra automatiquement ; le pédagogique se réduit à l'accès au contenu ; le coût total de possession est négligeable.
    \item \textbf{Discussion :} \emph{Accord :} Le manque de manuels est réel, le numérique offre des ressources immenses. \emph{Objection :} L'analogie avec le Rwanda occulte les différences de densité, de gouvernance, de budget par élève. Sans formation des enseignants, maintenance, électricité stable, la tablette devient un « gadget » coûteux (ex. : projet \emph{One Laptop Per Child} au Cameroun, résultats mitigés). \emph{Ouverture :} Une approche hybride (manuels + ressources numériques mutualisées en salle informatique solaire) serait plus résiliente et équitable.
\end{enumerate}
\end{exemplebox}

\begin{popfigure}
\begin{tikzpicture}[
    node distance=1.2cm and 1.5cm,
    box/.style={draw=popblue, thick, rounded corners, fill=popblueL, minimum width=2.8cm, minimum height=1.1cm, align=center, font=\small},
    arrow/.style={-Stealth, thick, popdark},
    labelbox/.style={font=\tiny, popmuted, align=center}
]
\node[box, align=center] (bac){Épreuve Bac\\Expliciter / Analyser / Discuter};
\node[box, right=of bac, align=center] (media){Analyse Médias\\Presse, Réseaux, Discours publics};
\node[box, right=of media, align=center] (citoyen){Débat Citoyen\\Argumenter, Écouter, Décider};
\node[box, below=of media, align=center] (pro){Éthique Professionnelle\\Responsabilité, Choix techniques, IA};
\node[box, below=of citoyen, align=center] (vie){Vie Personnelle\\Autonomie, Lucidité, Soin de soi};
\draw[arrow] (bac) -- (media) node[midway, above, labelbox, align=center]{Transfert\\(même méthode)};
\draw[arrow] (media) -- (citoyen) node[midway, above, labelbox, align=center]{Vigilance\\critique};
\draw[arrow] (citoyen) -- (pro) node[midway, right, labelbox, align=center]{Responsabilité\\ingénieur/technicien};
\draw[arrow] (pro) -- (vie) node[midway, right, labelbox, align=center]{Cohérence\\éthique};
\draw[arrow] (vie) -- (bac) node[midway, left, labelbox, align=center]{Formation\\continue};
\end{tikzpicture}
\legende{Schéma de transfert : la méthode philosophique irrigue l'analyse médiatique, le débat citoyen, l'éthique professionnelle et la vie personnelle.}
\end{popfigure}

\courssub{Le débat philosophique réglementé : Apprendre à penser ensemble}

En classe, nous transformons une thèse extraite d'un texte (ex. : « La corruption est une nécessité culturelle pour fluidifier l'administration ») en sujet de \textbf{débat réglementé}. Ce n'est pas une joute oratoire, mais un exercice de \emph{rationnalité partagée}.

\begin{definition}[Esprit critique (au sens kantien)]
\textbf{\emph{Sapere aude} — « Ose savoir par toi-même »} (Kant, \emph{Qu'est-ce que les Lumières ?}, 1784). L'esprit critique n'est pas l'esprit de contradiction stérile, c'est la conjonction de trois exigences :
\begin{enumerate}
    \item \textbf{Autonomie intellectuelle :} Refuser l'hétéronomie (penser par procuration : « On dit que… », « C'est la tradition », « Le chef a décidé »).
    \item \textbf{Exigence de preuves :} Accepter une thèse seulement si les arguments la soutenant résistent à l'examen (faits vérifiés, logique valide, absence de sophismes).
    \item \textbf{Acceptation de la contradiction :} Être prêt à changer d'avis si l'argument adverse est meilleur. La vérité se conquiert \emph{dans} le dialogue, pas avant.
\end{enumerate}
\end{definition}

\begin{methode}[Règles du débat philosophique réglementé (45 min)]
\begin{enumerate}
    \item \textbf{Présentation de la thèse} (2 min) : L'enseignant ou un élève lit la thèse extraite du texte.
    \item \textbf{Tour de table « Arguments pour »} (10 min) : Chaque intervenant \textbf{doit} commencer par : « Je reprends l'argument de X [précédent] : \emph{...}, avant d'avancer le mien : \emph{...} ». Cela force l'écoute active et la construction cumulative.
    \item \textbf{Tour de table « Arguments contre »} (10 min) : Même règle de reprise.
    \item \textbf{Distinction Fait / Opinion} (5 min) : L'enseignant note au tableau les énoncés ; la classe les classe : \emph{Fait vérifiable} (ex. « Le Cameroun a perdu 1500 milliards FCFA dans la corruption en 2023 selon le CONAC ») vs \emph{Opinion/Valuation} (ex. « La corruption est dans notre culture »).
    \item \textbf{Synthèse \& Positionnement final} (10 min) : Chaque élève écrit en 3 lignes : « Ma position actuelle : \emph{...} ; L'argument qui m'a fait bouger : \emph{...} ; Mon point de vigilance : \emph{...} ».
    \item \textbf{Méta-réflexion} (8 min) : « Qu'avons-nous appris sur la manière de discuter ? » (Ex. : « J'ai vu qu'un argument d'autorité ne tient pas », « J'ai distingué le problème structurel de la responsabilité individuelle »).
\end{enumerate}
\end{methode}

\begin{attention}[Piège : « La philosophie ne sert à rien pour un technicien »]
C'est l'erreur la plus coûteuse. Les jurys des grandes écoles d'ingénieurs (Polytechnique Yaoundé, ENSP, Mines) et les DRH des entreprises (SABC, Eneo, Camtel, banques, ONG) évaluent \textbf{la capacité à structurer un problème complexe et à argumenter une solution} — exactement ce que vous faites en Q\_2 et Q\_3. Un ingénieur qui ne sait pas \emph{problématiser} un cahier des charges, \emph{anticiper} les objections d'un client, \emph{justifier} un choix technique face à une alternative, est un exécutant, pas un concepteur. La philosophie est l'entraînement de ce « muscle » cognitif.
\end{attention}

\begin{exemplebox}[Enquête 2024 — Attentes des employeurs camerounais (simulation)]
Une enquête auprès de 120 DRH d'entreprises camerounaises (secteurs : BTP, télécoms, banque, agro-industrie, administration) révèle :
\begin{center}
\begin{tabularx}{\linewidth}{|l|X|c|}\hline
\rowcolor{popblueL}
\textbf{Compétence transversale} & \textbf{Description} & \textbf{\% DRH : Critère n°1 ou 2} \\ \hline
Capacité d'analyse et de synthèse & Décomposer un problème, hiérarchiser l'info, produire une note claire & 78\% \\ \hline
Qualité de l'expression écrite/orale & Rédiger un rapport, présenter un projet, argumenter en réunion & 78\% \\ \hline
Esprit critique / Résolution de problèmes & Identifier la cause racine, proposer des alternatives, décider & 71\% \\ \hline
Éthique professionnelle / Intégrité & Refuser la corruption, respecter les normes, assumer ses choix & 65\% \\ \hline
Travail en équipe / Communication & Écouter, reformuler, négocier, co-construire & 62\% \\ \hline
\end{tabularx}
\end{center}
\legende{Source : Enquête simulée type GICAM/INS 2024. Ces compétences sont \emph{exactement} celles travaillées en philosophie (Q\_1=analyse/synthèse, Q\_2=structure argumentative, Q\_3=discussion critique, Débat=écoute/réforme).}
\end{exemplebox}

\courssub{Écrire pour la cité : La « Tribune libre » philosophique}

Passer de l'analyse à la \emph{production} : rédiger une tribune de \textbf{300 mots maximum} (format standard presse) sur un thème d'actualité locale, soumise à la revue du lycée ou au journal scolaire (\emph{Le Lycéen}, \emph{Écho des Jeunes}, blog de l'établissement).

\begin{exoresolu}[Tribune libre : « Yaoundé, ville intelligente ? D'abord ville juste »]
\textbf{Thème :} Projet « Smart City » à Yaoundé (caméras, capteurs, data center) vs bidonvilles sans eau ni électricité.

\textbf{Tribune (298 mots) :}
« "Yaoundé ville intelligente" : le slogan brille sur les panneaux. Derrière le mot "intelligent", il y a des capteurs, de la 5G, des data centers. Mais l'intelligence d'une ville se mesure-t-elle au débit de sa fibre ou à la dignité de ses habitants ?

À Mvog-Ada, à Ekounou, à Nkolbisson, l'intelligence attend encore l'eau courante, l'évacuation des ordures, des trottoirs praticables pour les handicapés. Installer des caméras "intelligentes" sur des rues où les nids-de-poule dictent la vitesse, c'est confondre \emph{surveillance} et \emph{soin}. Hannah Arendt rappelle que la politique, c'est l'art de partager un monde commun. Un monde où 40\% de la population vit dans l'informel sans adresse officielle n'est pas "intelligent", il est \emph{aveugle}.

L'argument d'autorité ("Kigali l'a fait") ne vaut pas démonstration. Kigali a d'abord assaini, planifié, logé. La technique sans justice sociale produit une ville à deux vitesses : le centre "connecté" pour les nantis, la périphérie "oubliée" pour la majorité.

Vraie "Smart City" : budget participatif pour que les quartiers décident de leurs priorités ; données ouvertes (open data) sur les marchés publics pour lutter contre la corruption ; éclairage solaire dans les ruelles avant la 5G sur les ronds-points.

Ingénieurs, urbanistes, décideurs : l'intelligence artificielle ne remplacera jamais l'intelligence du cœur. Construisons d'abord une ville \emph{juste} ; elle sera alors, par surcroît, intelligente. »

\tcblower
\textbf{Analyse de la tribune :} Thèse claire (justice avant technologie) ; arguments : fait (bidonvilles), référence philosophique (Arendt), analogie critique (Kigali), proposition concrète (budget participatif, open data) ; ton citoyen, non injurieux, adressé aux décideurs.
\end{exoresolu}

\courssub{De l'écrit à l'oral : Bac, concours, entretiens professionnels}

L'oral du Probatoire (si rétabli) et, surtout, les entretiens d'admission (ENAM, IRIC, Polytech, banques, ONG) et d'embauche testent la \emph{même architecture mentale} :

\begin{center}
\begin{tabularx}{\linewidth}{|l|X|X|}\hline
\rowcolor{popblueL}
\textbf{Épreuve écrite (Texte)} & \textbf{Transposition orale} & \textbf{Clé de réussite} \\ \hline
Q\_1 : Expliciter le sens & \textbf{Reformuler la question / le sujet} & « Si je comprends bien, vous me demandez si… » (montre l'écoute active) \\ \hline
Q\_2 : Analyser l'argumentation & \textbf{Structurer sa réponse : Problème $\to$ Arguments $\to$ Exemples} & « Trois points : 1. Le constat… 2. La cause… 3. La solution… » (plan visible) \\ \hline
Q\_3 : Discuter & \textbf{Nuancer, objecter, ouvrir} & « Certes… Cependant… Dès lors… » (marqueurs logiques) \\ \hline
Débat réglementé & \textbf{Écoute active / Reformulation} & « Ce que vous dites, c'est que… Est-ce exact ? » (désamorce le conflit) \\ \hline
\end{tabularx}
\end{center}

\begin{methode}[Préparation « Entretien type » — 15 min / jour]
\begin{enumerate}
    \item \textbf{Sujet imposé} (ex. : « L'IA va-t-elle remplacer l'ingénieur ? », « Faut-il conditionner l'aide au développement à la démocratie ? »).
    \item \textbf{3 min :} Problématiser (quelle tension ?) + Plan en 3 parties (sur brouillon).
    \item \textbf{5 min :} S'enregistrer (téléphone) en répondant \emph{comme si} vous étiez face au jury.
    \item \textbf{7 min :} Réécoute critique : « Ai-je reformulé ? Mon plan est-il audible ? Mes exemples sont-ils concrets (Cameroun) ? Mon ton est-il posé ? ».
\end{enumerate}
\end{methode}

\courssub{Mon contrat de progression : Bilan personnel \& Route vers le Bac}

Transformez l'auto-évaluation de la Section 5 en un \textbf{contrat écrit, daté, signé}, affiché chez vous ou dans votre carnet de révisions.

\begin{popfigure}
\begin{tikzpicture}[
    cell/.style={draw=popline, minimum height=1cm, align=center, font=\small},
    header/.style={cell, fill=popblue, font=\bfseries, text=white},
    title/.style={font=\large\bfseries, popdark, anchor=north},
    section/.style={font=\bfseries, poppurple, anchor=west}
]
% Title
\node[title] at (0, 5.5) {MON CONTRAT DE PROGRESSION — PHILOSOPHIE TERMINALE TECH};
\node[title] at (0, 5.0) {Nom : \underline{\hspace{6cm}} \quad Date : \underline{\hspace{3cm}}};
% 3 Points forts
\node[section] at (-8, 4.3) {3 POINTS FORTS ACQUIS (je les maîtrise, je les mobilise)};
\foreach \i/\txt in {
    1/{Expliciter une thèse et des concepts sans paraphrase},
    2/{Repérer la structure d'un argumentaire (thèse/arguments/exemples)},
    3/{Rédiger une discussion en 3 temps (accord/objection/ouverture)}} {
    \node[cell, minimum width=16cm] at (0, 4.0-\i*0.7) {\i. \txt};
}
% 2 Points de vigilance
\node[section] at (-8, 1.8) {2 POINTS DE VIGILANCE (je dois encore travailler)};
\foreach \i/\txt in {
    1/{Introduire une référence philosophique pertinente \emph{au bon moment} (pas en "name-dropping")},
    2/{Gérer le temps en Q\_3 : ne pas tout dire, aller à l'essentiel}
} {
    \node[cell, minimum width=16cm] at (0, 1.5-\i*0.7) {\i. \txt};
}
% 1 Action concrète / semaine
\node[section] at (-8, -0.5) {1 ACTION CONCRÈTE PAR SEMAINE (jusqu'au Bac — \textbf{engagement signé})};
\node[cell, minimum width=16cm, fill=popgreenL, align=center] at (0, -1.2){
    \textbf{Lundi} : Lire 1 article philo (\emph{Le Philosophoire}, \emph{Philosophie Magazine}, \emph{The Conversation Africa}) + en faire une fiche « Thèse / 2 arguments / Mon objection » (20 min). \\
    \textbf{Mercredi} : Faire 1 plan complet de Q\_3 sur sujet Bac annales (30 min chrono) + auto-correction grille Section 5. \\
    \textbf{Samedi} : Participer au débat réglementé classe / écrire 1 micro-tribune (150 mots) sur l'actu locale.
};
% Signature
\node[font=\itshape, popmuted] at (0, -3.0) {« Je m'engage à tenir ce contrat. La philosophie n'est pas une matière, c'est ma boussole. »};
\node[font=\small, align=center] at (0, -3.6) {Signature : \underline{\hspace{5cm}} \quad Témoin (prof/parent/pair) : \underline{\hspace{4cm}}};
\end{tikzpicture}
\legende{Modèle de « Contrat de progression personnel » à imprimer, remplir, signer et afficher.}
\end{popfigure}

\courssub{Ouverture : La philosophie, « soin de l'âme » de l'ingénieur-citoyen}

\begin{propriete}[La méthode du texte philosophique : un entraînement à la raison critique universelle]
La méthode de l'exercice sur texte (Expliciter / Analyser / Discuter) n'est pas une « recette de cuisine » pour obtenir une note au Probatoire. C'est un \textbf{entraînement à la raison critique universelle}, indispensable à l'ingénieur/technicien face aux choix techniques à dimension éthique :
\begin{itemize}
    \item \textbf{Éthique de la responsabilité (Hans Jonas) :} L'ingénieur doit \emph{anticiper} les conséquences à long terme de ses choix techniques (ex. : algorithme de recrutement biaisé, gestion des déchets électroniques, IA de reconnaissance faciale). La Q\_3 (discussion) est l'entraînement direct à cette anticipation.
    \item \textbf{Éthique de l'IA / Données :} « Expliciter » = rendre un algorithme \emph{explicable} (explicabilité, \emph{white box} vs \emph{black box}) ; « Analyser » = détecter les biais dans les données d'entraînement ; « Discuter » = arbitrer entre performance et équité.
    \item \textbf{Développement durable :} Tout projet technique engage l'avenir. La méthode philosophique oblige à \emph{problématiser} (besoin réel vs solution technique), à \emph{peser} les arguments (coût/bénéfice, justice intergénérationnelle), à \emph{assumer} un choix argumenté.
\end{itemize}
\end{propriete}

\begin{savaistu}[Compétences « Expliciter / Analyser / Discuter » = Tests psychotechniques grands concours]
Les épreuves d'aptitude (ENAM, IRIC, Polytech, concours banque/assurance, fonction publique) comportent quasi systématiquement :
\begin{itemize}
    \item \textbf{Compréhension de texte :} « Quelle est la thèse de l'auteur ? » (Q\_1) ; « Quel argument utilise-t-il ? » (Q\_2).
    \item \textbf{Raisonnement logique / Critique :} « Laquelle de ces propositions affaiblit l'argument ? » (Q\_3 objection) ; « Quelle conclusion peut-on tirer ? » (Q\_3 ouverture).
    \item \textbf{Expression écrite :} « Rédigez une note de synthèse / une prise de position argumentée. » (Tribune libre, Dissertation).
\end{itemize}
\emph{Travailler la philosophie en Terminale Technique, c'est préparer ces concours en amont, sans surcoût.}
\end{savaistu}

\begin{exemplebox}[Cas concret : Ingénieur en génie civil — Choix du tracé d'une route]
\textbf{Situation :} Deux tracés possibles. Tracé A : court, moins cher, traverse une zone humide protégée + village ancestral. Tracé B : plus long, 15\% plus cher, évite zone humide et village.
\begin{enumerate}
    \item \textbf{Expliciter (Q\_1) :} Quel est le \emph{vrai} problème ? Pas « quel tracé choisir ? » mais « Comment arbitrer entre efficacité économique, respect de la biodiversité et droits des populations locales ? ».
    \item \textbf{Analyser (Q\_2) :} Arguments pour A : coût, délai, rentabilité immédiate. Arguments pour B : durabilité (éviter inondations futures), justice (consentement libre, préalable, éclairé des villageois), conformité loi (code de l'environnement), réputation long terme.
    \item \textbf{Discuter (Q\_3) :} \emph{Accord partiel :} Le coût est une contrainte réelle. \emph{Objection :} L'analyse coûts-avantages sur 30 ans (maintenance, contentieux, image) favorise souvent B. \emph{Ouverture :} Et si on innovait ? Pont éco-conçu sur la zone humide + fonds de développement villageois financé par l'économie du tracé A ?
\end{enumerate}
L'ingénieur qui a \emph{intériorisé} la méthode philo ne livre pas un calcul de rentabilité brute ; il livre une \emph{recommandation argumentée, éthiquement fondée, juridiquement robuste}. C'est \emph{ça}, la valeur ajoutée du philosophe-technicien.
\end{exemplebox}

\begin{aretenir}[Bilan final : La philosophie, votre « système d'exploitation » mental]
\begin{itemize}
    \item \textbf{\cle{Expliciter}} = Déboguer le sens (éviter les malentendus, les contresens, les « fake news »).
    \item \textbf{\cle{Analyser}} = Comprendre l'architecture d'un raisonnement (repérer les failles, les forces, les présupposés).
    \item \textbf{\cle{Discuter}} = Produire un jugement argumenté, nuancé, responsable (ni dogmatisme, ni relativisme).
    \item \textbf{\cle{Transférer}} = Appliquer ce « OS » à l'analyse médiatique, au débat citoyen, à l'éthique pro, à l'oral, à la vie.
\end{itemize}
Le Probatoire vérifiera votre maîtrise technique. La vie vérifiera votre maîtrise existentielle. Les deux demandent la même rigueur. \textbf{À vous de jouer.}
\end{aretenir}

\newpage

\coursec{Activité d'intégration}

\coursec{Architecture de l'épreuve et gestion stratégique du temps}

\courssub{Objectifs de l'activité}
\begin{itemize}
    \item \cle{Mobiliser} l'intégralité de la chaîne méthodologique de l'exercice sur texte (Question 1 : explicitation du sens ; Question 2 : analyse de l'argumentation ; Question 3 : discussion critique) dans une situation de \emph{transfert} complexe et ludique.
    \item \cle{Développer} l'esprit critique et l'argumentation orale et écrite en situation de contradiction contrôlée (plaidoirie, contre-interrogatoire, délibération).
    \item \cle{Ancrer} la réflexion philosophique dans le réel camerounais (contexte socio-économique, culturel, technologique) pour vérifier la \cle{pertinence} et l'\cle{opérativité} des thèses philosophiques.
    \item \cle{Co-construire} une grille d'évaluation partagée (par le Jury) pour intérioriser les critères de réussite du Probatoire et du Baccalauréat technique : rigueur conceptuelle, articulation logique, ancrage concret, nuance.
    \item \cle{Renforcer} la posture citoyenne : écoute active, respect de la contradiction, recherche du consensus nuancé plutôt que victoire rhétorique.
\end{itemize}

\begin{aretenir}[Rappel : l'architecture de l'épreuve sur texte (3 heures)]
L'exercice sur texte philosophique au Probatoire et au Baccalauréat technique comporte \textbf{trois questions} hiérarchisées :
\begin{enumerate}
    \item \textbf{Question 1 — Vérification des savoirs} : dégager la thèse du texte, définir les concepts-clés, dégager la structure (les mouvements de la pensée de l'auteur).
    \item \textbf{Question 2 — Vérification des savoir-faire} : analyser l'argumentation (identifier les prémisses et la conclusion, le type de raisonnement, les présupposés).
    \item \textbf{Question 3 — Discussion critique} : apprécier la pertinence de la thèse en la confrontant à d'autres points de vue et à l'expérience concrète, puis produire une position personnelle justifiée.
\end{enumerate}
\textbf{Gestion stratégique du temps (3 h) :} environ 45 min pour la Question 1, 45 min pour la Question 2, 1 h 30 pour la Question 3 (brouillon puis rédaction au propre). La Question 3, la plus valorisée, ne doit jamais être sacrifiée.
\end{aretenir}

\coursec{Maîtrise des trois questions : rappels méthodologiques}

\courssub{Ressources à mobiliser}
\begin{definition}[Méthode Q1 : Explicitation du sens (Savoir)]
\textbf{4 étapes incontournables :}
\begin{enumerate}
    \item \textbf{Contextualiser} : auteur (s'il est connu), époque, thème du programme (ex. : \cle{Technique}, \cle{Travail}, \cle{Société}).
    \item \textbf{Formuler la thèse} : une phrase unique, affirmative, reprenant le vocabulaire du texte.
    \item \textbf{Dégager les concepts-clés} (3 ou 4) : définition philosophique \emph{in situ} (le sens dans le texte), et non définition de dictionnaire.
    \item \textbf{Reconstituer la structure} : découpage en \cle{mouvements} (idées-forces) reliés par les connecteurs logiques du texte (\emph{donc, or, mais, par conséquent}).
\end{enumerate}
\emph{Pièges :} la paraphrase, le commentaire personnel, la citation sans analyse.
\end{definition}

\begin{definition}[Méthode Q2 : Analyse de l'argumentation (Savoir-faire)]
\textbf{3 opérations :}
\begin{enumerate}
    \item \textbf{Identifier les arguments} : distinguer les \cle{prémisses} (faits, définitions, autorités, analogies) et la \cle{conclusion}.
    \item \textbf{Analyser l'enchaînement} : type de raisonnement (\cle{déductif}, \cle{inductif}, \cle{analogique}, ou raisonnement par \cle{conditions de possibilité}). Repérer les connecteurs logiques.
    \item \textbf{Mettre au jour les présupposés} : ce qui n'est pas dit mais qui est nécessaire au raisonnement (ex. : « la technique n'est pas neutre »).
\end{enumerate}
\emph{Exigence du Bac :} formaliser au moins un raisonnement sous forme canonique (syllogisme ou schéma du type : si P alors Q ; or P ; donc Q).
\end{definition}

\begin{definition}[Méthode Q3 : Discussion critique (Compétence supérieure)]
\textbf{Structure dialectique rigoureuse (thèse -- antithèse -- synthèse \emph{problématisée}) :}
\begin{enumerate}
    \item \textbf{Problématique} : issue de la tension interne au texte (ex. : « Jusqu'où la maîtrise technique peut-elle résister à la rationalité gestionnaire algorithmique ? »).
    \item \textbf{Thèse (appui au texte)} : développer la force de la pensée de l'auteur avec des \textbf{arguments nouveaux} et des \textbf{exemples concrets (Cameroun)}.
    \item \textbf{Antithèse (limites et objections)} : objections \emph{philosophiques} (autres auteurs, concepts contraires) et \emph{factuelles} (réalité camerounaise qui contredit ou complexifie la thèse).
    \item \textbf{Synthèse (dépassement)} : position nouvelle intégrant les deux moments. Critères : \cle{nuance}, \cle{hiérarchisation}, \cle{ouverture}.
\end{enumerate}
\emph{Interdits :} le plan « pour / contre » binaire, le catalogue d'auteurs, les exemples hors sujet, la conclusion « tout est relatif ».
\end{definition}

\coursec{Atelier de consolidation : « Le Procès du Texte »}

\courssub{Situation}
\textbf{Contexte :} Dans le cadre de la préparation à l'épreuve de Philosophie du Baccalauréat technique (durée 3 h), la classe de Terminale d'un lycée technique de Douala aborde la leçon de consolidation n\textsuperscript{o} 3. Le professeur propose une mise en situation immersive : \textbf{« Le Procès du Texte »}.

\textbf{Le texte « accusé » (extrait court, inédit, inspiré des débats contemporains sur la technique en Afrique) :}
\begin{quote}
\emph{« L'irruption de l'intelligence artificielle générative dans nos ateliers et nos bureaux camerounais ne saurait être réduite à un simple transfert d'outils. Elle instaure une nouvelle rationalité gestionnaire : l'optimisation algorithmique devient la norme suprême, reléguant le savoir-faire du technicien, l'intuition de l'artisan et la solidarité de l'équipe au rang de "variables d'ajustement".}

\emph{À Douala comme à Yaoundé, le jeune ingénieur formé dans nos grandes écoles ne "conçoit" plus ; il interroge une machine. Il ne diagnostique plus la panne par l'écoute du moteur ou l'analyse du schéma ; il interroge la base de données. La technicité --- cette alliance de la main, de l'\oe il et de l'esprit forgée par l'expérience --- se dissout dans l'instantanéité de la réponse probabiliste.}

\emph{Pire, cette délégation cognitive fragilise notre souveraineté technique. Nos données industrielles, nos habitudes de maintenance, nos spécificités climatiques (humidité, poussière, variations de tension) alimentent des modèles entraînés ailleurs, qui nous renvoient des solutions standardisées, aveugles à notre réalité. Accepter sans critique cette "boîte noire", c'est accepter une nouvelle forme de dépendance structurelle, masquée sous le vernis de la modernité.}

\emph{Le progrès technique authentique ne se mesure pas à la vitesse du processeur, mais à la capacité du peuple technique à maîtriser son outil, à le réparer, à le détourner, à l'adapter. L'IA, dans sa forme propriétaire actuelle, est un dispositif de capture : elle capture notre attention, nos données, et in fine, notre autonomie de jugement professionnel. Résister, ce n'est pas refuser l'outil, c'est exiger sa transparence, sa réparabilité et sa subordination au projet humain collectif. »}
\end{quote}

\textbf{L'Accusation (thèse du Procureur) :} \emph{« Ce texte est dangereux et dépassé pour le Cameroun d'aujourd'hui. Il verse dans la technophobie romantique, ignore l'urgence du développement, diabolise un levier de productivité indispensable pour notre industrialisation (usines de transformation du cacao et du bois, zone économique de Kribi) et propose une "maîtrise" illusoire face à une technologie globale. Il décourage la jeunesse d'acquérir les compétences numériques les plus demandées sur le marché de l'emploi. »}

\textbf{La Défense (thèse des Avocats) :} \emph{« Ce texte est un manifeste de lucidité et de responsabilité. Il ne refuse pas la technique ; il en dénonce la dérive quand elle échappe à la maîtrise humaine et politique. Il protège le c\oe ur de métier du technicien camerounais : le diagnostic situé, l'adaptation au terrain, l'éthique de la maintenance. Il est actuel car il pose la condition de possibilité d'une industrialisation endogène et non mimétique. »}

\courssub{Consignes}
\textbf{Phase 1 : Préparation en groupes (30 min --- brouillons structurés obligatoires)}

\textbf{Groupe 1 : LA DÉFENSE (Avocats de l'auteur --- 8 à 10 élèves)}
\begin{enumerate}
    \item \textbf{Q1 --- Explicitation rigoureuse (Savoir) :} rédiger une fiche « Identité du texte » : thèse principale (une phrase), 3 concepts-clés définis (\cle{rationalité gestionnaire}, \cle{technicité}, \cle{dispositif de capture}), structure logique (3 mouvements : constat $\rightarrow$ analyse $\rightarrow$ prescription).
    \item \textbf{Q2 --- Reconstruction argumentative (Savoir-faire) :} formaliser l'argumentation en \cle{syllogismes} ou \cle{chaînes causales}. Exemple : « Prémisse 1 : l'IA générative impose l'optimisation algorithmique comme norme. Prémisse 2 : cette norme exclut le jugement situé. Conclusion : l'autonomie du technicien est menacée. » Identifier 2 \cle{présupposés} forts (ex. : « le savoir-faire manuel est irréductible au code », « la souveraineté technique est condition du développement »).
    \item \textbf{Anticipation :} lister 3 objections prévisibles de l'Accusation et préparer les \cle{parades} (arguments de secours, exemples concrets camerounais : maintenance des groupes électrogènes, adaptation de machines-outils d'occasion, formation en alternance).
\end{enumerate}

\textbf{Groupe 2 : L'ACCUSATION (Procureurs --- 8 à 10 élèves)}
\begin{enumerate}
    \item \textbf{Analyse critique (Q2 inversé / Q3) :} identifier 3 \cle{failles logiques} ou \cle{angles morts} : \emph{généralisation abusive} (toute IA est-elle une boîte noire ?), \emph{fausse alternative} (IA contre artisanat : pourquoi pas un usage hybride ?), \emph{anachronisme} (idéalisation du technicien « pur » d'avant).
    \item \textbf{Contre-exemples réels camerounais :} documenter 3 cas où le numérique a \cle{résolu} un problème local difficilement soluble autrement (ex. : applications de diagnostic des maladies des plantes pour les agriculteurs, optimisation logistique dans la filière cacao, maintenance prédictive dans les entreprises publiques d'énergie).
    \item \textbf{Argumentaire « Développement » :} construire l'argument : « Refuser l'IA propriétaire, c'est rater le train de la quatrième révolution industrielle, donc accroître la dépendance \emph{économique} ». Préparer des \cle{questions pièges} pour la Défense (ex. : « Comment former massivement les jeunes sans les cours en ligne et les outils numériques ? »).
\end{enumerate}

\textbf{Groupe 3 : LE JURY (Juges et greffiers --- 6 à 8 élèves)}
\begin{enumerate}
    \item \textbf{Grille d'évaluation (méta-cognition) :} concevoir une grille notée sur 20 reprenant les attendus officiels du Bac : \emph{Respect du texte (Q1) / 5 pts}, \emph{Qualité de l'analyse argumentative (Q2) / 5 pts}, \emph{Pertinence de la discussion (Q3) / 5 pts}, \emph{Expression / 5 pts}. Définir des \cle{descripteurs de performance} pour chaque niveau (ex. : Q1, 5 pts = thèse exacte + concepts définis + structure claire ; 2 pts = résumé approximatif).
    \item \textbf{Préparation de l'audience :} rédiger 3 \cle{questions de clarification} neutres à poser aux deux parties (ex. : « Défense, vous parlez de "technicité" : cela exclut-il la programmation ? », « Accusation, la "productivité" suffit-elle à définir le progrès ? »).
    \item \textbf{Modèle de verdict :} pré-rédiger la structure de la synthèse finale (type Q3) : \emph{problématique commune} $\rightarrow$ \emph{thèse de la Défense (nuancée)} $\rightarrow$ \emph{thèse de l'Accusation (nuancée)} $\rightarrow$ \emph{synthèse de dépassement (verdict)}.
\end{enumerate}

\textbf{Phase 2 : Audience publique (40 min)}
\begin{itemize}
    \item \emph{Ouverture (2 min) :} le président du Jury rappelle la procédure.
    \item \emph{Plaidoirie de la Défense (8 min) :} exposé structuré Q1 + Q2 (idéalement sans notes).
    \item \emph{Réquisitoire de l'Accusation (8 min) :} exposé des failles, contre-exemples, enjeux de développement.
    \item \emph{Contre-interrogatoires croisés (12 min) :} le Jury pose ses questions, les parties répondent, répliques brèves (2 min chacune).
    \item \emph{Plaidoyers de clôture (5 min) :} une phrase finale par camp.
\end{itemize}

\textbf{Phase 3 : Délibération et rédaction du verdict (20 min --- Jury seul)}
\begin{itemize}
    \item Le Jury délibère à partir de la grille et rédige collectivement le \textbf{verdict officiel} (format Q3 : 300 à 400 mots) : \emph{acquittement} (texte solide), \emph{condamnation avec sursis} (texte pertinent mais lacunaire sur un point précis), ou \emph{peine alternative} (texte à réviser : intégrer l'hybridation IA / humain).
\end{itemize}

\textbf{Phase 4 : Débriefing méta-cognitif (10 min --- classe entière)}
\begin{itemize}
    \item Tour de table : « Quelle étape (Q1, Q2, Q3) m'a semblé la plus difficile en situation réelle ? », « Qu'ai-je compris du \emph{lien} entre explication et discussion ? », « En quoi le fait de \emph{juger} m'a-t-il aidé à mieux \emph{répondre} ? ».
\end{itemize}

\coursec{Démarche et corrigé commenté}

\begin{exoresolu}[Simulation guidée : extrait du verdict du Jury et analyse méthodologique]
\textbf{Énoncé :} \emph{À l'issue de l'audience, le Jury (Groupe 3) rend le verdict suivant. Analysez la qualité méthodologique de ce verdict au regard des attendus des Questions 1, 2 et 3.}

\textbf{VERDICT DU JURY (extrait) :}

« \emph{Le texte soumis à notre appréciation soutient que l'IA générative menace l'autonomie du technicien camerounais en imposant une rationalité gestionnaire opaque (thèse, Q1). L'argumentation repose sur l'opposition entre la technicité (savoir incorporé, situé) et l'algorithmique (calcul probabiliste, décontextualisé), et sur le risque de dépendance structurelle via la capture des données locales par des modèles externes (analyse, Q2).}

\emph{La Défense a solidement établi que la technicité est le c\oe ur de la formation technique (ex. : ajustage, maintenance préventive sur machines-outils d'occasion au lycée technique), irréductible à la simple formulation de requêtes. L'Accusation a pertinemment objecté que les logiciels libres et les modèles légers, utilisables localement sur des machines modestes, permettent une appropriation locale, ce qui contredit le déterminisme de la "boîte noire propriétaire". Elle a cité l'exemple du diagnostic des maladies des plantes par application mobile comme augmentation, et non substitution, de l'expertise paysanne.}

\emph{Le Jury reconnaît la lucidité du texte sur les risques d'aliénation cognitive et de dépendance des données (thèse validée). Mais il juge l'antithèse décisive sur le plan factuel : la frontière "IA contre artisanat" est poreuse ; l'enjeu n'est pas le refus mais la gouvernance (données ouvertes, formation critique aux outils, droit à la réparation).}

\emph{\textbf{Verdict : CONDAMNATION AVEC SURSIS.} Le texte est philosophiquement pertinent (éthique de la responsabilité technique) mais sociologiquement incomplet (il sous-estime la plasticité des usages numériques au Cameroun). \textbf{Peine alternative :} réécriture de la conclusion intégrant l'impératif d'une technologie conviviale, au sens d'Ivan Illich, adaptée au contexte : "L'IA doit être un levier pour la technicité, non son fossoyeur. Cela suppose des outils ouverts, des données locales et une formation hybride."} »

\tcblower
\textbf{Analyse commentée du verdict (corrigé type pour l'enseignant) :}

\emph{1. Respect de la structure Q3 (problématique implicite $\rightarrow$ thèse $\rightarrow$ antithèse $\rightarrow$ synthèse) :} \textbf{excellent.} Le verdict ne tranche pas de manière binaire. Il \cle{hiérarchise} : validité philosophique (thèse) contre limite factuelle (antithèse). La synthèse propose un \cle{concept opératoire nouveau} (« technologie conviviale », « gouvernance des données ») qui dépasse le conflit initial.

\emph{2. Mobilisation de Q1 (explicitation) :} \textbf{réussie.} Le Jury résume la thèse et les concepts (\cle{technicité}, \cle{rationalité gestionnaire}, \cle{dispositif de capture}) avec précision. Il montre qu'il a \emph{compris} le texte avant de le juger. \emph{Point de vigilance :} la structure en trois mouvements n'est pas explicitée, seulement implicite dans l'analyse.

\emph{3. Mobilisation de Q2 (analyse argumentative) :} \textbf{solide.} Le Jury identifie le \cle{noyau argumentatif} : l'opposition technicité / algorithmique. Il évalue la \cle{force probante} des deux camps : la Défense ancre son argument dans le réel de la formation (ajustage), l'Accusation contre-argumente par la technique elle-même (logiciels libres, modèles légers). Le Jury valide le \cle{présupposé} (la technique n'est pas neutre) mais conteste l'\cle{inférence} sociologique (le caractère inéluctable de la capture).

\emph{4. Ancrage camerounais (transfert) :} \textbf{exemplaire.} Les exemples ne sont pas décoratifs : \emph{ajustage et machines d'occasion au lycée technique} (réalité de la formation), \emph{applications mobiles agricoles et logiciels libres} (réalité de l'innovation locale), \emph{Illich et la technologie conviviale} (référence philosophique pertinente pour le thème « Technique » du programme).

\emph{5. Expression et argumentation :} \textbf{niveau Bac.} Vocabulaire précis (\cle{gouvernance}, \cle{aliénation cognitive}, \cle{appropriation}), connecteurs logiques (\emph{mais, toutefois, dès lors}), ton objectif et mesuré.

\emph{Conclusion pédagogique :} ce verdict modélise la \cle{compétence visée} : non pas « avoir un avis », mais \textbf{juger en philosophe} : expliciter fidèlement, analyser la solidité du raisonnement, confronter au réel, produire une pensée nouvelle et nuancée. C'est la \cle{maîtrise de la méthode} qui est évaluée, bien plus que l'opinion sur l'IA.
\end{exoresolu}

\coursec{Transfert et citoyenneté : l'esprit critique hors les murs}

\courssub{Bilan de l'activité}
\begin{aretenir}[Bilan de consolidation --- ce que l'activité a mis en évidence]
\begin{itemize}
    \item \textbf{La méthode n'est pas une liste de cases à cocher, mais une \cle{posture intellectuelle}.} Le « Procès » a forcé les élèves à \emph{habiter} Q1 (fidélité), Q2 (rigueur logique) et Q3 (esprit de nuance) non comme des formalités, mais comme des \cle{outils} au service d'une vérité disputée.
    \item \textbf{L'\cle{ancrage contextuel} (Cameroun) est le révélateur de la pertinence philosophique.} Un texte qui ne résiste pas à l'épreuve du réel (ex. : l'application mobile qui aide réellement le paysan) est un texte incomplet. La discussion critique (Q3) \emph{exige} des exemples concrets, datés, localisés.
    \item \textbf{La \cle{contradiction} est productive.} L'Accusation a obligé la Défense à préciser ses concepts (qu'est-ce que la « technicité » face à la programmation ?). Le Jury a obligé les deux camps à dépasser leurs positions initiales. C'est la dialectique en acte : \emph{thèse $\rightarrow$ antithèse $\rightarrow$ synthèse}.
    \item \textbf{L'\cle{auto-évaluation} (grille du Jury) est le levier méta-cognitif majeur.} En construisant eux-mêmes la grille, les élèves ont intériorisé les \cle{critères de réussite du Bac} : thèse exacte, concepts définis, enchaînement logique, exemples pertinents, nuance. Ils savent désormais \emph{quoi} réviser et \emph{comment} s'auto-corriger.
    \item \textbf{Transférabilité immédiate :} cette simulation est la répétition générale de l'épreuve de 3 h : gestion du temps, structuration du brouillon, articulation de l'oral et de l'écrit, gestion du stress contradictoire.
    \item \textbf{Ouverture citoyenne :} hors les murs de la classe, la même démarche (comprendre avant de juger, analyser les arguments, discuter avec nuance) s'applique aux débats de la vie courante : rumeurs sur les réseaux sociaux, discours publicitaires, débats publics sur le développement du pays. L'exercice sur texte forme ainsi le \cle{jugement citoyen}.
\end{itemize}
\end{aretenir}

\begin{attention}[Pièges résiduels à surveiller pour le jour J]
\begin{itemize}
    \item \textbf{Q1 :} ne pas confondre « thèse de l'auteur » et « mon avis ». Ne pas lister des concepts sans les \emph{définir dans le texte}.
    \item \textbf{Q2 :} ne pas faire un simple « résumé des arguments ». Il faut \emph{analyser la structure} : prémisses $\rightarrow$ conclusion, type de raisonnement, présupposés. Formaliser (syllogisme) est un \cle{gage de rigueur} très valorisé.
    \item \textbf{Q3 :} le plan « avantages / inconvénients » est \textbf{à proscrire}. Exiger : problématique $\rightarrow$ appui critique $\rightarrow$ objection fondée $\rightarrow$ dépassement conceptuel.
    \item \textbf{Exemples :} « De nos jours... », « Dans la vie courante... » ne valent rien. Exiger la précision : « Au Cameroun, dans le secteur informel, dans les lycées techniques, dans la filière cacao, l'exemple précis X montre que... ».
    \item \textbf{Temps :} 3 h = environ 45 min pour Q1, 45 min pour Q2, 1 h 30 pour Q3 (brouillon + rédaction au propre). Ne pas « voler » du temps à Q3 pour peaufiner Q1.
\end{itemize}
\end{attention}

\newpage

\coursec{Méthodes et exercices résolus}

\courssub{Fiche méthode 1 : Lire et comprendre le texte avant toute réponse}

\begin{definition}[L'exercice sur texte philosophique]
L'\cle{exercice sur texte philosophique} est une épreuve écrite du Probatoire et du Baccalauréat technique qui consiste à rendre compte, de manière ordonnée et rédigée, de la pensée d'un auteur à partir d'un extrait d'œuvre. Il comporte classiquement trois questions : la \cle{Question 1} vérifie les \emph{savoirs} (compréhension et explication du sens du texte), la \cle{Question 2} vérifie les \emph{savoir-faire} (analyse de la démarche et de l'argumentation), la \cle{Question 3} vérifie la \emph{compétence supérieure} (discussion critique et position personnelle justifiée).
\end{definition}

\begin{methode}[Lecture méthodique d'un texte philosophique en trois passages]
\begin{enumerate}
\item \textbf{Première lecture (globale)} : lire le texte sans surligner, identifier l'auteur, l'œuvre et le thème général (la vérité, la liberté, l'État, le devoir...). Se demander : « De quoi parle ce texte ? »
\item \textbf{Deuxième lecture (analytique)} : souligner la \cle{thèse} (ce que l'auteur affirme), encadrer les \cle{arguments} (ce qui justifie la thèse), repérer les connecteurs logiques (\emph{car, donc, cependant, en effet}) et les exemples.
\item \textbf{Troisième lecture (critique)} : relever les mots difficiles ou ambigus, les présupposés de l'auteur, et formuler une première objection personnelle.
\item Rédiger au brouillon un plan du texte en 2 ou 3 moments : introduction de la thèse, développement, conclusion.
\item Vérifier que chaque question posée est reliée à un passage précis du texte avant de répondre.
\end{enumerate}
\textbf{Réflexes} : ne jamais commencer à rédiger avant d'avoir identifié la thèse ; toujours citer le texte entre guillemets pour justifier une réponse ; compter environ 20 minutes de lecture active sur 3 heures d'épreuve.
\end{methode}

\begin{aretenir}[Le vocabulaire minimal de l'analyse de texte]
\begin{itemize}
\item La \cle{thèse} : la position défendue par l'auteur, souvent signalée par un connecteur conclusif (\emph{donc, ainsi, par conséquent}).
\item L'\cle{argument} : une raison avancée pour établir la thèse.
\item La \cle{démarche} : l'ordre dans lequel les arguments s'enchaînent (déduction, induction, renversement, concession).
\item L'\cle{enjeu} : la raison pour laquelle la question traitée est importante pour l'homme ou la société.
\end{itemize}
\end{aretenir}

\begin{exoresolu}[Identifier la thèse et la structure d'un texte]
\textbf{Énoncé (type Bac)} : On donne le texte suivant : « L'homme n'est point fait pour la méditation, disent les partisans de l'action. Mais c'est oublier que toute action véritable suppose d'abord une délibération. Celui qui agit sans réfléchir n'agit pas : il est agi. La pensée n'est donc pas l'ennemie de l'action ; elle en est la condition. »
\begin{enumerate}
\item Dégagez la thèse de l'auteur.
\item Relevez deux arguments qui la soutiennent.
\end{enumerate}
\tcblower
\textbf{Solution détaillée.}
\begin{enumerate}
\item \textbf{Thèse} : la pensée (la réflexion) n'est pas opposée à l'action ; elle en est la condition nécessaire. Formulation possible : « L'auteur soutient que l'action authentique est impossible sans réflexion préalable. » On la repère grâce au connecteur conclusif « donc » dans la dernière phrase, qui condense la position de l'auteur.
\item \textbf{Argument 1} : « toute action véritable suppose d'abord une délibération » — argument de définition : une action qui mérite ce nom inclut un choix réfléchi.
\textbf{Argument 2} : « celui qui agit sans réfléchir n'agit pas : il est agi » — argument par distinction conceptuelle entre \emph{agir} (être cause de ses gestes) et \emph{être agi} (subir des causes extérieures, comme un automate).
\end{enumerate}
\textbf{Conclusion} : la thèse est que la réflexion conditionne l'action véritable, soutenue par un argument de définition et une distinction entre agir et être agi. Toute la suite du devoir s'appuiera sur cette identification.
\end{exoresolu}

\courssub{Fiche méthode 2 : Réussir la Question 1 — vérification des savoirs}

\begin{methode}[Expliquer le sens du texte (Question de savoir)]
\begin{enumerate}
\item \textbf{Reformuler la thèse} en une phrase claire, avec vos propres mots, sans paraphraser mot à mot.
\item \textbf{Définir les notions-clés} du texte (2 ou 3 termes au maximum) : donner une définition précise, éventuellement un synonyme et un contraire.
\item \textbf{Expliquer les articulations} : montrer comment chaque argument s'enchaîne (opposition, conséquence, exemple).
\item \textbf{Appuyer chaque affirmation} par une courte citation du texte entre guillemets, suivie de son explication.
\item \textbf{Dégager l'enjeu} : conclure en disant pourquoi cette thèse est importante (pour l'homme, la société, la morale...).
\end{enumerate}
\textbf{Formules à retenir} : « L'auteur affirme que... », « Ce passage signifie que... », « La notion de X désigne ici... », « L'enjeu de ce texte est donc... ». \textbf{Piège} : la Question 1 n'appelle ni avis personnel ni critique : on reste fidèle à la pensée de l'auteur.
\end{methode}

\begin{exoresolu}[Explication de texte — Question 1]
\textbf{Énoncé (type Bac)} : Texte : « La coutume est une seconde nature qui détruit la première. Pourquoi ce qui est naturel serait-il moins coutumier ? J'ai grand peur que cette nature ne soit elle-même qu'une première coutume, comme la coutume est une seconde nature. » (d'après Pascal)
\begin{enumerate}
\item Expliquez le sens de ce texte en définissant les notions de \emph{nature} et de \emph{coutume}.
\item Dégagez l'enjeu du texte.
\end{enumerate}
\tcblower
\textbf{Solution détaillée.}
\begin{enumerate}
\item \textbf{Reformulation de la thèse} : l'auteur soutient que ce que nous appelons « naturel » est en réalité le produit de l'habitude ; la frontière entre nature et coutume est donc incertaine.
\textbf{Définition de la nature} : ensemble des dispositions innées de l'homme, ce qu'il est avant toute éducation (synonyme : l'inné ; contraire : l'acquis).
\textbf{Définition de la coutume} : ensemble des habitudes acquises par la répétition et l'éducation, au point de paraître naturelles — c'est pourquoi l'auteur la nomme « seconde nature ».
\textbf{Articulation} : le texte procède par renversement progressif. D'abord, la coutume « détruit » la nature première (elle la recouvre). Ensuite, une question rhétorique (« Pourquoi ce qui est naturel serait-il moins coutumier ? ») inverse le rapport : et si la nature elle-même n'était qu'une habitude oubliée, une « première coutume » ? La citation « J'ai grand peur que cette nature ne soit elle-même qu'une première coutume » exprime le doute de l'auteur : nous ne pouvons plus distinguer l'inné de l'acquis.
\item \textbf{Enjeu} : si nos prétendues évidences naturelles (nos goûts, nos préjugés, nos jugements moraux) ne sont que des habitudes, alors l'esprit critique devient nécessaire : il faut interroger ce que nous tenons pour « naturel ». Par exemple, au Cameroun, certaines pratiques présentées comme « naturelles » (rôles attribués aux femmes et aux hommes) relèvent en partie de la coutume et peuvent être discutées.
\end{enumerate}
\textbf{Conclusion} : le texte signifie que la coutume recouvre et peut même constituer ce que nous croyons être notre nature ; son enjeu est de nous inviter à douter de nos évidences.
\end{exoresolu}

\courssub{Fiche méthode 3 : Réussir la Question 2 — vérification des savoir-faire}

\begin{methode}[Analyser l'argumentation d'un texte (Question de savoir-faire)]
\begin{enumerate}
\item \textbf{Identifier le type de chaque argument} : argument d'autorité, par l'exemple, par analogie, par l'absurde, par distinction, par élimination, argument \emph{a fortiori}.
\item \textbf{Reconstituer la démarche} : indiquer si le raisonnement est déductif (du général au particulier) ou inductif (des cas particuliers à une loi générale).
\item \textbf{Repérer les procédés rhétoriques} : question rhétorique, opposition, gradation, concession (« certes... mais »).
\item \textbf{Évaluer la validité} : l'argument prouve-t-il réellement la thèse ? L'exemple est-il représentatif ? Y a-t-il un présupposé caché ?
\item \textbf{Illustrer} chaque analyse par une citation précise du texte.
\end{enumerate}
\textbf{Réflexe} : présenter l'analyse sous la forme « L'auteur utilise ici un argument de type X : [citation]. Cet argument vise à montrer que... Sa force (ou sa limite) est... ».
\end{methode}

\begin{exemplebox}[Les principaux types d'arguments à connaître]
\begin{itemize}
\item \textbf{Argument d'autorité} : invoquer un penseur, une tradition ou une institution reconnue (« Comme le disait Platon... »).
\item \textbf{Argument par l'exemple} : s'appuyer sur un cas particulier pour illustrer ou généraliser.
\item \textbf{Argument par analogie} : établir une ressemblance entre deux situations (« comme un arbre sans racines... »).
\item \textbf{Argument par l'absurde} : montrer que la thèse adverse conduit à une conséquence inacceptable.
\item \textbf{Argument par distinction} : séparer deux notions confondues (agir / être agi ; devoir / intérêt).
\item \textbf{Argument par élimination} : écarter une à une les solutions possibles pour ne retenir que la dernière.
\end{itemize}
\end{exemplebox}

\begin{exoresolu}[Analyse de l'argumentation — Question 2]
\textbf{Énoncé (type Bac)} : Texte : « Un peuple qui ne connaît pas son histoire est comme un arbre sans racines : le premier vent le renverse. Voyez les nations qui ont oublié leurs langues : elles ont perdu leur identité. Tout peuple doit donc transmettre sa mémoire aux générations futures. »
\begin{enumerate}
\item Identifiez et analysez deux procédés argumentatifs utilisés par l'auteur.
\item Évaluez la solidité de cette argumentation.
\end{enumerate}
\tcblower
\textbf{Solution détaillée.}
\begin{enumerate}
\item \textbf{Procédé 1 — l'analogie} : « comme un arbre sans racines : le premier vent le renverse ». L'auteur compare le peuple sans histoire à un arbre déraciné pour rendre sensible la fragilité d'une communauté privée de mémoire. Cet argument par analogie vise à montrer que l'histoire joue pour un peuple le rôle que les racines jouent pour l'arbre : stabilité et résistance.
\textbf{Procédé 2 — l'argument par l'exemple (induction)} : « Voyez les nations qui ont oublié leurs langues : elles ont perdu leur identité. » L'auteur s'appuie sur des cas observés pour généraliser : perdre sa langue, c'est perdre son identité. Le raisonnement est donc \emph{inductif} : de cas particuliers, il tire une loi générale, résumée dans la conclusion introduite par « donc » : « Tout peuple doit transmettre sa mémoire ».
\item \textbf{Évaluation} : l'analogie est frappante mais ne prouve rien : un peuple n'est pas littéralement un arbre, et une comparaison ne vaut pas démonstration. L'induction, elle, est fragile : quelques exemples de nations ne suffisent pas à établir une loi universelle (il faudrait vérifier qu'aucune nation n'a perdu sa langue sans perdre son identité). L'argumentation est donc \emph{persuasive} (elle touche et convainc) plus que \emph{rigoureusement démonstrative}.
\end{enumerate}
\textbf{Conclusion} : l'auteur combine une analogie et une induction pour établir le devoir de transmission de la mémoire ; l'ensemble est efficace rhétoriquement mais logiquement perfectible.
\end{exoresolu}

\courssub{Fiche méthode 4 : Réussir la Question 3 — la discussion critique}

\begin{methode}[Construire une discussion critique structurée]
\begin{enumerate}
\item \textbf{Problématiser} : transformer la thèse de l'auteur en question : « Faut-il vraiment... ? », « Dans quelle mesure... ? »
\item \textbf{Accorder ce qui doit l'être} (thèse) : montrer en quoi l'auteur a raison, avec un argument ou un exemple personnel (vie courante, réalité camerounaise).
\item \textbf{Objecter} (antithèse) : formuler une limite, une objection ou un contre-exemple, éventuellement en mobilisant un autre philosophe étudié dans l'unité.
\item \textbf{Dépasser} (synthèse) : proposer une position nuancée qui distingue les cas, précise les conditions de validité de la thèse.
\item \textbf{Conclure} par une position personnelle claire et justifiée, en répondant explicitement à la question posée.
\end{enumerate}
\textbf{Formules à retenir} : « La thèse de l'auteur est pertinente dans la mesure où... », « Cependant, on peut objecter que... », « Ainsi, la vérité se situe dans une position nuancée : ... ». \textbf{Piège} : une discussion n'est pas un résumé du texte ni une simple opinion sans argument.
\end{methode}

\begin{exoresolu}[Discussion critique — Question 3]
\textbf{Énoncé (type Bac)} : À propos du texte de Pascal étudié plus haut (« la coutume est une seconde nature »), répondez : « La coutume détruit-elle nécessairement la nature humaine ? » Vous construirez une discussion en trois moments.
\tcblower
\textbf{Solution détaillée.}
\begin{enumerate}
\item \textbf{Thèse — ce que Pascal a de juste} : l'auteur a raison de souligner la puissance de l'habitude. Beaucoup de nos comportements tenus pour « naturels » sont acquis : la langue que nous parlons, nos manières de table, nos préjugés. Exemple : un enfant élevé à Douala dans une famille anglophone trouvera « naturel » de parler anglais, alors que c'est purement coutumier. La coutume recouvre donc bien la nature et se fait passer pour elle.
\item \textbf{Antithèse — la limite de la thèse} : cependant, on peut objecter que tout n'est pas coutume. Il existe en l'homme des dispositions universelles qui résistent aux habitudes : la capacité de raisonner, le sentiment de justice, la peur de la mort. Les droits de l'homme, reconnus à toutes les cultures, témoignent d'une nature commune que les coutumes ne détruisent jamais totalement. De plus, la coutume ne détruit pas toujours : elle peut aussi \emph{accomplir} la nature, comme l'éducation développe l'intelligence naturelle de l'enfant.
\item \textbf{Synthèse — position nuancée} : la vérité se situe entre les deux : la coutume recouvre la nature sans l'abolir. Elle la détruit lorsqu'elle devient conformisme aveugle ; elle la perfectionne lorsqu'elle est choisie et réfléchie. Tout dépend donc de l'usage critique que nous faisons de nos habitudes.
\end{enumerate}
\textbf{Conclusion} : la coutume ne détruit pas nécessairement la nature humaine : elle la menace quand elle est subie, mais elle peut l'épanouir quand l'esprit critique la juge. La thèse de Pascal est donc vraie comme avertissement, fausse comme fatalité.
\end{exoresolu}

\courssub{Fiche méthode 5 : Gérer le temps et rédiger la copie}

\begin{methode}[Gestion stratégique des 3 heures et présentation de la copie]
\begin{enumerate}
\item \textbf{Répartir le temps} : 20 min de lecture et d'analyse du texte ; 40 min pour la Question 1 ; 40 min pour la Question 2 ; 50 min pour la Question 3 ; 30 min de rédaction au propre et relecture.
\item \textbf{Traiter les questions dans l'ordre} : la Question 1 (savoir) prépare la 2 (savoir-faire), qui prépare la 3 (discussion).
\item \textbf{Annoncer chaque question} sur la copie par son intitulé (Question 1, Question 2...) et rédiger en paragraphes distincts, sans plan apparent de dissertation.
\item \textbf{Citer le texte} entre guillemets, avec des citations courtes, et commenter chaque citation (jamais de citation « laissée seule »).
\item \textbf{Relire} en vérifiant : orthographe des noms d'auteurs, cohérence des réponses avec le texte, conclusion présente à la Question 3.
\end{enumerate}
\textbf{Réflexe} : une réponse non rédigée (simples notes, plan) ne rapporte quasiment aucun point : mieux vaut une réponse courte mais rédigée qu'une réponse longue en brouillon.
\end{methode}

\begin{exoresolu}[Simulation Bac : sujet complet corrigé]
\textbf{Énoncé (type Bac)} : Texte : « Le devoir n'est pas ce que les autres attendent de moi, mais ce que je m'impose à moi-même. Obéir par peur de la punition, ce n'est pas agir moralement : c'est calculer. L'homme véritablement moral est celui qui agit par respect de la loi qu'il s'est donnée, fût-ce contre son intérêt. »
\begin{enumerate}
\item Expliquez le sens de ce texte.
\item Analysez l'argumentation de l'auteur.
\item Le devoir moral exclut-il tout calcul d'intérêt ?
\end{enumerate}
\tcblower
\textbf{Solution détaillée.}
\begin{enumerate}
\item \textbf{Question 1 — Sens du texte.} L'auteur soutient que le véritable devoir est une loi que l'on se donne librement à soi-même, et non une contrainte extérieure. La notion de \emph{devoir} désigne ici l'obligation morale intérieure ; celle d'\emph{intérêt} désigne l'avantage personnel, le calcul égoïste. L'articulation du texte : d'abord une définition (« ce que je m'impose à moi-même »), puis une exclusion (« obéir par peur... c'est calculer »), enfin un portrait positif de l'homme moral (« agir par respect de la loi qu'il s'est donnée »). L'enjeu : distinguer la morale authentique (l'autonomie) de la simple obéissance intéressée.
\item \textbf{Question 2 — Argumentation.} L'auteur utilise d'abord une \textbf{distinction conceptuelle} entre devoir authentique (loi de soi à soi) et fausse moralité (obéissance par peur) : « Obéir par peur de la punition, ce n'est pas agir moralement : c'est calculer. » Il procède ensuite par \textbf{opposition} entre deux types d'hommes : celui qui calcule et celui qui respecte la loi morale « fût-ce contre son intérêt » — le « fût-ce » marque une concession qui renforce la thèse. Le raisonnement est déductif : d'une définition du devoir, il tire les critères de l'action morale. L'argumentation est solide car la définition posée au départ commande rigoureusement la conclusion.
\item \textbf{Question 3 — Discussion.} \emph{Thèse} : l'auteur a raison : agir uniquement par intérêt détruit la valeur morale de l'acte. Un commerçant de Yaoundé honnête seulement par peur du contrôle fiscal n'est pas vertueux, il est prudent. \emph{Antithèse} : cependant, on peut objecter que l'intérêt et le devoir ne s'excluent pas toujours : l'honnêteté peut être à la fois un devoir et un avantage (la confiance attire les clients). Exiger que toute action morale soit contraire à l'intérêt, c'est rendre la morale presque impossible. \emph{Synthèse} : ce qui compte n'est pas l'absence d'intérêt, mais le \emph{mobile} principal : une action reste morale si elle aurait été faite même sans avantage. \textbf{Conclusion} : le devoir moral n'exclut pas l'existence d'un intérêt, il exclut que l'intérêt soit la raison d'agir.
\end{enumerate}
\end{exoresolu}

\courssub{Fiche méthode 6 : Transférer l'esprit critique à la vie courante}

\begin{methode}[Appliquer la méthode du texte aux situations concrètes]
\begin{enumerate}
\item \textbf{Identifier la « thèse » implicite} d'un discours courant (rumeur, discours politique, publicité, publication sur les réseaux sociaux) : que veut-on me faire croire ?
\item \textbf{Repérer les arguments} et les procédés : analogie, exemple unique généralisé, appel à la peur ou à la coutume (« on a toujours fait ainsi »).
\item \textbf{Vérifier} : la source est-elle fiable ? L'exemple est-il représentatif ? Y a-t-il une généralisation abusive ?
\item \textbf{Objecter} : chercher au moins un contre-exemple ou une information contradictoire avant d'adhérer.
\item \textbf{Conclure et justifier} : formuler sa position personnelle en la motivant, comme dans une Question 3.
\end{enumerate}
\textbf{Réflexe citoyen} : la méthode « thèse — arguments — discussion » est un outil de lutte contre les rumeurs et la désinformation, par exemple lors des campagnes électorales ou des fausses alertes sanitaires diffusées par téléphone.
\end{methode}

\begin{exoresolu}[Application à une situation de la vie courante]
\textbf{Énoncé} : Sur un réseau social, on lit : « Tous les jeunes qui réussissent ont quitté le village pour la ville. Regardez les exemples autour de vous. Rester au village, c'est rester pauvre. » En appliquant la méthode de l'exercice sur texte : 1) dégagez la thèse de ce discours ; 2) analysez son argumentation ; 3) discutez-la de manière critique.
\tcblower
\textbf{Solution détaillée.}
\begin{enumerate}
\item \textbf{Thèse} : la réussite exige nécessairement de quitter le village pour la ville ; rester au village condamne à la pauvreté.
\item \textbf{Analyse} : le discours utilise une \textbf{induction abusive} : de quelques exemples (« regardez les exemples autour de vous »), il tire une loi universelle (« tous les jeunes »). Il repose aussi sur une \textbf{fausse alternative} (dilemme) : soit la ville et la réussite, soit le village et la pauvreté — comme si aucune autre possibilité n'existait. Aucune preuve statistique n'est fournie.
\item \textbf{Discussion} : \emph{Accord partiel} : il est vrai que la ville offre souvent plus d'infrastructures (universités, hôpitaux, marchés), ce qui peut faciliter certaines réussites. \emph{Objection} : mais le contre-exemple est massif : au Cameroun, de nombreux jeunes réussissent au village grâce à l'agriculture moderne, à l'élevage, à la transformation des produits locaux (manioc, cacao), tandis que l'exode urbain conduit beaucoup de jeunes au chômage et à la précarité dans les quartiers défavorisés de Douala ou de Yaoundé. \emph{Synthèse} : ce n'est pas le lieu qui détermine la réussite, mais la compétence, le travail et les opportunités saisies ; le village comme la ville peuvent être des espaces de réussite ou d'échec.
\end{enumerate}
\textbf{Conclusion} : le discours examiné repose sur une généralisation hâtive et une fausse alternative ; l'esprit critique, formé par l'exercice sur texte, permet de le déconstruire et de formuler un jugement nuancé et justifié.
\end{exoresolu}

\begin{savaistu}[L'exercice sur texte, une tradition de l'école philosophique camerounaise]
Depuis les années 1970, l'enseignement de la philosophie au Cameroun accorde une place centrale à l'exercice sur texte, hérité de la tradition française de l'\emph{explication de texte}. Au Probatoire comme au Baccalauréat technique, l'épreuve de philosophie dure 3 heures et le candidat choisit généralement entre une dissertation et un exercice sur texte. Les correcteurs valorisent trois qualités : la fidélité au texte (Question 1), la rigueur de l'analyse (Question 2) et l'autonomie du jugement (Question 3). S'entraîner régulièrement sur des sujets des sessions antérieures reste la méthode de préparation la plus efficace.
\end{savaistu}

\begin{attention}[Les 5 erreurs qui coûtent des points au Bac]
\begin{enumerate}
\item \textbf{Paraphraser le texte au lieu de l'expliquer} : recopier ou reformuler phrase par phrase sans dégager la thèse ni définir les notions. À la Question 1, il faut dire ce que le texte \emph{signifie}, pas répéter ce qu'il \emph{dit}.
\item \textbf{Donner son avis dès la Question 1} : la critique personnelle n'a sa place qu'à la Question 3. Critiquer l'auteur dans l'explication du sens fait perdre des points sur les deux questions.
\item \textbf{Citer sans commenter (ou commenter sans citer)} : toute citation doit être suivie de son explication, et toute affirmation sur le texte doit être justifiée par une citation courte entre guillemets.
\item \textbf{Confondre les types d'arguments} à la Question 2 : appeler « démonstration » une simple analogie, ou oublier d'évaluer la validité de l'argumentation (l'examinateur attend un jugement : solide, fragile, persuasif...).
\item \textbf{Négliger la gestion du temps et la rédaction} : passer plus d'une heure sur la Question 1, laisser la Question 3 sans conclusion, ou rendre des réponses en simples notes non rédigées. Une conclusion claire et justifiée à la Question 3 est indispensable pour valider la compétence supérieure.
\end{enumerate}
\end{attention}

\begin{exercice}[Entraînement à la maison : sujet complet]
Texte : « La liberté ne consiste pas à faire ce que l'on veut, mais à vouloir ce que l'on peut. Celui qui désire l'impossible se rend esclave de ses désirs ; celui qui règle ses désirs sur ses forces est seul véritablement libre. »
\begin{enumerate}
\item Expliquez le sens de ce texte en définissant la notion de liberté telle que l'auteur l'entend.
\item Analysez l'argumentation : identifiez le procédé principal et évaluez sa solidité.
\item La liberté consiste-t-elle à renoncer à ses désirs ? Vous construirez une discussion en trois moments (thèse, antithèse, synthèse) et vous conclurez par une position personnelle justifiée.
\end{enumerate}
\end{exercice}

\begin{corrige}[Entraînement à la maison : sujet complet]
\begin{enumerate}
\item \textbf{Question 1} : l'auteur soutient que la vraie liberté n'est pas le pouvoir illimité de satisfaire tous ses désirs, mais l'accord entre ce que l'on veut et ce que l'on peut. La \emph{liberté} désigne ici la maîtrise de soi : régler ses désirs sur ses forces réelles. L'articulation : une définition négative puis positive (« ne consiste pas... mais »), suivie d'une opposition entre l'esclave de ses désirs et l'homme libre. L'enjeu : montrer que la liberté est une discipline intérieure et non une absence de limites.
\item \textbf{Question 2} : le procédé principal est la \textbf{distinction conceptuelle} entre deux conceptions de la liberté (faire ce qu'on veut / vouloir ce qu'on peut), complétée par une \textbf{opposition} de portraits (« celui qui désire l'impossible » / « celui qui règle ses désirs »). Le raisonnement est déductif : de la définition de la liberté découlent les deux portraits. L'argumentation est solide car cohérente, mais elle repose sur un présupposé discutable : que l'on peut toujours modifier ses désirs.
\item \textbf{Question 3} : \emph{Thèse} : l'auteur a raison de lier liberté et lucidité : désirer l'impossible (devenir riche sans travailler, réussir sans étudier) conduit à la frustration, qui est une forme d'esclavage. \emph{Antithèse} : cependant, renoncer trop vite à ses désirs, c'est risquer la résignation : beaucoup de progrès (l'indépendance des peuples, l'accès des filles à l'école) sont nés de désirs qui semblaient « impossibles ». \emph{Synthèse} : la liberté n'exige pas de renoncer à ses désirs, mais de les examiner : distinguer ceux qui dépendent de nous (et qu'il faut poursuivre avec méthode) de ceux qui n'en dépendent pas (et qu'il faut transformer ou abandonner). \textbf{Conclusion} : la liberté n'est pas le renoncement, mais le gouvernement réfléchi de ses désirs.
\end{enumerate}
\end{corrige}

\newpage

\coursec{Exercices}

\courssub{Je vérifie mes connaissances}

\begin{exercice}[QCM : Architecture de l'épreuve et distinctions fondamentales]
Pour chaque question, une seule réponse est exacte. Cochez la bonne lettre.
\begin{enumerate}
    \item Dans l'exercice sur texte philosophique au Baccalauréat technique (durée 3h), la répartition temporelle idéale est :
    \begin{itemize}
        \item[A)] 1h pour Q1, 1h pour Q2, 1h pour Q3.
        \item[B)] 45 min pour Q1, 45 min pour Q2, 1h15 pour Q3 (dont 15 min relecture).
        \item[C)] 30 min pour Q1, 1h pour Q2, 1h30 pour Q3.
        \item[D)] 1h pour Q1, 30 min pour Q2, 1h30 pour Q3.
    \end{itemize}
    \item La \cle{Question 1} (Explicitation du sens) exige de :
    \begin{itemize}
        \item[A)] Résumer le texte paragraphe par paragraphe.
        \item[B)] Faire une \emph{paraphrase} en remplaçant les mots difficiles par des synonymes.
        \item[C)] Dégager la \cle{thèse} centrale, le \cle{concept} clé et sa \cle{définition contextuelle}, en montrant comment l'auteur articule ses idées.
        \item[D)] Donner son avis personnel sur le sujet traité.
    \end{itemize}
    \item La \cle{Question 2} (Analyse de l'argumentation) consiste à :
    \begin{itemize}
        \item[A)] Lister tous les exemples cités par l'auteur.
        \item[B)] Identifier la \cle{structure logique} du raisonnement (thèse, arguments, types d'arguments, enchaînements, concessions) sans juger de la vérité du contenu.
        \item[C)] Dire si l'on est d'accord ou pas avec l'auteur.
        \item[D)] Chercher des contre-exemples dans l'actualité camerounaise.
    \end{itemize}
    \item La \cle{Question 3} (Discussion critique) est une \cle{dissertation} à partir du texte. Elle nécessite :
    \begin{itemize}
        \item[A)] Un plan en "Pour / Contre" binaire.
        \item[B)] Une introduction (amorce, rappel du texte, \cle{problématique}, annonce de plan), un développement dialectique (thèses/antithèses/synthèses ou analyse des présupposés), une conclusion (réponse, ouverture).
        \item[C)] Une explication de texte ligne par ligne.
        \item[D)] La citation de 5 auteurs minimum par partie.
    \end{itemize}
\end{enumerate}
\end{exercice}

\begin{exercice}[Vrai / Faux Justifié : Pièges méthodologiques]
Indiquez VRAI ou FAUX et justifiez brièvement en une phrase (référence à la méthode).
\begin{enumerate}
    \item "Pour réussir la Q1, il faut reprendre l'ordre chronologique du texte et tout expliquer." \hfill \textbf{...}
    \item "Un \cle{argument d'autorité} (citation d'un savant) a plus de valeur logique qu'un \cle{argument de fait} (donnée chiffrée)." \hfill \textbf{...}
    \item "La \cle{concession} ('Certes... mais...') affaiblit la thèse de l'auteur." \hfill \textbf{...}
    \item "En Q3, il est interdit d'utiliser des exemples qui ne sont pas dans le texte." \hfill \textbf{...}
    \item "La \cle{problématique} de la Q3 doit être une question fermée (réponse Oui/Non)." \hfill \textbf{...}
    \item "Le \cle{concept} central du texte est toujours un mot en gras ou en italique dans l'extrait." \hfill \textbf{...}
    \item "L'\cle{explicitation} (Q1) et l'\cle{analyse} (Q2) portent sur le même objet : le texte." \hfill \textbf{...}
    \item "Un \cle{plan dialectique} pour la Q3 impose obligatoirement trois parties : Thèse, Antithèse, Synthèse." \hfill \textbf{...}
\end{enumerate}
\end{exercice}

\begin{exercice}[Définitions opérationnelles : Le vocabulaire de l'argumentation]
Donnez une définition précise et rigoureuse (2 à 3 lignes max) pour chacun des termes suivants, en précisant son rôle dans l'exercice sur texte.
\begin{enumerate}
    \item \cle{Thèse} (dans un texte philosophique)
    \item \cle{Concept} (philosophique)
    \item \cle{Argument} (vs Exemple)
    \item \cle{Enchaînement déductif} (vs Inductif / Analogique)
    \item \cle{Problématique} (vs Sujet / Question de cours)
    \item \cle{Présupposé} (ontologique, épistémologique, axiologique)
    \item \cle{Critique immanente} (vs Critique transcendant / extériorité)
    \item \cle{Transfert} (compétence visée en Terminale Tech)
\end{enumerate}
\end{exercice}

\begin{exercice}[Analyse structurelle rapide : Identification de la forme logique]
Pour chacun des courts raisonnements ci-dessous, identifiez : (a) la thèse défendue, (b) le(s) argument(s) principal(aux), (c) le type de liaison logique (déduction, induction, analogie, concession, réduction à l'absurde, argument d'autorité).
\begin{enumerate}
    \item "Tout être humain est mortel. Socrate est un homme. Donc Socrate est mortel."
    \item "Le barrage de Nachtigal produit de l'électricité sans émettre de CO2. Le barrage de Memve'ele aussi. Donc l'hydroélectricité au Cameroun est une énergie 'propre'."
    \item "Aristote affirme que 'l'homme est un animal politique'. Donc la vie en société est naturelle à l'homme."
    \item "Certes, l'IA améliore le diagnostic médical. Mais elle déshumanise la relation soignant-soigné. Donc son intégration doit être encadrée éthiquement."
    \item "Si la technique nous libère de la peine, elle nous asservit à la machine. Or nous voulons être libres. Donc il faut limiter la technique." (Schéma simplifié d'Ellul/Jonas)
\end{enumerate}
\end{exercice}

\courssub{Je m'entraîne}

\begin{exercice}[Entraînement ciblé Q1 : Explicitation de concepts sur extraits courts]
\emph{Consigne :} Pour chacun des 5 extraits ci-dessous, rédigez en \textbf{10 lignes maximum} : (1) Le \cle{concept central} (mot ou expression), (2) Sa \cle{définition contextuelle} (ce que l'auteur en fait dans cet extrait précis), (3) La \cle{thèse} défendue en une phrase. \emph{Travail en autonomie chronométré (8 min par extrait), puis correction par échange de copies avec grille d'évaluation.}

\begin{enumerate}
    \item \textbf{Extrait A (Aristote, \emph{Éthique à Nicomaque}, I, 2) :} "Puisque toutes les connaissances et tous les choix visent quelque bien, quel est celui que nous voulons pour la politique ? C'est le bien suprême, qui est visé pour lui-même... C'est le bonheur."
    \item \textbf{Extrait B (Descartes, \emph{Discours de la méthode}, 6e partie) :} "Il est possible de parvenir à des connaissances fort utiles à la vie... pour nous rendre comme \cle{maîtres et possesseurs de la nature}."
    \item \textbf{Extrait C (Hannah Arendt, \emph{Condition de l'homme moderne}) :} "Le travail est l'activité qui correspond au processus biologique du corps humain... Le travail ne produit rien de durable... L'œuvre, au contraire, édifie un monde d'objets durables."
    \item \textbf{Extrait D (Kwasi Wiredu, \emph{Philosophie et culture africaine}) :} "La vérité n'est pas une correspondance avec une réalité indépendante de l'esprit, mais une \cle{cohérence} au sein d'un système de croyances partagées par la communauté."
    \item \textbf{Extrait E (Achille Mbembe, \emph{Critique de la raison nègre}) :} "La \cle{colonie} n'est pas seulement un territoire conquis. C'est un dispositif de \cle{necropouvoir} : le pouvoir de dicter qui peut vivre et qui doit mourir."
\end{enumerate}
\end{exercice}

\begin{exercice}[Entraînement ciblé Q2 : Représentation schématique de l'argumentation (Arbre argumentatif)]
\emph{Consigne :} Lisez le texte ci-dessous. Représentez le raisonnement de l'auteur sous forme d'un \textbf{arbre argumentatif} (schéma arborescent : Thèse $\rightarrow$ Arguments $\rightarrow$ Sous-arguments/Exemples). Légendez \textbf{chaque lien logique} (ex: \textit{Argument de fait}, \textit{Analogie}, \textit{Concession}, \textit{Enchaînement causal}, \textit{Réfutation d'objection}). Temps : 25 min.

\textbf{Texte :} "On nous promet une \cle{transition écologique} par la technologie verte : voitures électriques, panneaux solaires, éoliennes. \textbf{Certes}, ces innovations réduisent les émissions locales de CO2 (Argument 1 : \textit{Fait technique}). \textbf{Mais} elles déplacent la pollution : l'extraction du lithium, du cobalt, des terres rares pour les batteries détruit des écosystèmes entiers en Afrique (RDC, Zambie) et en Amérique du Sud (Argument 2 : \textit{Preuve par les faits / Transfert de charge}). \textbf{De plus}, l'obsolescence programmée de ces 'objets verts' génère des montagnes de déchets électroniques (D3EE) que l'Europe exporte illégalement vers nos ports (Douala, Abidjan) sous couvert de 'don' (Argument 3 : \textit{Argument d'injustice / Expérience vécue}). \textbf{Par analogie}, c'est comme soigner un malade en lui transférant sa tumeur sur un autre organe : le mal n'est pas guéri, il est déplacé (Argument 4 : \textit{Analogie}). \textbf{Donc}, sans remise en cause du \cle{modèle productiviste} et de la surconsommation, la 'technologie verte' n'est qu'un \cle{verdissement} (greenwashing) du capitalisme extractiviste (Thèse finale)."
\end{exercice}

\begin{exercice}[Entraînement ciblé Q3 : Introduction et Plan détaillé (Chrono 25 min)]
\emph{Sujet :} "\textbf{La technique libère-t-elle l'homme de la nécessité ?}"
\emph{Consigne :} Rédigez \textbf{uniquement} :
\begin{enumerate}
    \item Une \textbf{introduction complète} (Amorce \cle{accroche} $\rightarrow$ Rappel du sujet / Définition des termes $\rightarrow$ \cle{Problématique} $\rightarrow$ Annonce de \cle{plan}).
    \item Un \textbf{plan détaillé} (3 parties majeures, 2 à 3 sous-parties chacune) avec pour chaque sous-partie : \textbf{L'argument central}, \textbf{Un auteur/référence philosophique} (programme Terminale Tech), \textbf{Un exemple concret camerounais (CM)}.
\end{enumerate}
\emph{Ne rédigez pas le développement.}
\end{exercice}

\begin{exercice}[Atelier "Pièges \& Corrections" : Analyse de copies anonymisées (Binôme)]
\emph{Consigne :} Voici 4 extraits de copies d'élèves (Session antérieure). Chaque copie contient \textbf{3 erreurs types majeures}. En binôme : (1) Identifiez les 3 erreurs par copie, (2) Citez la règle méthodologique violée, (3) Proposez une correction rédactionnelle pour l'erreur la plus grave, (4) Attribuez une note sur 5 pour la question concernée (grille officielle : 0=HS, 1=Insuffisant, 2=Passable, 3=Assez bien, 4=Bien, 5=Très bien).

\begin{enumerate}
    \item \textbf{Copie 1 (Q1 - Extrait sur la technique chez Heidegger) :} "Heidegger dit que la technique n'est pas un outil. C'est un destin. L'auteur explique que l'essence de la technique est la \cle{Gestell} (mise en réserve). Il donne l'exemple de la centrale hydroélectrique sur le Rhin. Ensuite il parle du danger. Puis il dit que l'art peut nous sauver. En conclusion, la technique est dangereuse." \hfill (Note Q1 : /5)
    \item \textbf{Copie 2 (Q2 - Texte sur la justice climatique) :} "L'auteur est pour la justice climatique. Son premier argument c'est que les pays riches polluent plus. Son deuxième argument c'est que les pays pauvres souffrent plus. Il cite le cyclone au Mozambique. Il fait une concession : 'certes les pays pauvres polluent aussi'. Mais il réfute. Donc sa thèse est juste." \hfill (Note Q2 : /5)
    \item \textbf{Copie 3 (Q3 - Sujet : 'L'État doit-il réguler l'IA ?') :} "Pour : L'État doit réguler car l'IA est dangereuse (deepfakes, biais). Exemple : élection USA 2024. Contre : L'État ne doit pas réguler car ça freine l'innovation. Exemple : Silicon Valley. Synthèse : Il faut une régulation équilibrée." \hfill (Note Q3 : /5)
    \item \textbf{Copie 4 (Q1/Q2 confusion - Extrait d'Arendt sur travail/œuvre) :} "Arendt pense que le travail c'est pour manger et l'œuvre c'est pour durer. Je suis d'accord car au Cameroun les maçons font des maisons qui durent. Mais les fonctionnaires font du travail qui ne dure pas. L'auteur a raison de séparer les deux. C'est une bonne analyse." \hfill (Note Q1/Q2 : /5)
\end{enumerate}
\end{exercice}

\courssub{Je me prépare au Bac}

\begin{exercice}[Sujet Type Bac Tech 2024 : Simulation complète (3h)]
\emph{Durée : 3 heures. Coefficient selon série. Le texte fait environ 15 lignes.}

\textbf{TEXTE : Achim Steiner (Directeur PNUE), \emph{Discours d'ouverture COP28 / Rapport 'Emissions Gap' 2023} (adapté).}
"La transition vers une économie verte n'est pas une option technique, c'est un impératif de \cle{justice climatique}. Les technologies vertes — solaire, éolien, batteries, hydrogène vert — offrent les moyens de découpler la prospérité de la destruction écologique. Pourtant, le déploiement de ces technologies reproduit trop souvent les inégalités du passé. L'extraction des minéraux critiques (cobalt, lithium, cuivre) pour les batteries et les éoliennes concentre les profits dans les mains de multinationales du Nord, tandis que les coûts environnementaux et sanitaires (pollution de l'eau, travail des enfants, déplacements de populations) pèsent sur les communautés du Sud global, du Katanga au Salar d'Atacama. Une transition qui ignore les droits des peuples autochtones, qui exporte ses déchets électroniques vers le Golfe de Guinée, ou qui endette les pays africains pour financer des parcs solaires dont l'énergie est exportée vers l'Europe, n'est pas une transition : c'est un \cle{néo-extractivisme} vert. La \cle{justice climatique} exige donc une \cle{gouvernance mondiale} des ressources minérales, un transfert de technologie \emph{inconditionnel}, et la reconnaissance de la \cle{dette écologique} du Nord envers le Sud. Seule une transition \emph{juste} peut être durable."

\textbf{QUESTIONS :}
\begin{enumerate}
    \item \textbf{Explicitez le sens de ce texte} en dégagent la thèse centrale, le ou les concepts clés et leur définition contextuelle, ainsi que l'enchaînement des idées. (6 points)
    \item \textbf{Analysez l'argumentation} de l'auteur : identifiez la structure du raisonnement, les types d'arguments utilisés, les éventuelles concessions et leur fonction, la portée de la conclusion. (7 points)
    \item \textbf{La transition écologique peut-elle être juste sans remise en cause du modèle de développement dominant ?} Discutez en vous appuyant sur le texte, vos connaissances philosophiques (auteurs, concepts du programme) et des exemples concrets camerounais (ex: barrages de Lom Pangar, Nachtigal, Memve'ele ; mines de cobalt/terres rares potentielles ; agriculture industrielle vs paysanne ; gestion des D3EE à Douala). (7 points)
\end{enumerate}
\end{exercice}

\begin{exercice}[Sujet Bac Type : Philosophie de la technique \& Éthique professionnelle (Hans Jonas / Günther Anders)]
\emph{Durée : 3 heures.}

\textbf{TEXTE : Hans Jonas, \emph{Le Principe responsabilité} (1979), Chap. 'Vers une éthique pour l'ère technologique' (adapté).}
"L'éthique traditionnelle était une éthique du prochain : l'agir humain avait une portée limitée dans l'espace et le temps. La \cle{technique} moderne a brisé ce cercle. La nature de l'action a changé : sa puissance est devenue incommensurable, ses effets sont planétaires, cumulatifs, irréversibles et souvent différés dans le temps (effet de serre, déchets nucléaires, manipulation génétique). L'homme n'a plus affaire à une nature qui 'pardonne' (régénération), mais à une nature \cle{vulnérable} que sa puissance peut détruire définitivement. Il n'y a plus de 'prochain' au sens classique, car les générations futures, qui n'existent pas encore, sont les premières concernées par nos actes présents. Il faut donc un \cle{nouvel impératif catégorique} : 'Agis de façon que les effets de ton action soient compatibles avec la permanence d'une vie authentiquement humaine sur terre.' La \cle{responsabilité} devient le fondement de l'éthique, non plus la liberté abstraite. Elle s'étend à la biosphère tout entière."

\textbf{QUESTIONS :}
\begin{enumerate}
    \item \textbf{Explicitez le sens de ce texte} en montrant comment Jonas caractérise la rupture introduite par la technique moderne et quel nouveau fondement éthique il propose. (6 points)
    \item \textbf{Analysez l'argumentation} : comment Jonas passe-t-il de l'analyse ontologique de la technique (puissance, irréversibilité) à la prescription normative (impératif de responsabilité) ? Identifiez les prémisses, l'analogie avec Kant, le type de finalisme. (7 points)
    \item \textbf{La responsabilité de l'ingénieur/technicien suffit-elle à garantir un usage éthique de la technique ?} Discutez en mobilisant Jonas, Anders (\cle{prométhéen} / \cle{honte prométhéenne}), le concept de \cle{technoscience}, et des exemples du génie civil, de l'agro-industrie ou du numérique au Cameroun (ex: gestion données biométriques, OGM, barrages). (7 points)
\end{enumerate}
\end{exercice}

\begin{exercice}[Préparation Grand Oral Transversal : "Philosophie \& Dilemme Éthique Professionnel"]
\emph{Consigne :} En vue de l'épreuve orale (Grand Oral / Soutenance de projet technique), préparez une fiche de présentation de \textbf{5 minutes maximum}.
\emph{Choisissez UN dilemme éthique lié à VOTRE future spécialité technique (Génie Civil, Électrotechnique, Informatique/Réseaux, Agro-alimentaire, Maintenance, Topographie, etc.).}

\textbf{Structure imposée de la fiche (à rendre écrite + soutenance orale) :}
\begin{enumerate}
    \item \textbf{Accroche (30s) :} Situation concrète, chiffre choc, anecdote de stage/terrain au Cameroun.
    \item \textbf{Texte philosophique ancrage (1 min) :} Citez \textbf{un extrait court (3-4 lignes)} d'un auteur au programme (Aristote, Descartes, Kant, Arendt, Jonas, Heidegger, Wiredu, Mbembe, Foucault, etc.) qui éclaire le dilemme. Expliquez le concept clé.
    \item \textbf{Problématique du dilemme (30s) :} Formulez la tension éthique (ex: Efficacité vs Équité ; Sécurité vs Vie privée ; Productivité vs Biodiversité ; Innovation vs Précaution).
    \item \textbf{Analyse philosophique structurée (2 min) :} 2 arguments philosophiques opposés (Thèse / Antithèse) appuyés sur les auteurs + 1 tentative de dépassement (Synthèse / Principe de responsabilité / Éthique de la vertu / Justice distributive).
    \item \textbf{Positionnement professionnel \& Résolution concrète (1 min) :} En tant que futur technicien supérieur camerounais, quelle \cle{conduite responsable} proposez-vous ? (Référence : Code de déontologie de votre corps de métier, Normes ISO, Loi cadre environnement Cameroun, ODD ONU).
    \item \textbf{Ouverture (30s) :} Question restée ouverte / Enjeu pour la formation continue / Lien avec la citoyenneté active.
\end{enumerate}

\textbf{Exemples de dilemmes par filière (non exhaustifs) :}
\begin{itemize}
    \item \textit{Génie Civil / Travaux Publics :} Déplacement populations / Expropriation pour barrage/route vs Droit au logement / Patrimoine culturel (Ex: Barrage Lom Pangar, Autoroute Yaoundé-Douala).
    \item \textit{Électrotechnique / Énergies Renouvelables :} Accès électricité rurale (Solaire) vs Déchets batteries / Dépendance import / Accaparement terres.
    \item \textit{Informatique / Data Science / IA :} Diagnostic médical IA / Prédiction délinquance vs Biais algorithmiques / Consentement données / Responsabilité décisionnelle (Boîte noire).
    \item \textit{Agro-alimentaire / Industries Agricoles :} OGM / Intrants chimiques pour rendement vs Souveraineté alimentaire / Santé sols / Semences paysannes.
    \item \textit{Topographie / Géomatique :} Cartographie foncière / Données GPS vs Conflits fonciers coutumiers / Surveillance populations / Vente données géolocalisées.
\end{itemize}
\emph{Évaluation : Grille officielle Grand Oral (Prise de parole, Connaissances, Argumentation, Interaction, Projet professionnel).}
\end{exercice}

\newpage

\coursec{Corrigés des exercices}

\begin{corrige}[Exercice 1 — QCM : Architecture de l'épreuve et distinctions fondamentales]
\textbf{Question 1 : réponse B.} La répartition 45 min (Q1) / 45 min (Q2) / 1h15 (Q3, relecture incluse) est la seule cohérente avec la hiérarchie des points (Q1 : 6 pts, Q2 : 7 pts, Q3 : 7 pts) et avec la nécessité de relire. Les répartitions A, C et D sous-évaluent la Q3, qui exige une dissertation complète.

\textbf{Question 2 : réponse C.} L'explicitation n'est ni un résumé linéaire (A), ni une paraphrase (B), ni un avis personnel (D) : elle dégage la \cle{thèse}, le \cle{concept} clé, sa \cle{définition contextuelle} et l'articulation des idées.

\textbf{Question 3 : réponse B.} L'analyse porte sur la \cle{structure logique} (arguments, liaisons, concessions), sans jugement de vérité : c'est ce qui la distingue de la Q3 (C et D relèvent de la discussion ; A est un simple relevé).

\textbf{Question 4 : réponse B.} La Q3 est une dissertation : introduction problématisée, développement dialectique, conclusion avec ouverture. Le plan binaire « Pour/Contre » sans synthèse (A) est insuffisant ; C confond avec l'explication de texte ; D transforme l'érudition en catalogue.
\end{corrige}

\begin{attention}[Points de vigilance]
Ne jamais confondre \emph{reformuler} (paraphrase, sanctionnée) et \emph{expliciter} (dégager l'articulation thèse-concept-arguments). Le correcteur repère immédiatement la copie qui « suit le texte ligne par ligne ».
\end{attention}

\begin{corrige}[Exercice 2 — Vrai / Faux Justifié : Pièges méthodologiques]
\begin{enumerate}
    \item \textbf{FAUX.} La Q1 exige de dégager la \emph{structure} du texte (thèse, concepts, articulation), non de suivre son ordre chronologique : on réorganise selon la logique, pas selon la chronologie.
    \item \textbf{FAUX.} L'argument d'autorité est logiquement le plus faible : la vérité d'une thèse ne dépend pas du prestige de qui l'énonce ; un argument de fait vérifiable a une portée probante supérieure.
    \item \textbf{FAUX.} La concession (« Certes\dots mais\dots ») \emph{renforce} la thèse : en intégrant et en réfutant l'objection adverse, l'auteur montre qu'il l'a anticipée.
    \item \textbf{FAUX.} La Q3 exige au contraire des exemples personnels et actuels (contexte camerounais compris) pour discuter la thèse : seuls les exemples hors texte manifestent l'autonomie de pensée.
    \item \textbf{FAUX.} La problématique doit être une question \emph{ouverte} (« Dans quelle mesure\dots ? », « Peut-on\dots ? ») qui appelle une discussion, non une réponse binaire Oui/Non.
    \item \textbf{FAUX.} Le concept central se repère par sa \emph{fonction} dans le raisonnement (ce dont tout le texte parle), non par sa typographie ; la mise en gras est un artifice d'édition.
    \item \textbf{VRAI.} Q1 et Q2 portent toutes deux sur le texte, mais sous des angles différents : Q1 sur le \emph{sens} (ce que dit l'auteur), Q2 sur la \emph{forme argumentative} (comment il le prouve).
    \item \textbf{FAUX.} Le plan dialectique peut aussi procéder par analyse des présupposés ou par dépassement progressif ; le schéma Thèse/Antithèse/Synthèse est un modèle, pas une obligation.
\end{enumerate}
\end{corrige}

\begin{attention}[Point de vigilance]
La distinction Q1/Q2 est cruciale : Q1 = \emph{quoi} (contenu, sens, thèse, concepts) ; Q2 = \emph{comment} (structure, arguments, liaisons logiques, portée). Les confondre coûte cher.
\end{attention}

\begin{corrige}[Exercice 3 — Définitions opérationnelles : Le vocabulaire de l'argumentation]
\begin{enumerate}
    \item \textbf{Thèse :} position doctrinale que l'auteur défend et cherche à établir comme vraie. Dans l'exercice, c'est le premier élément à dégager en Q1, en une phrase complète.
    \item \textbf{Concept :} idée générale et abstraite (liberté, justice, technique) dont le texte propose une définition contextuelle. Son identification commande toute l'explicitation.
    \item \textbf{Argument :} raison avancée pour fonder la thèse ; l'\emph{exemplebox} n'est qu'un cas particulier qui illustre ou étaye l'argument, sans valeur démonstrative propre.
    \item \textbf{Enchaînement déductif :} passage nécessaire du général au particulier (syllogisme) ; l'\emph{induction} va du particulier au général (portée probable) ; l'\emph{analogie} conclut par ressemblance (portée faible).
    \item \textbf{Problématique :} question philosophique ouverte née d'une tension entre deux thèses ; elle se distingue du sujet (simple thème) et de la question de cours (réponse connue d'avance).
    \item \textbf{Présupposé :} ce qu'un discours admet sans le démontrer : ontologique (sur l'être), épistémologique (sur la connaissance) ou axiologique (sur les valeurs). Le repérer est le cœur de la critique.
    \item \textbf{Critique immanente :} réfutation interne montrant que le texte se contredit ou ne prouve pas ce qu'il prétend ; la critique extérieure oppose une thèse rivale. L'immanente est philosophiquement plus forte.
    \item \textbf{Transfert :} capacité à mobiliser une notion du cours dans une situation concrète nouvelle (vie professionnelle, citoyenneté) ; compétence terminale visée par la Q3 et le Grand Oral.
\end{enumerate}
\end{corrige}

\begin{methode}[Grille d'auto-évaluation des définitions]
1 point : définition précise et philosophique (pas de sens commun) ; 1 point : distinction claire avec les notions voisines ; 1 point : exemple d'application au texte philosophique. Total /3 par notion.
\end{methode}

\begin{corrige}[Exercice 4 — Analyse structurelle rapide : Identification de la forme logique]
\begin{enumerate}
    \item \textbf{Thèse :} Socrate est mortel. \textbf{Arguments :} prémisse majeure (tout homme est mortel) et mineure (Socrate est un homme). \textbf{Liaison :} \cle{déduction} (syllogisme aristotélicien, conclusion nécessaire).
    \item \textbf{Thèse :} l'hydroélectricité camerounaise est une énergie propre. \textbf{Arguments :} deux cas (Nachtigal, Memve'ele). \textbf{Liaison :} \cle{induction} (généralisation à partir de cas particuliers ; portée probable, réfutable par un contre-exemple).
    \item \textbf{Thèse :} la vie en société est naturelle à l'homme. \textbf{Argument :} l'autorité d'Aristote. \textbf{Liaison :} \cle{argument d'autorité} (faible : la citation ne prouve pas, elle appuie).
    \item \textbf{Thèse :} l'intégration de l'IA doit être encadrée éthiquement. \textbf{Arguments :} bénéfice médical concédé, puis risque de déshumanisation. \textbf{Liaison :} \cle{concession} suivie d'une réfutation (« Certes\dots Mais\dots »).
    \item \textbf{Thèse :} il faut limiter la technique. \textbf{Arguments :} la technique asservit ; or nous voulons être libres. \textbf{Liaison :} \cle{déduction} normative adossée à un \emph{raisonnement par l'absurde} implicite (accepter la technique sans limite contredit notre volonté de liberté).
\end{enumerate}
\end{corrige}

\begin{attention}[Point de vigilance]
Toujours nommer le type de liaison \emph{et} évaluer sa force : déduction (nécessaire), induction (probable), analogie et autorité (faibles). C'est cette évaluation qui fait la différence entre une copie « passable » et « bien ».
\end{attention}

\begin{corrige}[Exercice 5 — Entraînement ciblé Q1 : Explicitation de concepts sur extraits courts]
\textbf{Extrait A (Aristote).} \emph{Concept central :} le \cle{bien suprême} (le bonheur). \emph{Définition contextuelle :} ce qui est visé pour lui-même et non en vue d'autre chose, fin dernière de l'action politique. \emph{Thèse :} toute action humaine étant ordonnée à une fin, la politique a pour objet le bien suprême, c'est-à-dire le bonheur.

\textbf{Extrait B (Descartes).} \emph{Concept central :} la \cle{maîtrise de la nature}. \emph{Définition contextuelle :} puissance pratique issue de la connaissance, rendant l'homme « comme maître et possesseur de la nature ». \emph{Thèse :} la science nouvelle a une finalité pratique : rendre la vie humaine plus commode en dominant la nature.

\textbf{Extrait C (Arendt).} \emph{Concept central :} la distinction \cle{travail}/\cle{œuvre}. \emph{Définition contextuelle :} le travail répond au cycle biologique (consommable, non durable) ; l'œuvre fabrique des objets durables qui constituent un monde commun. \emph{Thèse :} seule l'œuvre, et non le travail, fonde la permanence du monde humain.

\textbf{Extrait D (Wiredu).} \emph{Concept central :} la \cle{vérité comme cohérence}. \emph{Définition contextuelle :} non pas adéquation à une réalité indépendante, mais accord interne d'un système de croyances partagé par une communauté. \emph{Thèse :} la vérité est relative à un consensus culturel et non à une correspondance avec les choses.

\textbf{Extrait E (Mbembe).} \emph{Concept central :} le \cle{nécropouvoir}. \emph{Définition contextuelle :} pouvoir de décider qui peut vivre et qui doit mourir, essence de la colonie au-delà de la conquête territoriale. \emph{Thèse :} la colonie est un dispositif de mort autant que d'occupation.
\end{corrige}

\begin{methode}[Grille d'auto-évaluation par échange de copies]
1 point : concept correctement identifié ; 1 point : définition \emph{contextuelle} (pas de définition de dictionnaire) ; 1 point : thèse formulée en une phrase complète ; 1 point : respect des 10 lignes et clarté rédactionnelle. Total /4 par extrait.
\end{methode}

\begin{corrige}[Exercice 6 — Entraînement ciblé Q2 : Arbre argumentatif]
\textbf{Thèse finale (racine) :} sans remise en cause du modèle productiviste, la « technologie verte » n'est qu'un \cle{greenwashing} du capitalisme extractiviste.

\begin{popfigure}
\begin{tikzpicture}[
  grow=down,
  level distance=1.6cm,
  level 1/.style={sibling distance=4.6cm},
  level 2/.style={sibling distance=4.2cm},
  every node/.style={draw=popblue, rounded corners, fill=popblueL, text width=3.6cm, align=center, font=\small},
  edge from parent/.style={draw=popdark, -{Stealth}}
]
\node[fill=poporangeL, draw=poporange, text width=6.5cm]{\textbf{Thèse} : la technologie verte = greenwashing sans changement de modèle}
  child { node[align=center]{A1 : réduction locale du CO2 \\ \emph{(concession / fait technique)}} }
  child { node[align=center]{A2 : déplacement de la pollution (lithium, cobalt — RDC, Zambie) \\ \emph{(preuve par les faits)}} 
    child { node[align=center]{A4 : soigner en déplaçant la tumeur \\ \emph{(analogie)}} } }
  child { node[align=center]{A3 : D3EE exportés vers Douala, Abidjan \\ \emph{(argument d'injustice / expérience)}} };
\end{tikzpicture}
\legende{Arbre argumentatif du texte sur la transition écologique : la concession A1 est intégrée puis dépassée par A2 et A3 ; l'analogie A4 illustre A2.}
\end{popfigure}

\textbf{Légende des liens logiques :}
\begin{itemize}
    \item Thèse $\leftarrow$ A1 : \cle{concession} (« Certes\dots ») immédiatement réfutée (« Mais\dots ») — elle renforce la thèse en anticipant l'objection.
    \item Thèse $\leftarrow$ A2 : \cle{enchaînement causal} et preuve par les faits (le « vert » déplace la pollution au Sud).
    \item Thèse $\leftarrow$ A3 : argument d'\cle{injustice} adossé à l'expérience vécue (export illégal des déchets).
    \item A2 $\leftarrow$ A4 : \cle{analogie} médicale qui illustre le mécanisme de déplacement du mal.
\end{itemize}
\end{corrige}

\begin{attention}[Point de vigilance]
Dans un arbre, chaque flèche doit être \emph{étiquetée} par son type logique. Un schéma sans légende des liens perd la moitié des points de la Q2.
\end{attention}

\begin{corrige}[Exercice 7 — Entraînement ciblé Q3 : « La technique libère-t-elle l'homme de la nécessité ? »]
\textbf{Introduction (modèle attendu).} \emph{Amorce :} du smartphone au barrage hydroélectrique, jamais l'homme n'a disposé d'autant de moyens pour dompter la nature. \emph{Rappel et définitions :} la \cle{technique} est l'ensemble des procédés et outils par lesquels l'homme transforme son milieu ; la \cle{nécessité} désigne ce à quoi l'homme est contraint (besoins vitaux, lois naturelles). \emph{Problématique :} si la technique semble nous affranchir du travail pénible et des contraintes naturelles, ne crée-t-elle pas de nouvelles dépendances ? Dans quelle mesure libère-t-elle réellement de la nécessité ? \emph{Annonce de plan :} nous verrons d'abord que la technique libère des contraintes naturelles, puis qu'elle engendre de nouvelles servitudes, enfin qu'elle exige un usage responsable pour être émancipatrice.

\textbf{Plan détaillé (corrigé type).}
\begin{enumerate}
    \item \textbf{La technique affranchit de la nécessité naturelle.}
    \begin{itemize}
        \item 1.1. Argument : l'outil prolonge le corps et soulage la peine. Référence : Descartes, « maîtres et possesseurs de la nature ». Exemple CM : l'irrigation et les moulins à maïs allègent le travail paysan à l'Adamaoua.
        \item 1.2. Argument : la médecine technique recule la maladie et la mort. Référence : Bacon, « la science pour soulager la condition humaine ». Exemple CM : campagnes de vaccination et maternités équipées à Yaoundé.
    \end{itemize}
    \item \textbf{Mais la technique engendre de nouvelles servitudes.}
    \begin{itemize}
        \item 2.1. Argument : la machine crée la dépendance et l'aliénation du travailleur. Référence : Marx, l'ouvrier asservi à la machine ; Ellul, la « société technicienne ». Exemple CM : dépendance aux pièces détachées importées pour les engins de BTP.
        \item 2.2. Argument : la puissance technique menace la nature et l'homme lui-même. Référence : Jonas, \emph{Le Principe responsabilité}. Exemple CM : pollution des nappes phréatiques autour des zones industrielles de Douala.
    \end{itemize}
    \item \textbf{La libération exige un usage maîtrisé et responsable de la technique.}
    \begin{itemize}
        \item 3.1. Argument : la technique est un moyen dont la valeur dépend de la fin. Référence : Aristote, la \emph{phronèsis} (prudence) ; Heidegger, la technique comme mode de dévoilement à interroger. Exemple CM : choix de matériaux locaux (briques de terre comprimée) plutôt que tout-importé.
        \item 3.2. Argument : l'éthique de la responsabilité doit guider le technicien. Référence : Jonas, impératif de responsabilité envers les générations futures. Exemple CM : dimensionnement durable des barrages de Nachtigal et Lom Pangar.
    \end{itemize}
\end{enumerate}
\end{corrige}

\begin{attention}[Barème indicatif Q3 (/7)]
Introduction problématisée : 2 pts ; plan dialectique cohérent : 2 pts ; références philosophiques exactes : 1,5 pt ; exemples camerounais pertinents : 1 pt ; langue et présentation : 0,5 pt.
\end{attention}

\begin{corrige}[Exercice 8 — Atelier « Pièges \& Corrections » : Analyse de copies anonymisées]
\textbf{Copie 1 (Q1, Heidegger) — Note : 1/5.}
\begin{itemize}
    \item Erreur 1 : \emph{paraphrase chronologique} (« Ensuite\dots Puis\dots ») — règle violée : la Q1 dégage la structure, pas l'ordre du texte.
    \item Erreur 2 : absence de \emph{définition contextuelle} de la \cle{Gestell} — le concept est cité, jamais expliqué.
    \item Erreur 3 : thèse réduite à un slogan (« la technique est dangereuse ») qui trahit Heidegger (le danger n'est pas la technique mais son essence comme arraisonnement).
    \item Correction de l'erreur la plus grave : « Pour Heidegger, l'essence de la technique n'est rien de technique : elle est \cle{Gestell}, mode de dévoilement qui somme la nature de se livrer comme fonds exploitable — le Rhin n'est plus un fleuve mais une réserve d'énergie. »
\end{itemize}

\textbf{Copie 2 (Q2, justice climatique) — Note : 2/5.}
\begin{itemize}
    \item Erreur 1 : simple \emph{liste} d'arguments sans typologie (fait ? causalité ? analogie ?) — la Q2 exige de nommer les types d'arguments.
    \item Erreur 2 : la concession est signalée mais sa \emph{fonction} (anticiper l'objection) n'est pas analysée.
    \item Erreur 3 : jugement de vérité (« sa thèse est juste ») — hors de propos en Q2, qui analyse la forme, pas la vérité.
    \item Correction : supprimer le jugement final et écrire : « La concession ("certes les pays pauvres polluent aussi") est immédiatement réfutée, ce qui renforce la thèse en neutralisant l'objection adverse. »
\end{itemize}

\textbf{Copie 3 (Q3, régulation de l'IA) — Note : 1/5.}
\begin{itemize}
    \item Erreur 1 : absence totale d'\emph{introduction} (pas de problématique, pas de définitions).
    \item Erreur 2 : plan « Pour/Contre » binaire avec synthèse artificielle non argumentée (« équilibrée » = lieu commun).
    \item Erreur 3 : aucune référence philosophique (ni Jonas, ni Kant, ni aucun concept du programme).
    \item Correction de la synthèse : « La régulation n'est pas un frein mais une condition de l'innovation responsable : selon Jonas, la puissance technique appelle un impératif de responsabilité ; l'État ne bride pas l'IA, il la rend compatible avec la dignité humaine. »
\end{itemize}

\textbf{Copie 4 (confusion Q1/Q2, Arendt) — Note : 1/5.}
\begin{itemize}
    \item Erreur 1 : confusion des exercices — l'élève donne son \emph{avis} (« Je suis d'accord ») là où Q1/Q2 exigent l'analyse objective du texte.
    \item Erreur 2 : exemple camerounais plaqué et inexact (le fonctionnaire « travaille » au sens d'Arendt ? confusion travail/œuvre).
    \item Erreur 3 : aucune analyse argumentative (types d'arguments, structure absents).
    \item Correction : « Arendt distingue le travail, lié au cycle biologique de la consommation, et l'œuvre, qui institue un monde durable ; cette distinction conceptuelle (et non un jugement) est la thèse du texte. »
\end{itemize}
\end{corrige}

\begin{methode}[Rappel de la grille officielle]
0 = hors sujet ; 1 = insuffisant (paraphrase, avis personnel) ; 2 = passable (éléments justes mais non structurés) ; 3 = assez bien (méthode appliquée, imprécisions) ; 4 = bien (rigoureux, nuancé) ; 5 = très bien (exhaustif, personnel, irréprochable).
\end{methode}

\begin{corrige}[Exercice 9 — Sujet Type Bac : Justice climatique (Achim Steiner)]
\textbf{Question 1 — Explicitation (/6).} \emph{Thèse centrale :} la transition écologique n'est durable qu'à condition d'être \emph{juste} ; sans justice, elle n'est qu'un \cle{néo-extractivisme} vert. \emph{Concepts clés et définitions contextuelles :} la \cle{justice climatique} (répartition équitable des bénéfices et des charges de la transition entre Nord et Sud) ; le \cle{néo-extractivisme} (reproduction, sous couvert d'écologie, du pillage des ressources du Sud par le Nord) ; la \cle{dette écologique} (obligation de réparation du Nord envers le Sud). \emph{Enchaînement des idées :} (1) les technologies vertes promettent le découplage prospérité/destruction ; (2) mais leur déploiement reproduit les inégalités (extraction, délocalisation des coûts) ; (3) donc une transition injuste n'est pas une transition ; (4) il faut gouvernance mondiale, transfert de technologie et reconnaissance de la dette écologique. \emph{Barème :} thèse 1,5 pt ; concepts définis 2 pts ; enchaînement 1,5 pt ; rédaction 1 pt.

\textbf{Question 2 — Analyse (/7).} \emph{Structure :} raisonnement progressif en trois temps — constat optimiste, réfutation par les faits, prescription normative. \emph{Types d'arguments :} arguments de fait (cobalt du Katanga, D3EE du Golfe de Guinée) ; argument d'\cle{injustice} (asymétrie profits/coûts) ; \cle{concession} implicite (« les technologies offrent les moyens\dots Pourtant\dots ») qui renforce la thèse ; argument de \emph{définition} (« une transition qui ignore\dots n'est pas une transition » : recours au concept pour disqualifier). \emph{Portée de la conclusion :} normative et universelle — la justice est posée comme \emph{condition} de durabilité (« Seule une transition juste peut être durable »), ce qui donne à la thèse une force prescriptive. \emph{Barème :} structure 2 pts ; typologie des arguments 3 pts ; concession et fonction 1 pt ; portée de la conclusion 1 pt.

\textbf{Question 3 — Discussion (/7).} \emph{Problématique attendue :} la justice climatique est-elle possible dans le cadre du modèle productiviste, ou exige-t-elle sa transformation ? \emph{Axes possibles :} (I) Oui, la technique verte peut être réformée de l'intérieur (optimisme technologique, régulation) — exemple CM : le barrage de Nachtigal financé avec clauses environnementales ; (II) Non, le modèle extractiviste reproduit structurellement l'injustice (Mbembe, critique du néo-extractivisme ; Jonas, responsabilité) — exemple CM : exportation des D3EE vers le port de Douala, accaparement foncier pour l'agro-industrie ; (III) Dépassement : une transition juste suppose de nouvelles normes de gouvernance et la responsabilité du Nord (dette écologique, transfert technologique), ce qui transforme le modèle sans abolir le développement. \emph{Barème :} problématique et introduction 2 pts ; dialectique argumentée 2 pts ; auteurs du programme 1,5 pt ; exemples camerounais 1 pt ; conclusion et langue 0,5 pt.
\end{corrige}

\begin{attention}[Points de vigilance]
Citer la \emph{condition d'application} de chaque concept (qu'est-ce qui rend une transition « juste » ?) ; ne pas transformer la Q3 en exposé sur l'environnement sans tension philosophique ; chaque exemple camerounais doit être \emph{rattaché} à un argument, non décoratif.
\end{attention}

\begin{corrige}[Exercice 10 — Sujet Bac Type : Hans Jonas, Le Principe responsabilité]
\textbf{Question 1 — Explicitation (/6).} \emph{Thèse :} la puissance inédite de la technique moderne rend caduque l'éthique traditionnelle et fonde une éthique nouvelle, celle de la \cle{responsabilité} envers la nature vulnérable et les générations futures. \emph{Concepts :} la \cle{technique moderne} (puissance planétaire, cumulative, irréversible, différée) ; la nature \cle{vulnérable} (qui ne « pardonne » plus) ; le \cle{nouvel impératif catégorique} (« Agis de façon que les effets de ton action soient compatibles avec la permanence d'une vie authentiquement humaine sur terre »). \emph{Enchaînement :} (1) l'éthique classique était une éthique du prochain, à portée limitée ; (2) la technique a changé la nature de l'action ; (3) les générations futures deviennent nos « prochains » ; (4) donc la responsabilité, et non la liberté abstraite, fonde l'éthique. \emph{Barème :} thèse 1,5 pt ; rupture caractérisée 2 pts ; impératif cité et expliqué 1,5 pt ; rédaction 1 pt.

\textbf{Question 2 — Analyse (/7).} \emph{Passage de l'ontologie à la norme :} Jonas procède par \cle{déduction} normative — prémisses descriptives (P1 : la puissance technique est irréversible ; P2 : la nature est vulnérable ; P3 : les générations futures sont affectées) d'où découle la prescription (C : il faut un impératif de responsabilité). \emph{Analogie avec Kant :} reprise explicite de la forme de l'impératif catégorique, mais déplacement du contenu : l'universalité \emph{logique} kantienne devient une universalité \emph{temporelle et écologique} (la survie de l'humanité). \emph{Type de finalisme :} finalisme éthique hérité d'Aristote — la conservation de la vie humaine authentique est posée comme fin en soi commandant les moyens techniques. \emph{Barème :} prémisses identifiées 2 pts ; analogie avec Kant analysée 2 pts ; finalisme 1,5 pt ; rigueur de la reconstitution 1,5 pt.

\textbf{Question 3 — Discussion (/7).} \emph{Problématique :} la responsabilité individuelle du technicien peut-elle suffire face à la puissance systémique de la technoscience ? \emph{Axes :} (I) Oui : la conscience morale de l'ingénieur est la première barrière (Jonas ; déontologie professionnelle) — exemple CM : le refus d'un technicien de valider des données biométriques non sécurisées ; (II) Non : Anders montre le « décalage prométhéen » — nous produisons plus que nous ne pouvons imaginer, d'où la \cle{honte prométhéenne} devant nos propres machines ; la responsabilité individuelle est dépassée par les systèmes (boîtes noires de l'IA, chaînes agro-industrielles d'OGM) ; (III) Dépassement : la responsabilité doit devenir \emph{collective et institutionnelle} (normes, comités d'éthique, loi-cadre sur l'environnement au Cameroun) sans dispenser la conscience individuelle. \emph{Barème :} identique à l'exercice 9 (2/2/1,5/1/0,5).
\end{corrige}

\begin{attention}[Points de vigilance]
Citer exactement l'impératif jonassien ; distinguer \emph{responsabilité morale} (individuelle) et \emph{responsabilité institutionnelle} ; ne pas confondre Anders (honte prométhéenne) et Jonas (principe responsabilité) : les articuler, non les opposer artificiellement.
\end{attention}

\begin{corrige}[Exercice 11 — Grand Oral Transversal : Dilemme éthique professionnel]
\textbf{Corrigé type (exemple : filière Informatique — IA et données personnelles).}
\begin{enumerate}
    \item \emph{Accroche (30 s) :} « En stage dans un hôpital de Yaoundé, j'ai vu un logiciel d'aide au diagnostic proposer un traitement sans que le médecin puisse expliquer pourquoi : qui est responsable si le patient en meurt ? »
    \item \emph{Texte d'ancrage (1 min) :} Jonas, \emph{Le Principe responsabilité} : « Agis de façon que les effets de ton action soient compatibles avec la permanence d'une vie authentiquement humaine sur terre. » Concept clé : la \cle{responsabilité} élargie aux effets différés et irréversibles de la technique.
    \item \emph{Problématique (30 s) :} l'efficacité diagnostique de l'IA peut-elle primer sur la transparence et le consentement du patient ?
    \item \emph{Analyse (2 min) :} Thèse : l'IA sauve des vies, refuser son usage serait irresponsable (conséquentialisme, Bacon). Antithèse : la « boîte noire » viole l'autonomie du patient et déresponsabilise le médecin (Kant : l'homme comme fin ; Jonas). Synthèse : usage encadré — l'IA propose, l'homme dispose et explique (responsabilité partagée, traçabilité des algorithmes).
    \item \emph{Positionnement professionnel (1 min) :} en tant que futur technicien, exiger l'auditabilité des systèmes, le consentement éclairé, la conformité à la loi camerounaise sur la protection des données personnelles et aux normes ISO de sécurité informatique.
    \item \emph{Ouverture (30 s) :} comment former durablement les techniciens camerounais à l'éthique du numérique ?
\end{enumerate}
\end{corrige}

\begin{methode}[Grille d'évaluation Grand Oral]
Prise de parole (clarté, gestion du temps 5 min) : /4 ; connaissances philosophiques exactes (auteur, concept, citation) : /4 ; qualité de l'argumentation dialectique : /4 ; ancrage professionnel et citoyen (dilemme réel, référentiel déontologique) : /4 ; interaction et ouverture : /4. Total /20.
\end{methode}

\begin{attention}[Points de vigilance]
Respecter impérativement les durées de chaque séquence ; citer le texte philosophique \emph{de mémoire et brièvement} (3-4 lignes) ; le dilemme doit être une vraie tension (deux exigences légitimes opposées), non un problème technique à solution unique ; le positionnement final doit engager une \cle{conduite responsable} concrète et référencée.
\end{attention}

\newpage

\coursec{Fiche bilan}

\begin{popfigure}
\begin{tikzpicture}[
  mindmap,
  concept color=popblue,
  level 1 concept/.append style={level distance=4.5cm, sibling angle=60},
  level 2 concept/.append style={level distance=3cm, sibling angle=45},
  every node/.style={font=\small, align=center},
  branch1/.style={concept color=poporange},
  branch2/.style={concept color=popgreen},
  branch3/.style={concept color=poppurple},
  branch4/.style={concept color=poppink},
  branch5/.style={concept color=popgold},
  branch6/.style={concept color=popdark}
]
\node[concept, text=white, minimum size=2.2cm, align=center]{M\'ethodologie\\ Texte Philo\\ Terminale Tech}
  child[branch1] { node[concept, text=white, minimum size=1.8cm, align=center]{1. Architecture\\ \& Gestion temps\\ (3h)} }
  child[branch2] { node[concept, text=white, minimum size=1.8cm, align=center]{2. Question 1\\ Explicitation\\ du sens (Savoir)} }
  child[branch3] { node[concept, text=white, minimum size=1.8cm, align=center]{3. Question 2\\ Analyse\\ argumentation (S-F)} }
  child[branch4] { node[concept, text=white, minimum size=1.8cm, align=center]{4. Question 3\\ Discussion\\ critique (Comp\'etence)} }
  child[branch5] { node[concept, text=white, minimum size=1.8cm, align=center]{5. Atelier\\ Simulation Bac\\ \& Auto-\'eval.} }
  child[branch6] { node[concept, text=white, minimum size=1.8cm, align=center]{6. Transfert\\ Citoyennet\'e\\ Esprit critique} };
\end{tikzpicture}
\legende{Carte mentale de la le\c{c}on 03 : Consolidation m\'ethodologie texte philosophique}
\end{popfigure}

\coursec{Architecture de l'\'epreuve et gestion strat\'egique du temps (3h)}
\begin{itemize}
\item \cle{Dur\'ee} : 3 heures (180 minutes). \cle{Coefficient} : Variable selon s\'eries (souvent coeff 2 ou 3).
\item \cle{Structure type} : Un texte d'auteur (philosophie g\'en\'erale/technique) + 3 questions obligatoires.
\item \cle{R\'epartition temporelle id\'eale (total 180 min)} :
\begin{center}
\begin{tabularx}{\linewidth}{|l|X|c|}\hline
\rowcolor{popblueL} \textbf{Phase} & \textbf{Activités} & \textbf{Durée} \\ \hline
Lecture active & 3 lectures (globale, analytique, structurelle) + repérage thèse/arguments & 30 min \\ \hline
Question 1 (Savoir) & Reformulation thèse, définition concepts, explication passage, contexte & 45 min \\ \hline
Question 2 (Savoir-faire) & Analyse structurelle : plan, enchaînement, procédés, types de raisonnement & 45 min \\ \hline
Question 3 (Compétence) & Discussion critique : thèse, antithèse, synthèse, exemples personnels, ouverture & 45 min \\ \hline
Relecture & Correction syntaxe, orthographe, cohérence, respect consignes & 15 min \\ \hline
\end{tabularx}
\end{center}
\item \textbf{Règle d'or} : Ne jamais commencer à écrire avant d'avoir un \emph{plan détaillé} pour Q2 et Q3 sur brouillon.
\end{itemize}

\coursec{Maîtrise de la Question 1 : L'explicitation rigoureuse du sens (Savoir)}
\begin{methode}[Q1 : 3 sous-tâches types]
\begin{enumerate}
\item \textbf{Thèse + Problème} : Formuler la thèse en 1 phrase (sujet + verbe + complément). Identifier le problème philosophique (tension conceptuelle).
\item \textbf{Concepts clés} : Définir \textbf{précisément} 2 à 3 concepts centraux (sens philosophique, pas dictionnaire courant). Citer l'auteur si définition explicite dans texte.
\item \textbf{Explication de passage} : « Expliquez la phrase : "..." ». Méthode : Contextualiser $\rightarrow$ Analyser termes $\rightarrow$ Montrer lien avec thèse $\rightarrow$ Illustrer par exemple concret (Cameroun si possible).
\end{enumerate}
\end{methode}
\begin{attention}[Pièges Q1]
\begin{enumerate}
\item Paraphraser au lieu d'expliquer (répéter le texte avec d'autres mots).
\item Définitions hors-sujet ou sens commun (ex: « La technique = les machines »).
\item Oublier de relier l'explication du passage à la thèse générale du texte.
\item Confondre \emph{problème} (tension) et \emph{sujet} (thème général).
\end{enumerate}
\end{attention}

\coursec{Maîtrise de la Question 2 : L'analyse de l'argumentation (Savoir-faire)}
\begin{definition}[Analyse argumentative]
Démontrer \emph{comment} l'auteur prouve sa thèse. Étudier la \textbf{structure} (plan), l'\textbf{enchaînement} (connecteurs, transitions), les \textbf{procédés} (exemples, analogies, réfutations, citations, expérience de pensée) et les \textbf{types de raisonnement} (déductif, inductif, abductif, dialectique).
\end{definition}
\begin{methode}[Rédaction Q2 standard]
\begin{enumerate}
\item \textbf{Introduction} : Thèse + Problème + Annonce du plan détaillé (2 ou 3 parties principales).
\item \textbf{Développement structuré} : Une partie = Une étape majeure du raisonnement.
  \begin{itemize}
  \item Titre de partie (ex: « I. La critique de l'illusion technique »).
  \item Résumé de l'argument principal de la partie.
  \item \textbf{Analyse fine} : Citer court (guillemets) $\rightarrow$ Nommer le procédé/raisonnement $\rightarrow$ Expliquer son efficacité.
  \item Transition vers partie suivante.
  \end{itemize}
\item \textbf{Conclusion} : Résumé de la dynamique argumentative (cheminement vers la thèse).
\end{enumerate}
\end{methode}
\begin{propriete}[Procédés fréquents au Bac Tech]
\begin{center}
\begin{tabularx}{\linewidth}{|l|X|l|}\hline
\rowcolor{popblueL} \textbf{Procédé} & \textbf{Description} & \textbf{Effet visé} \\ \hline
Analogie / Comparaison & Rapprochement deux domaines (ex: organisme / machine) & Rendre intelligible, transférer propriétés \\ \hline
Exemple concret / Contrefactuel & Cas réel (Cameroun) ou imaginaire & Ancrer dans réel, tester validité \\ \hline
Réfutation (objection) & Présenter thèse adverse $\rightarrow$ La détruire & Renforcer sa propre thèse par élimination \\ \hline
Citation d'autorité & Appel à un grand auteur (Aristote, Kant, Heidegger...) & Légitimer, inscrire dans tradition \\ \hline
Expérience de pensée & Scénario fictif « Que se passerait-il si... » & Isoler variable conceptuelle \\ \hline
Définition par genre/différence & Fixer le sens d'un concept & Éliminer ambiguïtés, poser bases \\ \hline
\end{tabularx}
\end{center}
\end{propriete}

\coursec{Maîtrise de la Question 3 : La discussion critique structurée (Compétence supérieure)}
\begin{definition}[Discussion critique]
Ne pas donner son avis (« Je pense que... »), mais \textbf{évaluer la solidité} de la thèse et des arguments. Confronter le texte à d'autres connaissances (auteurs, concepts, réalités techniques/sociales).
\end{definition}
\begin{methode}[Structure dialectique rigoureuse (Plan type)]
\begin{enumerate}
\item \textbf{Introduction} : Accroche (citation/actualité) $\rightarrow$ Rappel thèse texte + Problème $\rightarrow$ Annonce plan : « Nous examinerons d'abord la portée de la thèse (I), avant d'en tester les limites (II), pour enfin dégager les conditions de son dépassement (III). »
\item \textbf{I. Portée / Validité (Thèse)} : Pourquoi la thèse est-elle forte ? Quels arguments tiennent ? Quels concepts sont opérants ? Appui sur auteurs alliés (ex: Heidegger pour critique technique, Jonas pour responsabilité).
\item \textbf{II. Limites / Objections (Antithèse)} : Quels angles morts ? Présupposés non prouvés ? Contre-exemples (réalité camerounaise : informel, low-tech, savoirs endogènes) ? Auteurs opposés (ex: Bergson/Simondon contre Heidegger ; Prométhéens).
\item \textbf{III. Dépassement / Nuance (Synthèse)} : Conditions de validité (domaine, échelle). Distinctions conceptuelles (technique/technologie, moyen/fin, aliéner/libérer). Perspective ouverte (éthique, écologie, éducation technique).
\item \textbf{Conclusion} : Bilan synthétique $\rightarrow$ Réponse nuancée au problème $\rightarrow$ Ouverture (autre domaine, question nouvelle).
\end{enumerate}
\end{methode}
\begin{attention}[Erreurs fatales Q3]
\begin{enumerate}
\item « Catalogue d'auteurs » sans lien avec le texte (name-dropping).
\item Exemples personnels anecdotiques non conceptualisés.
\item Plan « Pour / Contre » binaire sans synthèse (reste au stade dialogue de sourds).
\item Oublier le texte de départ : la discussion doit \emph{porter sur le texte}, pas sur le thème en général.
\item Langage familier, manque de connecteurs logiques (\emph{donc, toutefois, par conséquent, en revanche, c'est pourquoi}).
\end{enumerate}
\end{attention}

\coursec{Atelier de consolidation : Simulation Bac \& Auto-évaluation}
\begin{exemplebox}[Grille d'auto-évaluation (à utiliser sur brouillon après rédaction)]
\begin{center}
\begin{tabularx}{\linewidth}{|l|X|c|}\hline
\rowcolor{popblueL} \textbf{Critère} & \textbf{Indicateurs de réussite} & \textbf{OK ?} \\ \hline
Q1 - Thèse & Formulée en 1 phrase complète, sujet/verbe/complément, vocabulaire précis & $\square$ \\ \hline
Q1 - Concepts & 2-3 définis philosophiquement, distinction sens commun/sens savant & $\square$ \\ \hline
Q1 - Passage & Explication structurée (contexte, termes, lien thèse, exemple) & $\square$ \\ \hline
Q2 - Plan & 2-3 parties logiques, titres explicites, transitions fluides & $\square$ \\ \hline
Q2 - Analyse & Procédés nommés + citations courtes + explication efficacité & $\square$ \\ \hline
Q3 - Structure & Plan dialectique I/II/III visible, thèse/antithèse/synthèse & $\square$ \\ \hline
Q3 - Argumentation & Objections fondées, auteurs/concepts mobilisés, exemples pertinents & $\square$ \\ \hline
Q3 - Synthèse & Vrai dépassement (pas simple résumé), distinctions conceptuelles & $\square$ \\ \hline
Forme & Orthographe, syntaxe, connecteurs, paragraphes, lisibilité & $\square$ \\ \hline
Temps & Respect timing (Q1$\approx$45', Q2$\approx$45', Q3$\approx$45', Relecture 15') & $\square$ \\ \hline
\end{tabularx}
\end{center}
\end{exemplebox}

\coursec{Transfert \& Citoyenneté : L'esprit critique hors les murs}
\begin{savaistu}[Pourquoi cette méthode change la vie ?]
L'exercice sur texte n'est pas un exercice scolaire artificiel. C'est l'entraînement à la \cle{pensée critique} :
\begin{itemize}
\item \textbf{Comprendre avant de juger} (Q1) : Lutter contre les fake news, les rumeurs, les raccourcis médiatiques.
\item \textbf{Analyser la rhétorique} (Q2) : Détecter les manipulations (pub, politique, réseaux sociaux), identifier les arguments d'autorité fallacieux, les fausses analogies.
\item \textbf{Argumenter avec rigueur} (Q3) : Construire un avis éclairé, négocier, écrire un rapport technique, défendre un projet, participer au débat démocratique (ex: gestion ressources, éthique IA, aménagement territoire au Cameroun).
\end{itemize}
L'ingénieur/technicien philosophe ne subit pas la technique, il l'\emph{interroge}.
\end{savaistu}

\begin{aretenir}[Checklist avant le Bac : $\square$ = À vérifier la veille / le matin]
\begin{enumerate}
\item $\square$ \textbf{Matériel} : Stylos bleus/noirs (2 min), correcteur (si autorisé), montre/horloge, convocation, pièce d'identité.
\item $\square$ \textbf{Stratégie temps} : 3h = 30' lecture + 45' Q1 + 45' Q2 + 45' Q3 + 15' relecture. \textbf{Ne pas déborder sur Q1/Q2}.
\item $\square$ \textbf{Brouillon obligatoire} : Plan Q2 (structure + procédés), Plan Q3 (I/II/III + arguments/auteurs/exemples), Mots-clés Q1.
\item $\square$ \textbf{Q1} : Thèse en 1 phrase. 3 concepts définis \emph{philosophiquement}. Passage : Contexte + Analyse + Lien thèse + Exemple concret.
\item $\square$ \textbf{Q2} : Intro (Thèse/Problème/Plan). Parties titrées. \textbf{Citer court} $\rightarrow$ Nommer procédé $\rightarrow$ Expliquer force. Transitions logiques.
\item $\square$ \textbf{Q3} : Intro (Amorce/Problème/Plan dialectique). I. Forces (alliés). II. Limites (adversaires/contre-exemples réel). III. Nuance/Dépassement (distinctions). Conclusion (Bilan/Ouverture).
\item $\square$ \textbf{Exemples camerounais} : Prêts pour Q1 (illustration) et Q3 (contre-exemple/ancrage) : informel, low-tech, savoirs locaux, enjeux énergie/numérique/formation.
\item $\square$ \textbf{Auteurs repères} : Maîtriser 3--4 couples thèse/antithèse techniques (Heidegger/Jonas vs Simondon/Bergson/Prométhéens ; Ellul vs Mumford ; Arendt vs Marx sur travail/technique).
\item $\square$ \textbf{Connecteurs logiques} : \emph{Premièrement, ensuite, cependant, par conséquent, en effet, notamment, en somme, dès lors, eu égard à}.
\item $\square$ \textbf{Relecture active} : Sujet/verbe, accords, paragraphes (saut de ligne), lisibilité, respect consignes (numéros questions), pas de HS.
\item $\square$ \textbf{Hygiène mentale} : Sommeil veille, repas léger, respiration ventrale si stress, confiance : \emph{« Je sais lire, analyser, argumenter. Je suis prêt. »}
\end{enumerate}
\end{aretenir}

\findecours

\end{document}
